% !TEX program = lualatex
% ---------------------------------------------------------------------------------------------
%  The illustrated blueprint of the architecture of the announced proof of RH  (riemann-rev)
%
%  Build:   make -C blueprint            (figures, generated parts, two LuaLaTeX passes)
%  Frame:   this file (hand-written); every number, statement, name and status shown comes from
%           generated/*.tex, which scripts/build_blueprint_tex.py writes from data/model.json,
%           data/announcement_dates.json, data/source/*, verification/logs/axioms.json, the SVG
%           figures and the Lean files (quoted verbatim).
% ---------------------------------------------------------------------------------------------
\documentclass[11pt,a4paper,oneside]{article}
% reproducible output: do not copy the info dictionaries (creation dates) of included PDFs, the
% file names or the banner into the result; the dates themselves come from SOURCE_DATE_EPOCH
\ifdefined\pdfvariable\pdfvariable suppressoptionalinfo \numexpr 1+2+4+8\relax\fi
\usepackage[a4paper,hmargin=24mm,top=24mm,bottom=27mm,headheight=14pt,headsep=7mm,footskip=11mm]{geometry}

%% ------------------------------------------------------------------ fonts (LuaLaTeX, TeX Live)
\usepackage{fontspec}
\directlua{
  luaotfload.add_fallback("bpfallback", {"STIXTwoMath-Regular.otf:mode=node;", "DejaVuSans.ttf:mode=node;"})
  luaotfload.add_fallback("bpmonofallback", {"DejaVuSansMono.ttf:mode=node;", "STIXTwoMath-Regular.otf:mode=node;"})
}
\setmainfont{LibertinusSerif}[Extension=.otf,UprightFont=*-Regular,ItalicFont=*-Italic,
  BoldFont=*-Semibold,BoldItalicFont=*-SemiboldItalic,RawFeature={fallback=bpfallback}]
\setsansfont{Inter}[Extension=.otf,UprightFont=*-Regular,ItalicFont=*-Italic,
  BoldFont=*-SemiBold,BoldItalicFont=*-SemiBoldItalic,RawFeature={fallback=bpfallback}]
\setmonofont{JetBrainsMono}[Extension=.otf,UprightFont=*-Regular,ItalicFont=*-Italic,
  BoldFont=*-SemiBold,BoldItalicFont=*-SemiBoldItalic,Scale=MatchLowercase,
  RawFeature={-calt;-liga;fallback=bpmonofallback}]
\usepackage{amsmath}
\usepackage[math-style=TeX]{unicode-math}
\setmathfont{LibertinusMath-Regular.otf}

%% ------------------------------------------------------------------ packages
\usepackage{xcolor}
\usepackage{graphicx}
\usepackage{tikz}
\usetikzlibrary{calc,arrows.meta,backgrounds,positioning}
\usepackage{tcolorbox}
\tcbuselibrary{skins,breakable}
\usepackage{array,booktabs,longtable,tabularx,colortbl}
\usepackage{enumitem}
\usepackage{titlesec}
\usepackage{fancyhdr}
\usepackage{needspace}
\usepackage{caption}
\usepackage{fvextra}
\usepackage{pdflscape}
\usepackage{etoolbox}
\usepackage[protrusion=true,expansion=true]{microtype}
\usepackage[hyphens]{xurl}
\usepackage[pdfusetitle,bookmarksnumbered,bookmarksopen,bookmarksopenlevel=1,pdfdisplaydoctitle]{hyperref}
\usepackage{bookmark}

%% ------------------------------------------------------------------ colours
\definecolor{ink}{HTML}{0F172A}
\definecolor{inksoft}{HTML}{475569}
\definecolor{inkfaint}{HTML}{94A3B8}
\definecolor{rulegrey}{HTML}{CBD5E1}
\definecolor{wash}{HTML}{F8FAFC}
\definecolor{linkblue}{HTML}{1E40AF}
\definecolor{numerics}{HTML}{D97706}
% Generated by scripts/build_blueprint_tex.py from data/model.json -- do not edit by hand.
\definecolor{stratum-arith}{HTML}{C2410C}
\expandafter\def\csname bp@strattitle@arith\endcsname{Arithmetic supply}
\definecolor{stratum-est}{HTML}{7C3AED}
\expandafter\def\csname bp@strattitle@est\endcsname{The estimate}
\definecolor{stratum-ana}{HTML}{0891B2}
\expandafter\def\csname bp@strattitle@ana\endcsname{The analytic engine}
\definecolor{stratum-dest}{HTML}{E11D48}
\expandafter\def\csname bp@strattitle@dest\endcsname{The destination}
\definecolor{tier-two-kernels}{HTML}{047857}
\expandafter\def\csname bp@tierlabel@two-kernels\endcsname{Checked by two kernels}
\expandafter\def\csname bp@tiericon@two-kernels\endcsname{✓✓}
\definecolor{tier-checked}{HTML}{059669}
\expandafter\def\csname bp@tierlabel@checked\endcsname{Kernel-checked at the archive pin}
\expandafter\def\csname bp@tiericon@checked\endcsname{✓}
\definecolor{tier-partial}{HTML}{10B981}
\expandafter\def\csname bp@tierlabel@partial\endcsname{Partly checked}
\expandafter\def\csname bp@tiericon@partial\endcsname{◐}
\definecolor{tier-covered}{HTML}{34D399}
\expandafter\def\csname bp@tierlabel@covered\endcsname{Inside a checked closure}
\expandafter\def\csname bp@tiericon@covered\endcsname{⊂}
\definecolor{tier-read}{HTML}{4F46E5}
\expandafter\def\csname bp@tierlabel@read\endcsname{Source-read}
\expandafter\def\csname bp@tiericon@read\endcsname{◇}
\definecolor{tier-reported}{HTML}{64748B}
\expandafter\def\csname bp@tierlabel@reported\endcsname{Full replay reported}
\expandafter\def\csname bp@tiericon@reported\endcsname{↗}
\definecolor{tier-motivation}{HTML}{94A3B8}
\expandafter\def\csname bp@tierlabel@motivation\endcsname{Motivation}
\expandafter\def\csname bp@tiericon@motivation\endcsname{○}
\expandafter\def\csname bp@nodetitle@main-rh\endcsname{Mathlib RH}
\expandafter\def\csname bp@nodetier@main-rh\endcsname{two-kernels}
\expandafter\def\csname bp@nodetitle@main-zeta\endcsname{Zeta in expanded form}
\expandafter\def\csname bp@nodetier@main-zeta\endcsname{two-kernels}
\expandafter\def\csname bp@nodetitle@main-right\endcsname{Nonvanishing right of one half}
\expandafter\def\csname bp@nodetier@main-right\endcsname{two-kernels}
\expandafter\def\csname bp@nodetitle@main-strip\endcsname{Dirichlet GRH in the open strip}
\expandafter\def\csname bp@nodetier@main-strip\endcsname{two-kernels}
\expandafter\def\csname bp@nodetitle@main-prim\endcsname{Primitive critical-strip zeros}
\expandafter\def\csname bp@nodetier@main-prim\endcsname{two-kernels}
\expandafter\def\csname bp@nodetitle@stmt-zeta\endcsname{Zeta target type}
\expandafter\def\csname bp@nodetier@stmt-zeta\endcsname{read}
\expandafter\def\csname bp@nodetitle@stmt-right\endcsname{Half-plane target type}
\expandafter\def\csname bp@nodetier@stmt-right\endcsname{read}
\expandafter\def\csname bp@nodetitle@stmt-strip\endcsname{Dirichlet strip target type}
\expandafter\def\csname bp@nodetier@stmt-strip\endcsname{read}
\expandafter\def\csname bp@nodetitle@stmt-prim\endcsname{Primitive strip target type}
\expandafter\def\csname bp@nodetier@stmt-prim\endcsname{read}
\expandafter\def\csname bp@nodetitle@g-m\endcsname{The weighted prime sum}
\expandafter\def\csname bp@nodetier@g-m\endcsname{covered}
\expandafter\def\csname bp@nodetitle@window-def\endcsname{Uniform families on power windows}
\expandafter\def\csname bp@nodetier@window-def\endcsname{covered}
\expandafter\def\csname bp@nodetitle@vv\endcsname{Weighted lengths of local-global kernels}
\expandafter\def\csname bp@nodetier@vv\endcsname{read}
\expandafter\def\csname bp@nodetitle@vwin-def\endcsname{The cohomological window}
\expandafter\def\csname bp@nodetier@vwin-def\endcsname{covered}
\expandafter\def\csname bp@nodetitle@row-def\endcsname{The local–global Galois row}
\expandafter\def\csname bp@nodetier@row-def\endcsname{read}
\expandafter\def\csname bp@nodetitle@rh-all\endcsname{The two zeta exports share a proof}
\expandafter\def\csname bp@nodetier@rh-all\endcsname{read}
\expandafter\def\csname bp@nodetitle@all-chi\endcsname{All-character half-plane theorem}
\expandafter\def\csname bp@nodetier@all-chi\endcsname{read}
\expandafter\def\csname bp@nodetitle@crit-all\endcsname{Primitive critical-line wrapper}
\expandafter\def\csname bp@nodetier@crit-all\endcsname{read}
\expandafter\def\csname bp@nodetitle@claim-all\endcsname{Three endpoint propositions}
\expandafter\def\csname bp@nodetier@claim-all\endcsname{read}
\expandafter\def\csname bp@nodetitle@rh-every\endcsname{Specialize at modulus one}
\expandafter\def\csname bp@nodetier@rh-every\endcsname{read}
\expandafter\def\csname bp@nodetitle@rh-right\endcsname{Reflect a zeta zero}
\expandafter\def\csname bp@nodetier@rh-right\endcsname{checked}
\expandafter\def\csname bp@nodetitle@vw-crit\endcsname{Reflect a primitive character}
\expandafter\def\csname bp@nodetier@vw-crit\endcsname{read}
\expandafter\def\csname bp@nodetitle@gamma\endcsname{Gamma factors in the strip}
\expandafter\def\csname bp@nodetier@gamma\endcsname{read}
\expandafter\def\csname bp@nodetitle@strip\endcsname{Reflect an arbitrary character}
\expandafter\def\csname bp@nodetier@strip\endcsname{checked}
\expandafter\def\csname bp@nodetitle@prim-zero\endcsname{Remove imprimitive Euler factors}
\expandafter\def\csname bp@nodetier@prim-zero\endcsname{covered}
\expandafter\def\csname bp@nodetitle@factor\endcsname{Nonzero removed Euler factor}
\expandafter\def\csname bp@nodetier@factor\endcsname{covered}
\expandafter\def\csname bp@nodetitle@prim-reflect\endcsname{Transfer a primitive zero to 1−s}
\expandafter\def\csname bp@nodetier@prim-reflect\endcsname{covered}
\expandafter\def\csname bp@nodetitle@bundle-zeta\endcsname{Zeta component of the conjunction}
\expandafter\def\csname bp@nodetier@bundle-zeta\endcsname{checked}
\expandafter\def\csname bp@nodetitle@every\endcsname{Closed half-plane assembly}
\expandafter\def\csname bp@nodetier@every\endcsname{read}
\expandafter\def\csname bp@nodetitle@from-row\endcsname{The half-plane theorem from one row}
\expandafter\def\csname bp@nodetier@from-row\endcsname{read}
\expandafter\def\csname bp@nodetitle@row-prod\endcsname{Producer of the Galois row}
\expandafter\def\csname bp@nodetier@row-prod\endcsname{read}
\expandafter\def\csname bp@nodetitle@row-geo\endcsname{Galois representations from theta lifting}
\expandafter\def\csname bp@nodetier@row-geo\endcsname{read}
\expandafter\def\csname bp@nodetitle@tlift\endcsname{Theta-lift existence at high rank}
\expandafter\def\csname bp@nodetier@tlift\endcsname{read}
\expandafter\def\csname bp@nodetitle@geo-one\endcsname{Weight-one geometric realization}
\expandafter\def\csname bp@nodetier@geo-one\endcsname{read}
\expandafter\def\csname bp@nodetitle@rec\endcsname{Existence of local Langlands parameters}
\expandafter\def\csname bp@nodetier@rec\endcsname{read}
\expandafter\def\csname bp@nodetitle@lit\endcsname{Assemble the automorphic inputs}
\expandafter\def\csname bp@nodetier@lit\endcsname{read}
\expandafter\def\csname bp@nodetitle@open-record\endcsname{Three fields in the literature input}
\expandafter\def\csname bp@nodetier@open-record\endcsname{read}
\expandafter\def\csname bp@nodetitle@eigen-package\endcsname{The unitary and eigenvalue package}
\expandafter\def\csname bp@nodetier@eigen-package\endcsname{read}
\expandafter\def\csname bp@nodetitle@theta-package\endcsname{The local theta package}
\expandafter\def\csname bp@nodetier@theta-package\endcsname{read}
\expandafter\def\csname bp@nodetitle@hodge-package\endcsname{The Hodge-bundle package}
\expandafter\def\csname bp@nodetier@hodge-package\endcsname{read}
\expandafter\def\csname bp@nodetitle@local-rec\endcsname{The local parameter input at the call site}
\expandafter\def\csname bp@nodetier@local-rec\endcsname{read}
\expandafter\def\csname bp@nodetitle@local-rec-proj\endcsname{Read the parameter field}
\expandafter\def\csname bp@nodetier@local-rec-proj\endcsname{read}
\expandafter\def\csname bp@nodetitle@green\endcsname{Regular characters yield the Green input}
\expandafter\def\csname bp@nodetier@green\endcsname{read}
\expandafter\def\csname bp@nodetitle@linked-two-three\endcsname{Bind the estimates at 2 and 3}
\expandafter\def\csname bp@nodetier@linked-two-three\endcsname{read}
\expandafter\def\csname bp@nodetitle@goal-shape\endcsname{Apply the shape to the same chosen data}
\expandafter\def\csname bp@nodetier@goal-shape\endcsname{read}
\expandafter\def\csname bp@nodetitle@shape-fields\endcsname{Supply the actual fields of the shape}
\expandafter\def\csname bp@nodetier@shape-fields\endcsname{read}
\expandafter\def\csname bp@nodetitle@tlift-def\endcsname{Quantifiers of the theta lift}
\expandafter\def\csname bp@nodetier@tlift-def\endcsname{read}
\expandafter\def\csname bp@nodetitle@theta-match\endcsname{Match global integrals and local data}
\expandafter\def\csname bp@nodetier@theta-match\endcsname{read}
\expandafter\def\csname bp@nodetitle@covol\endcsname{Produce the covolume bound}
\expandafter\def\csname bp@nodetier@covol\endcsname{read}
\expandafter\def\csname bp@nodetitle@seven\endcsname{Seven inputs for every character}
\expandafter\def\csname bp@nodetier@seven\endcsname{read}
\expandafter\def\csname bp@nodetitle@definite\endcsname{Build a nonempty definite-side package}
\expandafter\def\csname bp@nodetier@definite\endcsname{read}
\expandafter\def\csname bp@nodetitle@every-green\endcsname{Pass from primitive to all characters}
\expandafter\def\csname bp@nodetier@every-green\endcsname{read}
\expandafter\def\csname bp@nodetitle@imprimitive\endcsname{Restore all characters}
\expandafter\def\csname bp@nodetier@imprimitive\endcsname{checked}
\expandafter\def\csname bp@nodetitle@zero-green\endcsname{Select the window width 3/2}
\expandafter\def\csname bp@nodetier@zero-green\endcsname{read}
\expandafter\def\csname bp@nodetitle@shape\endcsname{Choose a common geometric family}
\expandafter\def\csname bp@nodetier@shape\endcsname{read}
\expandafter\def\csname bp@nodetitle@family-design\endcsname{Combine divisibility, slopes and windows}
\expandafter\def\csname bp@nodetier@family-design\endcsname{read}
\expandafter\def\csname bp@nodetitle@divin\endcsname{Control the local divisibility loss}
\expandafter\def\csname bp@nodetier@divin\endcsname{read}
\expandafter\def\csname bp@nodetitle@vv-sum\endcsname{Sum the local losses}
\expandafter\def\csname bp@nodetier@vv-sum\endcsname{read}
\expandafter\def\csname bp@nodetitle@supply\endcsname{Produce cyclic character hosts}
\expandafter\def\csname bp@nodetier@supply\endcsname{checked}
\expandafter\def\csname bp@nodetitle@vwin-row\endcsname{Cohomological window from the row}
\expandafter\def\csname bp@nodetier@vwin-row\endcsname{read}
\expandafter\def\csname bp@nodetitle@window-inputs\endcsname{Supply the automorphic window inputs}
\expandafter\def\csname bp@nodetier@window-inputs\endcsname{read}
\expandafter\def\csname bp@nodetitle@depth-window\endcsname{Build the rank-N package}
\expandafter\def\csname bp@nodetier@depth-window\endcsname{read}
\expandafter\def\csname bp@nodetitle@lit-window\endcsname{Window from rank-N automorphy}
\expandafter\def\csname bp@nodetier@lit-window\endcsname{read}
\expandafter\def\csname bp@nodetitle@patch-window\endcsname{Pass through the descent hypotheses}
\expandafter\def\csname bp@nodetier@patch-window\endcsname{read}
\expandafter\def\csname bp@nodetitle@four-produce\endcsname{Produce the four counting bounds}
\expandafter\def\csname bp@nodetier@four-produce\endcsname{read}
\expandafter\def\csname bp@nodetitle@four-use\endcsname{Combine the cohomology counts}
\expandafter\def\csname bp@nodetier@four-use\endcsname{read}
\expandafter\def\csname bp@nodetitle@zero-top\endcsname{Produce primitive zero-free levels}
\expandafter\def\csname bp@nodetier@zero-top\endcsname{checked}
\expandafter\def\csname bp@nodetitle@tied-cyc\endcsname{Assemble the cyclic levels}
\expandafter\def\csname bp@nodetier@tied-cyc\endcsname{covered}
\expandafter\def\csname bp@nodetitle@tied\endcsname{Pass from levels to characters}
\expandafter\def\csname bp@nodetier@tied\endcsname{covered}
\expandafter\def\csname bp@nodetitle@window-consumer\endcsname{Choose epsilon after a candidate zero}
\expandafter\def\csname bp@nodetier@window-consumer\endcsname{covered}
\expandafter\def\csname bp@nodetitle@prime-bound\endcsname{The prime-sum inequality on tied families}
\expandafter\def\csname bp@nodetier@prime-bound\endcsname{covered}
\expandafter\def\csname bp@nodetitle@artin-window\endcsname{Estimate the weighted level on a window}
\expandafter\def\csname bp@nodetier@artin-window\endcsname{checked}
\expandafter\def\csname bp@nodetitle@core-gm\endcsname{From the height discrepancy to the prime sum}
\expandafter\def\csname bp@nodetier@core-gm\endcsname{checked}
\expandafter\def\csname bp@nodetitle@ch-top\endcsname{Bound the discrepancy per index}
\expandafter\def\csname bp@nodetier@ch-top\endcsname{checked}
\expandafter\def\csname bp@nodetitle@ch-core\endcsname{Average over characters}
\expandafter\def\csname bp@nodetier@ch-core\endcsname{checked}
\expandafter\def\csname bp@nodetitle@form-b\endcsname{The analytic criterion for a level}
\expandafter\def\csname bp@nodetier@form-b\endcsname{covered}
\expandafter\def\csname bp@nodetitle@landau\endcsname{The character-family power-window theorem}
\expandafter\def\csname bp@nodetier@landau\endcsname{checked}
\expandafter\def\csname bp@nodetitle@landau-lemma\endcsname{Apply the pole obstruction}
\expandafter\def\csname bp@nodetier@landau-lemma\endcsname{covered}
\expandafter\def\csname bp@nodetitle@pole\endcsname{Poles of the finite-family transform}
\expandafter\def\csname bp@nodetier@pole\endcsname{checked}
\expandafter\def\csname bp@nodetitle@oscillation\endcsname{Rule out a pole right of \(θ\)}
\expandafter\def\csname bp@nodetier@oscillation\endcsname{covered}
\expandafter\def\csname bp@nodetitle@alternate-rh\endcsname{A second zeta deduction}
\expandafter\def\csname bp@nodetier@alternate-rh\endcsname{read}
\expandafter\def\csname bp@nodetitle@rank\endcsname{The geometric ranks}
\expandafter\def\csname bp@nodetier@rank\endcsname{read}
\expandafter\def\csname bp@nodetitle@minor\endcsname{A local source of divisibility}
\expandafter\def\csname bp@nodetier@minor\endcsname{checked}
\expandafter\def\csname bp@nodetitle@behrend\endcsname{Behrend’s collision-free indices}
\expandafter\def\csname bp@nodetier@behrend\endcsname{checked}
\expandafter\def\csname bp@nodetitle@empty-data\endcsname{An existence check on a period identity}
\expandafter\def\csname bp@nodetier@empty-data\endcsname{read}
\expandafter\def\csname bp@nodetitle@slope\endcsname{The lattice slope inequality}
\expandafter\def\csname bp@nodetier@slope\endcsname{read}
\expandafter\def\csname bp@nodetitle@paper-chain\endcsname{The square-root multiplicity analogy}
\expandafter\def\csname bp@nodetier@paper-chain\endcsname{motivation}
\expandafter\def\csname bp@nodetitle@paper-behrend\endcsname{The tower selection in the timeline}
\expandafter\def\csname bp@nodetier@paper-behrend\endcsname{motivation}
\expandafter\def\csname bp@nodetitle@tower-window-from-lemma-k\endcsname{Assemble the arithmetic tower window}
\expandafter\def\csname bp@nodetier@tower-window-from-lemma-k\endcsname{checked}
\expandafter\def\csname bp@nodetitle@tower-window-constants\endcsname{Choose one set of tower constants}
\expandafter\def\csname bp@nodetier@tower-window-constants\endcsname{checked}
\expandafter\def\csname bp@nodetitle@tower-bound-from-rank\endcsname{From a finitely generated tower to two length bounds}
\expandafter\def\csname bp@nodetier@tower-bound-from-rank\endcsname{checked}
\expandafter\def\csname bp@nodetitle@tower-operator-bound\endcsname{Control exceptional twists by the fixed layer}
\expandafter\def\csname bp@nodetier@tower-operator-bound\endcsname{covered}
\expandafter\def\csname bp@nodetitle@tower-class-rank\endcsname{Uniform ℓ-rank of class groups in the tower}
\expandafter\def\csname bp@nodetier@tower-class-rank\endcsname{read}
\expandafter\def\csname bp@nodetitle@tower-fixed-growth\endcsname{Supply the absolute bound for the reference layer}
\expandafter\def\csname bp@nodetier@tower-fixed-growth\endcsname{read}
\expandafter\def\csname bp@nodetitle@counts-i\endcsname{Case I gives one bound across coefficient primes}
\expandafter\def\csname bp@nodetier@counts-i\endcsname{read}
\expandafter\def\csname bp@nodetitle@counts-ii-iii\endcsname{Cases II and III from the rank-four symplectic argument}
\expandafter\def\csname bp@nodetier@counts-ii-iii\endcsname{read}
\expandafter\def\csname bp@nodetitle@counts-v\endcsname{Case V at twist one}
\expandafter\def\csname bp@nodetier@counts-v\endcsname{read}
\expandafter\def\csname bp@nodetitle@counts-common-integer\endcsname{A nonzero integer controls one character orbit}
\expandafter\def\csname bp@nodetier@counts-common-integer\endcsname{read}
\expandafter\def\csname bp@nodetitle@counts-window\endcsname{Convert the joint window into \bpcode{VWin}}
\expandafter\def\csname bp@nodetier@counts-window\endcsname{read}
\expandafter\def\csname bp@nodetitle@counts-joint\endcsname{Reuse the tower interval for the positive part}
\expandafter\def\csname bp@nodetier@counts-joint\endcsname{read}
\expandafter\def\csname bp@nodetitle@counts-height\endcsname{Evaluate the pointwise estimate on admissible windows}
\expandafter\def\csname bp@nodetier@counts-height\endcsname{read}
\expandafter\def\csname bp@nodetitle@counts-pointwise\endcsname{Combine the two prime sums at a point}
\expandafter\def\csname bp@nodetier@counts-pointwise\endcsname{read}
\expandafter\def\csname bp@nodetitle@counts-inside\endcsname{Add contributions inside the exceptional set}
\expandafter\def\csname bp@nodetier@counts-inside\endcsname{read}
\expandafter\def\csname bp@nodetitle@counts-outside\endcsname{Bound the sum over arbitrary finite sets of other primes}
\expandafter\def\csname bp@nodetier@counts-outside\endcsname{read}
\expandafter\def\csname bp@nodetitle@moment-target\endcsname{The first moment of the Chebyshev error}
\expandafter\def\csname bp@nodetier@moment-target\endcsname{covered}
\expandafter\def\csname bp@nodetitle@moment-theorem\endcsname{Evaluate the first moment}
\expandafter\def\csname bp@nodetier@moment-theorem\endcsname{checked}
\expandafter\def\csname bp@nodetitle@moment-calculation\endcsname{Match the two limits at one}
\expandafter\def\csname bp@nodetier@moment-calculation\endcsname{covered}
\expandafter\def\csname bp@nodetitle@moment-majorant\endcsname{Dominate the centered character sum}
\expandafter\def\csname bp@nodetier@moment-majorant\endcsname{covered}
\expandafter\def\csname bp@nodetitle@moment-dominated\endcsname{Pass the Mellin integral to \(σ\) = 1}
\expandafter\def\csname bp@nodetier@moment-dominated\endcsname{covered}
\expandafter\def\csname bp@nodetitle@moment-laurent\endcsname{Read Euler’s constant from the logarithmic derivative}
\expandafter\def\csname bp@nodetier@moment-laurent\endcsname{covered}
\expandafter\def\csname bp@comptitle@endpoints\endcsname{The five endpoints}
\expandafter\def\csname bp@compnum@endpoints\endcsname{1}
\expandafter\def\csname bp@compstrat@endpoints\endcsname{dest}
\expandafter\def\csname bp@comptitle@reflection\endcsname{Reflection}
\expandafter\def\csname bp@compnum@reflection\endcsname{2}
\expandafter\def\csname bp@compstrat@reflection\endcsname{dest}
\expandafter\def\csname bp@comptitle@descent\endcsname{Conductor descent}
\expandafter\def\csname bp@compnum@descent\endcsname{3}
\expandafter\def\csname bp@compstrat@descent\endcsname{dest}
\expandafter\def\csname bp@comptitle@statistic\endcsname{The centered prime sum}
\expandafter\def\csname bp@compnum@statistic\endcsname{4}
\expandafter\def\csname bp@compstrat@statistic\endcsname{ana}
\expandafter\def\csname bp@comptitle@residues\endcsname{Mellin transform and positive residues}
\expandafter\def\csname bp@compnum@residues\endcsname{5}
\expandafter\def\csname bp@compstrat@residues\endcsname{ana}
\expandafter\def\csname bp@comptitle@criterion\endcsname{The window criterion}
\expandafter\def\csname bp@compnum@criterion\endcsname{6}
\expandafter\def\csname bp@compstrat@criterion\endcsname{ana}
\expandafter\def\csname bp@comptitle@choreography\endcsname{The order of choices}
\expandafter\def\csname bp@compnum@choreography\endcsname{7}
\expandafter\def\csname bp@compstrat@choreography\endcsname{ana}
\expandafter\def\csname bp@comptitle@comparison\endcsname{The finite character comparison}
\expandafter\def\csname bp@compnum@comparison\endcsname{8}
\expandafter\def\csname bp@compstrat@comparison\endcsname{est}
\expandafter\def\csname bp@comptitle@family\endcsname{The geometric family and the theta chain}
\expandafter\def\csname bp@compnum@family\endcsname{9}
\expandafter\def\csname bp@compstrat@family\endcsname{est}
\expandafter\def\csname bp@comptitle@hosts\endcsname{Hosts for every character}
\expandafter\def\csname bp@compnum@hosts\endcsname{10}
\expandafter\def\csname bp@compstrat@hosts\endcsname{est}
\expandafter\def\csname bp@comptitle@cohomology\endcsname{The cohomological window}
\expandafter\def\csname bp@compnum@cohomology\endcsname{11}
\expandafter\def\csname bp@compstrat@cohomology\endcsname{arith}
\expandafter\def\csname bp@comptitle@primes\endcsname{Auxiliary primes and Behrend towers}
\expandafter\def\csname bp@compnum@primes\endcsname{12}
\expandafter\def\csname bp@compstrat@primes\endcsname{arith}
\expandafter\def\csname bp@comptitle@galois\endcsname{The Galois row}
\expandafter\def\csname bp@compnum@galois\endcsname{13}
\expandafter\def\csname bp@compstrat@galois\endcsname{arith}
\expandafter\def\csname bp@comptitle@inputs\endcsname{Seven inputs at the decisive call}
\expandafter\def\csname bp@compnum@inputs\endcsname{14}
\expandafter\def\csname bp@compstrat@inputs\endcsname{arith}
\newcommand\ModelTitle{The architecture of the announced proof of the Riemann hypothesis}
\newcommand\ModuleCount{175,187}
\newcommand\ArchiveParts{14}
\newcommand\NumComponents{14}
\newcommand\NumNodes{109}
\newcommand\NumDecls{107}
\newcommand\NumIdeas{2}
\newcommand\NumEdges{113}
\newcommand\NumInterfaces{14}
\newcommand\NumCuts{2}
\newcommand\NumTimeline{20}
\newcommand\NumStrataWord{four}
\newcommand\NumCutsWord{two}
\newcommand\NumItinerariesWord{five}
\newcommand\NumExportsWord{five}
\newcommand\TimelinePages{3 and 4}
\newcommand\CutITitle{One-sided prime windows imply RH}
\newcommand\CutIITitle{\bpcode{TopG} and the cohomological windows imply RH and Dirichlet GRH}
\newcommand\CutITier{two-kernels}
\newcommand\CutIITier{checked}
\newcommand\StatusHeadline{Two independent kernels accepted the complete proof, and all five theorems matched the canonical Mathlib challenge.}
\newcommand\StatusDetail{The first theorem has Mathlib’s type \bpcode{RiemannHypothesis}. The other endpoints give the zeta zero statement and the half-plane and critical-strip results for Dirichlet characters at every positive modulus. The arithmetic proof constructs the windows it needs for a prescribed character. Each exported proof uses precisely the three standard explicit axioms.}
\newcommand\StatusExternal{Andrew Sutherland and Simon \bpcode{DeDeo} each report independently building and replaying all 175,187 original modules with Lean 4.33.1. \bpcode{DeDeo} reports zero errors and crashes. Sutherland additionally records a con-leche comparison of the five targets and a lean4lean check of every original module. Five modules required higher reduction limits. The individual logs identify the names, flags and successful reruns.}
\newcommand\StatusExternalFirst{Andrew Sutherland and Simon \bpcode{DeDeo} each report independently building and replaying all 175,187 original modules with Lean 4.33.1.}
\newcommand\StatusExternalThird{Andrew Sutherland and Simon \bpcode{DeDeo} each report independently building and replaying all 175,187 original modules with Lean 4.33.1. \bpcode{DeDeo} reports zero errors and crashes. Sutherland additionally records a con-leche comparison of the five targets and a lean4lean check of every original module. Five modules required higher reduction limits. The individual logs identify the names, flags and successful reruns.}
\newcommand\StatusExternalOwn{}
\newcommand\StatusNative{At Lean 4.33.1 we checked 24 named original results and 21 reading deductions across 2,244 distinct original modules. We also checked five endpoint equivalences and an Euler-factor boundary example. The one-sided prime-window implication passed fresh Lean replay with Mathlib and a second kernel on a portable proof of 61,697 declarations. The two-input deduction of ordinary RH and Dirichlet GRH, four steps from finite character comparison to the prime-sum bound and the auxiliary-prime theorem for three changing quadratic twists all passed. \bpcode{LemmaK}, \bpcode{FixedIn} and \bpcode{TwoIn} yield one record of cohomological window constants for every cutoff exponent Ct ≥ 1. From Poitou–Tate duality, local finiteness and local Tate duality, a tower gives its cohomological bound. Adding Ferrero–Washington and \bpcode{LemmaK} yields the joint window under those same hypotheses. Eight further teaching declarations use Lean 4.35.0. At each fixed primitive character, \bpcode{TopG} and just its own cohomological estimate yield a cyclic-family bound with coefficient one. For the principal character modulo one this is the exact centered ordinary prime sum. The archive supplies the ordinary first moment unconditionally. Fresh kernel replay establishes that Schoenfeld’s RH-to-Chebyshev estimate suffices for the converse from RH to the window bound.}
\newcommand\StatusFull{Building the original library targets and Main with 48 Lake workers and persistent group objects}
\newcommand\StatusAsOf{20 November 2026, 23:51 UTC}
\newcommand\StatusAsOfShort{2026-11-20 23:51 UTC}
\newcommand\StatusJob{26647685}
\newcommand\ArchiveSHA{9631b4f020199e580ba86178cb56c6e9b09f690c2369a80dc4334aa93f49c7fc}
\newcommand\PdfSHA{3fc97784b71ddd44ff8aa0a177e99622a1c61b13cfa843a0989ee4a7a0cfeb8d}
\newcommand\MirrorURL{https://github.com/Deicyde/riemann-hypothesis}
\newcommand\MirrorName{Deicyde/riemann-hypothesis}
\newcommand\MirrorCommit{40e56bab026552e089d7afcfeff08d2defb9c005}
\newcommand\MirrorCommitShort{40e56bab}
\newcommand\Toolchain{leanprover/lean4:v4.33.1}
\newcommand\ToolchainVersion{4.33.1}
\newcommand\MathlibCommit{0df444a360eaa60ab8c11dca51a86af692955474}
\newcommand\MathlibCommitShort{0df444a3}
\newcommand\SourceCommit{348ec48748e0ea536aba7d64351b2ee2f7d116d7}
\newcommand\SourceCommitShort{348ec487}
\newcommand\ImportedAt{21 November 2026, 00:17 UTC}
\newcommand\SkeletonLean{Lean 4.35.0}
\newcommand\SkeletonBuild{passed}
\newcommand\SkeletonWarnings{0}
\newcommand\SkeletonDecls{18}
\newcommand\SkeletonAxiomFree{10}
\newcommand\SkeletonGenerated{2026-11-21 00:21:04 UTC}
\newcommand\HeroCaption{\textbf{The centered prime sum, normalized.} Computed: (G(y) − log 2\(π\))/\(\sqrt{y}\) for 10 ≤ y ≤ \(10^{7}\), where G(y) = \(\sum_{n≤y}\) \(Λ\)(n)(y/n-1) − y log y + (1+\(γ\)) y. Dotted lines: ±(2+\(γ\)-log 4\(π\)), the band that RH forces. Shaded column: one power window [Y, \(Y^{3/2}\)], Y = \(10^{4}\), with an illustrative one-sided ceiling C \(y^{θ}\) (\(θ\) = 0.55) in the same normalization.}


\hypersetup{colorlinks=true,linkcolor=ink,urlcolor=linkblue,citecolor=ink,filecolor=linkblue,
  pdftitle={\ModelTitle: an illustrated blueprint},
  pdfsubject={Components, interfaces, machine-checked implications and pinned Lean declarations of the
    announced proof (20 November 2026)},
  pdfkeywords={Riemann hypothesis, Dirichlet L-functions, Lean, blueprint, formal verification}}
\urlstyle{same}

%% ------------------------------------------------------------------ page and text layout
\setlength{\parindent}{0pt}
\setlength{\parskip}{0.55em plus 0.15em minus 0.05em}
\setlength{\jot}{4pt}
\setlist{leftmargin=1.4em,topsep=0.2em,itemsep=0.15em,parsep=0pt}
\raggedbottom
\AtBeginDocument{\color{ink}}
\captionsetup{font={small},labelfont={sf,bf},labelsep=period,justification=justified,
  singlelinecheck=true,skip=4pt,hypcap=false}
\setcounter{tocdepth}{2}

\newcommand\bpsf[1]{{\sffamily #1}}
\newcommand\bpsmallhead[1]{{\sffamily\bfseries\fontsize{7.2}{9}\selectfont\addfontfeatures{LetterSpace=6}\MakeUppercase{#1}}}
\newcommand\bpcode[1]{\texttt{#1}}
\newcommand\bphash[1]{\texttt{\footnotesize\expandafter\seqsplit\expandafter{#1}}}
\makeatletter
% break long hashes anywhere (no package needed)
\def\seqsplit#1{\bp@seqsplit#1\@nil}
\def\bp@seqsplit#1{\ifx#1\@nil\else#1\hskip0pt plus 0.3pt\relax\expandafter\bp@seqsplit\fi}
\makeatother

% section headings
\newcommand\bpseccolor{ink}
\titleformat{\section}{\sffamily\bfseries\fontsize{19}{23}\selectfont\color{ink}}
  {\textcolor{\bpseccolor}{\thesection}}{0.7em}{}
\titlespacing*{\section}{0pt}{0pt}{4mm}
\titleformat{\subsection}{\sffamily\bfseries\fontsize{13}{16}\selectfont\color{ink}}{\thesubsection}{0.6em}{}
\titlespacing*{\subsection}{0pt}{5mm}{2mm}
\titleformat{\subsubsection}[runin]{\sffamily\bfseries\fontsize{10}{12}\selectfont\color{ink}}{}{0pt}{}[.]
\titlespacing*{\subsubsection}{0pt}{3mm}{0.6em}
\titleformat{\paragraph}[runin]{\normalfont\bfseries}{}{0pt}{}
\titlespacing*{\paragraph}{0pt}{1.6mm}{0.6em}

% running heads
\pagestyle{fancy}
\fancyhf{}
\renewcommand{\headrulewidth}{0.4pt}
\renewcommand{\headrule}{{\color{rulegrey}\hrule width\headwidth height\headrulewidth}}
\fancyhead[L]{\sffamily\fontsize{7.5}{9}\selectfont\color{inksoft}\ModelTitle}
\fancyhead[R]{\sffamily\fontsize{7.5}{9}\selectfont\color{inksoft}\rightmark}
\fancyfoot[L]{\sffamily\fontsize{7}{9}\selectfont\color{inkfaint}Status as of \StatusAsOfShort\enspace
  \textperiodcentered\enspace Lean links pinned to \MirrorCommitShort}
\fancyfoot[R]{\sffamily\fontsize{8}{9}\selectfont\color{inksoft}\thepage}
\renewcommand{\sectionmark}[1]{\markright{\thesection\enspace #1}}
\fancypagestyle{plain}{\fancyhf{}\renewcommand{\headrulewidth}{0pt}%
  \fancyfoot[R]{\sffamily\fontsize{8}{9}\selectfont\color{inksoft}\thepage}}

%% ------------------------------------------------------------------ TikZ: layers and styles
\pgfdeclarelayer{panels}
\pgfdeclarelayer{cuts}
\pgfsetlayers{panels,cuts,main}
\tikzset{
  bpdiagram/.style={font=\sffamily},
  bppanelhead/.style={text=white,font=\sffamily\bfseries\fontsize{8.6}{10}\selectfont},
  bpflow/.style={-{Stealth[length=2.4mm,width=2.6mm]},line width=1.6pt,draw=black!16},
  bpcomp/.style={draw=#1,line width=0.8pt,rounded corners=1.8mm,fill=white,text width=43mm,
    minimum width=47mm,minimum height=19.5mm,inner sep=2mm,align=flush left,font=\sffamily},
  bpif/.style={-{Stealth[length=2.1mm,width=1.8mm]},line width=0.85pt,draw=#1!80!black,rounded corners=1.2mm},
  bplabel/.style={font=\sffamily\itshape\fontsize{6.3}{7.3}\selectfont,text=ink,fill=white,
    fill opacity=0.93,text opacity=1,inner xsep=0.7mm,inner ysep=0.35mm,rounded corners=0.6mm},
  bpcutII/.style={fill=tier-checked,fill opacity=0.07,draw=tier-checked,draw opacity=1,line width=1pt,
    dash pattern=on 3.4pt off 2.2pt,rounded corners=4mm},
  bpcutI/.style={fill=tier-two-kernels,fill opacity=0.09,draw=tier-two-kernels,draw opacity=1,
    line width=1.2pt,dash pattern=on 7pt off 2pt on 1.4pt off 2pt,rounded corners=3mm},
  bpcuttag/.style={fill=#1,text=white,font=\sffamily\bfseries\fontsize{6.4}{7.4}\selectfont,
    rounded corners=1.6mm,inner xsep=1.5mm,inner ysep=0.75mm},
  bpcutroute/.style={-{Stealth[length=2mm,width=1.7mm]},line width=0.95pt,draw=tier-two-kernels,
    dash pattern=on 3pt off 1.6pt},
  bplegendbox/.style={draw=#1,line width=0.7pt,rounded corners=1.8mm,fill=white,text width=95mm,
    inner sep=2.2mm,align=flush left,font=\sffamily\fontsize{6.6}{8}\selectfont},
  bptierbadgestyle/.style={bpinner,rounded corners=0.9mm,fill=#1!13,text=#1!72!black,inner xsep=1.3mm,
    inner ysep=0.7mm,font=\sffamily\bfseries\fontsize{7}{8}\selectfont},
  bpnumstyle/.style={bpinner,circle,fill=#1,text=white,inner sep=0pt},
  % nodes of pictures nested in a node of another picture inherit its options: reset them
  bpinner/.style={text width=,minimum width=0pt,minimum height=0pt,text height=,text depth=,
    draw=none,line width=0.4pt,rounded corners=0pt,inner sep=0.333em,outer sep=0pt,fill opacity=1,
    text opacity=1,rotate=0,xshift=0pt,yshift=0pt,font=},
}
\newcommand\bpfootfont{\sffamily\fontsize{6.2}{7.2}\selectfont\color{inksoft}}

%% ------------------------------------------------------------------ tiers, numbers, references
\newcommand\bptierlabel[1]{\csname bp@tierlabel@#1\endcsname}
\newcommand\bptiericon[1]{\csname bp@tiericon@#1\endcsname}
\newcommand\bptiericonplain[1]{\textcolor{tier-#1}{\bptiericon{#1}}}
\newcommand\bptierbadge[1]{\tikz[baseline=0pt]\node[anchor=base,bptierbadgestyle=tier-#1]
  {\bptiericon{#1}\enspace\bptierlabel{#1}};}
\newcommand\bptierbadgesmall[1]{\tikz[baseline=0pt]\node[anchor=base,bptierbadgestyle=tier-#1,
  font=\sffamily\bfseries\fontsize{5.7}{6.5}\selectfont,inner xsep=1mm,inner ysep=0.55mm]
  {\bptiericon{#1}\enspace\bptierlabel{#1}};}
% circled component numbers in the colour of the stratum: large (headings), small (picture), tiny (text)
\newcommand\bpnumlarge[2]{\tikz[baseline=0pt]\node[anchor=base,bpnumstyle=stratum-#1,minimum size=6.6mm,
  font=\sffamily\bfseries\fontsize{10.5}{10.5}\selectfont,yshift=-0.2mm]{#2};}
\newcommand\bpnumsmall[2]{\tikz[baseline=0pt]\node[anchor=base,bpnumstyle=stratum-#1,minimum size=4.3mm,
  font=\sffamily\bfseries\fontsize{6.8}{6.8}\selectfont]{#2};}
\newcommand\bpnumtiny[2]{\tikz[baseline=0pt]\node[anchor=base,bpnumstyle=stratum-#1,minimum size=3.5mm,
  font=\sffamily\bfseries\fontsize{5.6}{5.6}\selectfont,yshift=-0.15mm]{#2};}
\makeatletter
\newcommand\bpcompnum[1]{\edef\bp@s{\csname bp@compstrat@#1\endcsname}%
  \bpnumtiny{\bp@s}{\csname bp@compnum@#1\endcsname}}
\newcommand\bpcompref[1]{\hyperref[comp:#1]{\bpcompnum{#1}\nobreak\,\csname bp@comptitle@#1\endcsname}}
\newcommand\bppath[1]{{\ttfamily\nolinkurl{#1}}}
\newcommand\bpnoderef[1]{\hyperlink{node:#1}{\textcolor{tier-\csname bp@nodetier@#1\endcsname}%
  {\bptiericon{\csname bp@nodetier@#1\endcsname}}\,\csname bp@nodetitle@#1\endcsname}}
\makeatother
\makeatletter\newcommand\bptoc{\@starttoc{toc}}\makeatother
\newcommand\bpstep{\enspace{\color{inkfaint}$\rightarrow$}\enspace}

%% ------------------------------------------------------------------ the picture: node contents
\newcommand\bpdiagramtitle{\begin{minipage}{277mm}
  {\sffamily\bfseries\fontsize{15}{17}\selectfont The proof in one picture}\hfill
  {\sffamily\fontsize{7.4}{9}\selectfont\color{inksoft}status as of \StatusAsOf\enspace
   \textperiodcentered\enspace Lean names link to the mirror at \MirrorCommitShort}\\[1mm]
  {\fontsize{8.6}{10.4}\selectfont\itshape\color{inksoft}From the arithmetic supply (left) to the exported
   statements (right). Arrows: the \NumInterfaces{} interfaces, labelled by the proposition they carry.
   Dashed frontiers: the two machine-checked implications. Click a title for its section, a monospace
   name for its pinned source.}\end{minipage}}
\newcommand\bpnodehead[3]{\begin{tabular}{@{}p{4.4mm}@{\hspace{1.2mm}}p{37.4mm}@{}}
  \bpnumsmall{#1}{#2} & \raggedright\fontsize{7.9}{9.1}\selectfont\bfseries #3\end{tabular}\par}
\newcommand\bpnodekicker[1]{\vspace{0.7mm}{\raggedright\fontsize{6.1}{7.4}\selectfont\color{inksoft}#1\par}}
\newcommand\bpkickertext[1]{{\itshape #1}}
\newcommand\bpkickersep{\,{\color{inkfaint}\textperiodcentered}\,}
\newcommand\bpnodeexports[1]{\vspace{0.6mm}{\fontsize{5.6}{6.9}\selectfont #1\par}}
\newcommand\bpnodebadge[1]{\vspace{1mm}\bptierbadgesmall{#1}}
\newcommand\bpcutlegend[6]{%
  {\color{tier-#3}\bfseries CUT #2\enspace\bptiericon{#3}\enspace\bptierlabel{#3}}\hfill
  {\color{inksoft}\hyperref[cut:#1]{details \textrightarrow}}\\[0.5mm]
  {\fontsize{7.4}{9}\selectfont\bfseries\hyperref[cut:#1]{#4}}\\[0.9mm]
  {\fontsize{7.2}{9}\selectfont premise\enspace\(#5\)}\\[0.4mm]
  {\fontsize{7.2}{9}\selectfont\(\Longrightarrow\)\enspace #6}}
\newcommand\bptierlegendline[1]{\makebox[5.4mm][l]{\bptiericonplain{#1}}\bptierlabel{#1}}
\newcommand\bptierlegendbox[1]{{\bfseries\hyperref[sec:tiers]{Verification tiers}}\\[0.6mm]#1}
\newcommand\bpdiagramfootbox{\begin{minipage}{277mm}\raggedright\bpdiagramfoot\end{minipage}}
\newcommand\bpdiagramfoot{Numbers follow the reading order of Sections~\ref{stratum:dest}--\ref{stratum:arith}
  (destination first). Green dots: where the premises of Cut~II cross its frontier.
  (W): where the premise of Cut~I enters the window criterion; its zeta route bypasses the order of
  choices and the conductor descent. Tiers as of \StatusAsOfShort.}

%% ------------------------------------------------------------------ component pages
\newcommand\bpstratcomp[1]{\item[\bpcompnum{#1}] \hyperref[comp:#1]{\csname bp@comptitle@#1\endcsname}}
\newcommand\bpstratflow[3]{\item[] \raggedright\bpcompref{#1}\enspace\textrightarrow\enspace\bpcompref{#2}\newline
  {\color{inksoft}\emph{#3}}}
\newcommand\bpstratum[5]{%
  \clearpage
  \renewcommand\bpseccolor{stratum-#1}%
  \section{#2}\label{stratum:#1}%
  {\color{stratum-#1}\hrule height 2.2pt}\vspace{3mm}%
  {\fontsize{12.5}{16}\selectfont\itshape #3\par}\vspace{2mm}%
  \begin{minipage}[t]{0.42\linewidth}\bpsmallhead{Components}\par\vspace{0.6mm}
    \begin{itemize}[leftmargin=6mm,labelsep=2mm]#4\end{itemize}\end{minipage}\hfill
  \begin{minipage}[t]{0.55\linewidth}\bpsmallhead{Across the stratum boundary}\par\vspace{0.6mm}
    {\small\begin{itemize}[leftmargin=0mm,labelsep=0mm,itemsep=0.5mm]#5\end{itemize}}\end{minipage}\par
  \vspace{3mm}{\color{rulegrey}\hrule}\vspace{2mm}%
  \renewcommand\bpseccolor{ink}}

\newcommand\bpcomponent[5]{% id, stratum, number, title, tier
  \par\Needspace{0.32\textheight}\addvspace{7mm}%
  \phantomsection\label{comp:#1}%
  \addcontentsline{toc}{subsection}{\protect\numberline{\texorpdfstring{\protect\bpnumtiny{#2}{#3}}{#3}}#4}%
  \markright{#3\enspace #4}%
  \noindent\begin{tabular}{@{}p{8.6mm}@{}p{\dimexpr\linewidth-8.6mm\relax}@{}}
    \bpnumlarge{#2}{#3} & \raggedright\sffamily\bfseries\fontsize{15.5}{18.5}\selectfont #4
  \end{tabular}\par\nobreak\vspace{1.3mm}%
  {\color{stratum-#2}\hrule height 0.9pt}\nobreak\vspace{1.6mm}%
  \noindent\bptierbadge{#5}}
\newcommand\bpkicker[1]{\hfill{\sffamily\fontsize{7.6}{9}\selectfont\color{inksoft}#1}\par\nobreak}
\newcommand\bpgist[1]{\vspace{1.2mm}{\fontsize{12.3}{15.6}\selectfont\itshape #1\par}\vspace{0.4mm}}
\newcommand\bpdisplay[1]{{\centering\(\displaystyle #1\)\par}}
\newtcolorbox{bpstatement}[1]{enhanced,colback=stratum-#1!3,colframe=stratum-#1,boxrule=0pt,
  leftrule=2.4pt,arc=0pt,outer arc=0pt,left=3.5mm,right=3.5mm,top=3.4mm,bottom=2.6mm,
  before skip=2.5mm,after skip=3.5mm,
  overlay={\node[anchor=north east,font=\sffamily\bfseries\fontsize{6.2}{7}\selectfont,
    text=stratum-#1,inner sep=1.1mm] at (frame.north east) {STATEMENT};}}
\newcommand\bpfigure[4]{\par\vspace{2mm}\noindent\begin{minipage}{\linewidth}\centering
  \includegraphics[width=#2\linewidth]{build/fig/#1.pdf}\par
  \captionof{figure}{\textbf{#3.} #4}\label{fig:#1}\end{minipage}\par\vspace{2mm}}
\newcommand\bpfigref[2]{\par{\small\color{inksoft}Figure~\ref{fig:#1} (\emph{#2}) illustrates this
  component as well.}\par}
\newenvironment{bpleantable}[1]{\par\vspace{1.5mm}\noindent
  {\color{stratum-#1}\bpsmallhead{Lean declarations}}\enspace
  {\sffamily\fontsize{7}{9}\selectfont\color{inksoft}pinned to \MirrorName{} at \MirrorCommitShort}\par
  \nobreak\vspace{0.4mm}\small\setlength{\tabcolsep}{0pt}\renewcommand{\arraystretch}{1.18}%
  \tabularx{\linewidth}{@{}>{\raggedright\arraybackslash}X@{\hspace{3mm}}>{\raggedleft\arraybackslash}p{44mm}@{}}
  \arrayrulecolor{rulegrey}\toprule}
  {\bottomrule\endtabularx\par}
\newcommand\bpleanrow[4]{\href{#1}{#2}\ifstrempty{#4}{}{\enspace{\footnotesize\color{inksoft}— #4}}
  & \leavevmode{\sffamily\fontsize{7}{9}\selectfont\color{inksoft}#3}\\}
\newcommand\bpiface[3]{\par\noindent\hangindent=4mm\hangafter=1\bpcompref{#1}\enspace
  {\color{inksoft}\textperiodcentered}\enspace\emph{#2}\par\noindent\hspace*{4mm}{\footnotesize #3}\par\vspace{0.6mm}}
\newcommand\bpnone{{\color{inkfaint}—}}
\newcommand\bpinterfaces[2]{\par\vspace{2.4mm}\noindent
  \begin{minipage}[t]{0.485\linewidth}\bpsmallhead{Uses}\enspace{\color{inkfaint}\textleftarrow}\par
    \vspace{0.8mm}\small\ifstrempty{#1}{\bpnone{} {\footnotesize\color{inksoft}no incoming interface}}{#1}\end{minipage}\hfill
  \begin{minipage}[t]{0.485\linewidth}\bpsmallhead{Supplies}\enspace{\color{inkfaint}\textrightarrow}\par
    \vspace{0.8mm}\small\ifstrempty{#2}{\bpnone{} {\footnotesize\color{inksoft}the exported statements}}{#2}\end{minipage}\par}
\newcommand\bpnodeitem[3]{\mbox{\textcolor{tier-#2}{\bptiericon{#2}}\,\hyperlink{node:#1}{#3}}\hskip 1.1em plus 0.6em\relax}
\newcommand\bpownednodes[3]{\par\vspace{2mm}\noindent\bpsmallhead{Blueprint nodes}\enspace
  {\footnotesize\color{inksoft}#1 declarations:\enspace #2\enspace(Section~\ref{sec:table})}\par\nobreak
  \vspace{0.5mm}{\footnotesize\raggedright #3\par}}
\newtcolorbox{bpinfobox}[1]{enhanced,breakable,colback=#1!5,colframe=#1,boxrule=0pt,leftrule=2.2pt,
  arc=0pt,outer arc=0pt,left=3mm,right=3mm,top=1.6mm,bottom=1.6mm,before skip=2.5mm,after skip=1.5mm,
  fontupper=\small}
\newcommand\bpstatus[2]{\begin{bpinfobox}{tier-#1}\bpsmallhead{Status}\enspace
  {\sffamily\fontsize{7}{9}\selectfont\color{inksoft}as of \StatusAsOfShort}\hfill\bptierbadge{#1}\par
  \vspace{0.6mm}#2\end{bpinfobox}}
\newcommand\bpcheck[2]{\begin{bpinfobox}{numerics}\bpsmallhead{What remains to check}\par\vspace{0.6mm}#1%
  #2\end{bpinfobox}}
\newcommand\bpaudit[3]{\par\vspace{0.6mm}{\footnotesize\textbf{Audit criterion} (\bpnoderef{#1}):
  \emph{#2.} #3}}
\newcommand\bptlitem[4]{\par\noindent\hyperlink{tl:#1}{Entry~#1} ({#2}):\enspace #3\enspace
  {\color{inksoft}\textperiodcentered\enspace\emph{#4}}}
\newcommand\bpannouncement[1]{\par\vspace{1.6mm}\noindent\bpsmallhead{In the announcement}\par\nobreak
  \vspace{0.3mm}{\small #1\par}}

%% ------------------------------------------------------------------ legend environments (Section 1)
\newenvironment{bptierlegend}{\par\small\renewcommand{\arraystretch}{1.25}\setlength{\tabcolsep}{0pt}%
  \tabularx{\linewidth}{@{}>{\raggedright\arraybackslash}p{50mm}@{\hspace{4mm}}>{\raggedright\arraybackslash}X@{}}
  \arrayrulecolor{rulegrey}\toprule}{\bottomrule\endtabularx\par}
\newcommand\bptierrow[2]{\bptierbadge{#1} & #2\\}
\newenvironment{bpstratalegend}{\par\small\begin{itemize}[leftmargin=7mm,itemsep=1.2mm]}{\end{itemize}}
\newcommand\bpstratumrow[3]{\item[\textcolor{stratum-#1}{\rule[-0.3mm]{3.2mm}{3.2mm}}]
  \textbf{\hyperref[stratum:#1]{\csname bp@strattitle@#1\endcsname}.} #2\newline
  {\footnotesize\raggedright #3\par}}
\newenvironment{bpitineraries}{\par\small\begin{enumerate}[leftmargin=6mm,itemsep=1.4mm]}{\end{enumerate}}
\newcommand\bpitinerary[3]{\item \textbf{#1.} \emph{#2}\newline{\footnotesize #3}}

%% ------------------------------------------------------------------ cuts and Lean quotations
\newtcolorbox{bpcut}[4]{enhanced,breakable,colback=tier-#3!4,colframe=tier-#3,boxrule=0.6pt,
  leftrule=3pt,arc=1.6mm,left=3.5mm,right=3.5mm,top=2mm,bottom=2mm,before skip=4mm,after skip=3mm,
  colbacktitle=tier-#3,coltitle=white,
  title={\sffamily\bfseries Cut #2\enspace\bptiericon{#3}\enspace\bptierlabel{#3}\enspace
    \textperiodcentered\enspace #4},
  fonttitle=\fontsize{9.5}{12}\selectfont,phantom={\phantomsection\label{cut:#1}}}
\newcommand\bpcutpremise[1]{\noindent\bpsmallhead{Premise}\par\nobreak\vspace{1mm}\bpdisplay{#1}\vspace{1.4mm}}
\newcommand\bpcutconclusion[1]{\noindent\bpsmallhead{Conclusion}\par\nobreak\vspace{0.6mm}
  {\centering\(\Longrightarrow\)\enspace #1\par}\vspace{1.4mm}}
\newcommand\bpcutmeta[4]{{\small\raggedright\noindent\textbf{Lean.} #1, in \bppath{#2} (this repository;
  it compiles only against the archive, at Lean \ToolchainVersion).\par\noindent\textbf{Evidence.} #3\par
  \noindent\textbf{Components inside the cut.} #4\par}}
\newcommand\bpquotecaption[2]{\par\vspace{1mm}{\small\raggedright\color{inksoft}Quoted verbatim from
  \bppath{#1}: #2.\par}}
\newcommand\LKw[1]{\textbf{\textcolor[HTML]{1E3A8A}{#1}}}
\newcommand\LDecl[1]{\textbf{\textcolor[HTML]{B45309}{#1}}}
\newcommand\LTac[1]{\textcolor[HTML]{0E7490}{#1}}
\newcommand\LSort[1]{\textcolor[HTML]{7C2D12}{#1}}
\newcommand\LCom[1]{\textit{\textcolor[HTML]{64748B}{#1}}}
\newcommand\LStr[1]{\textcolor[HTML]{047857}{#1}}
\DefineVerbatimEnvironment{bplean}{Verbatim}{commandchars=\\\{\},fontsize=\fontsize{8.4}{10.4}\selectfont,
  numbers=left,numbersep=2.5mm,breaklines=true,breakindent=4ex,
  breaksymbolleft={\tiny\color{inkfaint}\ensuremath{\hookrightarrow}},xleftmargin=6mm}
\renewcommand\theFancyVerbLine{\sffamily\fontsize{6}{7}\selectfont\color{inkfaint}\arabic{FancyVerbLine}}
\newtcolorbox{bpleanquote}[3]{enhanced,breakable,colback=wash,colframe=rulegrey,boxrule=0.5pt,arc=1.2mm,
  left=1mm,right=2mm,top=1.2mm,bottom=1mm,before skip=2mm,after skip=3mm,
  title={\bpcode{#1}\hfill lines #2--#3},fonttitle=\sffamily\fontsize{7}{8}\selectfont,
  colbacktitle=wash,coltitle=inksoft,toptitle=0.4mm,bottomtitle=0.4mm,titlerule=0.4pt,
  titlerule style={rulegrey}}
\DefineVerbatimEnvironment{bpshell}{Verbatim}{fontsize=\fontsize{8.4}{10.6}\selectfont,
  breaklines=true,xleftmargin=3mm,frame=leftline,framerule=1.2pt,rulecolor=\color{rulegrey}}
\newenvironment{bpaxiomtable}{\par\small\renewcommand{\arraystretch}{1.16}\setlength{\tabcolsep}{0pt}%
  \tabularx{\linewidth}{@{}>{\raggedright\arraybackslash}p{80mm}@{\hspace{3mm}}>{\raggedright\arraybackslash}X@{\hspace{2mm}}c@{}}
  \arrayrulecolor{rulegrey}\toprule
  \bpsmallhead{Declaration} & \bpsmallhead{Axioms (transitive)} & \bpsmallhead{Free}\\\midrule}
  {\bottomrule\endtabularx\par}
\newcommand\bpaxiomrow[3]{#1 & \leavevmode #2 & \ifstrequal{#3}{1}{\textcolor{tier-checked}{✓}}{}\\}
\newcommand\bpnoaxioms{{\color{tier-checked}none}}

%% ------------------------------------------------------------------ the 87-declaration table
\newcommand\bpnodecl{{\footnotesize\color{inkfaint}no declaration (motivation)}}
\newenvironment{bpnodetable}{\par\footnotesize\setlength{\tabcolsep}{1.2mm}\renewcommand{\arraystretch}{1.12}%
  \begin{longtable}{@{}>{\centering\arraybackslash}p{3.2mm}>{\raggedright\arraybackslash}p{33mm}%
    >{\raggedright\arraybackslash}p{51mm}>{\raggedright\arraybackslash}p{\dimexpr\linewidth-94.8mm\relax}@{}}
  \arrayrulecolor{rulegrey}\toprule
  & \bpsmallhead{Node} & \bpsmallhead{Declaration (pinned)} & \bpsmallhead{Contract}\\\midrule\endfirsthead
  \toprule & \bpsmallhead{Node} & \bpsmallhead{Declaration (pinned)} & \bpsmallhead{Contract}\\\midrule\endhead
  \bottomrule\endfoot}
  {\end{longtable}\par}
\newcommand\bpnodegroup[5]{\multicolumn{4}{@{}>{\columncolor{stratum-#2!7}}l@{}}{\rule{0pt}{4.2mm}%
  \hyperref[comp:#1]{\bpnumtiny{#2}{#3}\enspace{\sffamily\bfseries\small #4}}\enspace
  {\sffamily\fontsize{7}{9}\selectfont\color{inksoft}#5 nodes}}\\[0.5mm]}
\newcommand\bpnoderow[6]{\leavevmode\hypertarget{node:#1}{}\textcolor{tier-#2}{\bptiericon{#2}} & #3\newline
  {\sffamily\fontsize{6}{7.5}\selectfont\color{inkfaint}#1} & #4\newline{\sffamily\fontsize{6}{7.5}\selectfont
  \color{inkfaint}#5} & #6\\}
\newcommand\bpgraph[8]{% trim of left half (right part), of right half (left part), top, bottom, widths, caption
  \clearpage\newgeometry{margin=10mm}
  \begin{landscape}\thispagestyle{empty}\centering
  \includegraphics[width=#5\linewidth,trim=0 #4 #1 #3,clip]{build/fig/blueprint-graph.pdf}\par
  \vspace{1.5mm}{\color{rulegrey}\hrule}\vspace{1.5mm}
  \includegraphics[width=#6\linewidth,trim=#2 #4 0 #3,clip]{build/fig/blueprint-graph.pdf}\par
  \vspace{1mm}\captionof{figure}{\textbf{#7}, in two halves (left: the suppliers; right: continued
    towards the endpoints). #8 The node titles are those of the table above.}\label{fig:graph}
  \end{landscape}\restoregeometry}

%% ------------------------------------------------------------------ timeline, pins
\tikzset{
  bptimeline/.style={font=\sffamily},
  bptlaxis/.style={line width=0.8pt,draw=inksoft},
  bptlmonth/.style={font=\sffamily\bfseries\fontsize{7}{8}\selectfont,text=inksoft},
  bptlbreak/.style={draw=none,fill=white,minimum width=3mm,minimum height=3mm,
    path picture={\draw[inksoft,line width=0.6pt] ([xshift=-0.8mm]path picture bounding box.south)
    -- ++(1mm,3mm) ([xshift=0.3mm]path picture bounding box.south) -- ++(1mm,3mm);}},
  bptltick/.style={draw=inksoft,line width=0.6pt},
  bptlday/.style={font=\sffamily\fontsize{6.4}{7}\selectfont,text=inksoft,inner sep=0.6mm},
  bptlann/.style={draw=stratum-dest,line width=0.9pt,dash pattern=on 2pt off 1.4pt},
  bptlannlabel/.style={font=\sffamily\bfseries\fontsize{6.4}{7}\selectfont,text=stratum-dest},
  bptldot/.style={circle,fill=#1,text=white,minimum size=4.2mm,inner sep=0pt,
    font=\sffamily\bfseries\fontsize{6}{6}\selectfont},
  bptlfrom/.style={-{Stealth[length=1.4mm]},draw=#1,line width=0.8pt},
}
\newcommand\bptldotinline[2]{\tikz[baseline=0pt]\node[anchor=base,bptldot=#1,yshift=-0.2mm]{#2};}
\newcommand\bpnocomp{{\footnotesize\color{inkfaint}historical; no component}}
\newenvironment{bptimelinetable}{\par\footnotesize\setlength{\tabcolsep}{1.3mm}\renewcommand{\arraystretch}{1.1}%
  \begin{longtable}{@{}>{\centering\arraybackslash}p{4.5mm}>{\raggedright\arraybackslash}p{17mm}%
    >{\raggedright\arraybackslash}p{\dimexpr\linewidth-74mm\relax}>{\raggedright\arraybackslash}p{40mm}@{}}
  \arrayrulecolor{rulegrey}\toprule
  & \bpsmallhead{Date} & \bpsmallhead{Entry, relation and assessment} & \bpsmallhead{Component}\\\midrule\endfirsthead
  \toprule & \bpsmallhead{Date} & \bpsmallhead{Entry, relation and assessment} & \bpsmallhead{Component}\\\midrule\endhead
  \bottomrule\endfoot}
  {\end{longtable}\par}
\newcommand\bptlrow[9]{\leavevmode\hypertarget{tl:#1}{}\bptldotinline{#9}{#1} & \bpcode{#2}\newline
  {\fontsize{6.4}{7.6}\selectfont\color{inksoft}#3} & \textbf{#4}\enspace{\color{inksoft}(page~#5)}\newline
  \emph{#6.}\enspace{\color{inksoft}#8} & #7\\[0.8mm]}
\newenvironment{bppintable}{\par\small\renewcommand{\arraystretch}{1.25}\setlength{\tabcolsep}{0pt}%
  \tabularx{\linewidth}{@{}>{\raggedright\arraybackslash}p{56mm}@{\hspace{4mm}}>{\raggedright\arraybackslash}X@{}}
  \arrayrulecolor{rulegrey}\toprule}{\bottomrule\endtabularx\par}
\newcommand\bppinrow[2]{#1 & #2\\}
\newenvironment{bpsourcelist}{\begin{itemize}\small}{\end{itemize}}
\newenvironment{bpprovenance}{\par\footnotesize\renewcommand{\arraystretch}{1.18}\setlength{\tabcolsep}{0pt}%
  \tabularx{\linewidth}{@{}>{\raggedright\arraybackslash}p{62mm}@{\hspace{3mm}}>{\raggedright\arraybackslash}X@{\hspace{3mm}}>{\raggedleft\arraybackslash}p{27mm}@{}}
  \arrayrulecolor{rulegrey}\toprule
  \bpsmallhead{In this repository} & \bpsmallhead{Imported from the reading project} & \bpsmallhead{SHA-256 (prefix)}\\\midrule}
  {\bottomrule\endtabularx\par}
\newcommand\bpprovrow[3]{#1 & #2 & \bpcode{#3}\\}

%% =================================================================================== document
\begin{document}

%% ------------------------------------------------------------------ title page
\newgeometry{hmargin=20mm,top=18mm,bottom=16mm}
\begin{titlepage}
\noindent\begin{tikzpicture}[x=1mm,y=1mm]
  \foreach \s [count=\i from 0] in {arith,est,ana,dest} {\fill[stratum-\s] (\i*42.5,0) rectangle ++(42.5,1.6);}
\end{tikzpicture}\par\vspace{5mm}
{\sffamily\bfseries\fontsize{8.5}{10}\selectfont\addfontfeatures{LetterSpace=12}\color{inksoft}
  AN ILLUSTRATED BLUEPRINT\par}\vspace{3mm}
{\sffamily\bfseries\fontsize{23}{27.5}\selectfont\color{ink}The architecture of the announced proof\\
  of the Riemann hypothesis\par}\vspace{3.5mm}
{\fontsize{11.8}{14.8}\selectfont\itshape\color{inksoft}Levent Alpöge's announcement of 20~November 2026
  and its Lean archive of \ModuleCount{} modules: \NumComponents{} components in \NumStrataWord{} strata,
  \NumInterfaces{} interfaces, \NumCutsWord{} machine-checked implications and \NumDecls{} pinned declarations\par}
\vspace{2.5mm}
{\sffamily\fontsize{7.4}{9.4}\selectfont\color{inksoft}
Companion: the interactive atlas \bpcode{atlas/index.html} in this repository (\bpcode{riemann-rev}).
Data: \bpcode{data/model.json}, built from the reading project at \SourceCommitShort{} (imported
\ImportedAt). Lean sources pinned to \MirrorName{} at \MirrorCommitShort{} (archive toolchain
\Toolchain, Mathlib \MathlibCommitShort). All times UTC.\par}
\vspace{3.5mm}
{\centering\includegraphics[width=0.9\linewidth,trim=0 22 0 30,clip]{build/fig/hero.pdf}\par}\vspace{1.2mm}
{\sffamily\fontsize{6.6}{8}\selectfont\color{inksoft}\HeroCaption\par}
\vspace{2.5mm}
\begin{tcolorbox}[enhanced,colback=wash,colframe=tier-reported,boxrule=0pt,leftrule=3pt,arc=0pt,
  outer arc=0pt,left=4mm,right=4mm,top=2.4mm,bottom=2.4mm,fontupper=\footnotesize]
\bpsmallhead{Status}\enspace{\sffamily\fontsize{7.4}{9}\selectfont\color{inksoft}as of \StatusAsOf{}
  (headline of the reading project's status record)}\par\vspace{1.4mm}
{\fontsize{10.6}{13.2}\selectfont\itshape ``\StatusHeadline''\par}\vspace{1mm}
\renewcommand{\arraystretch}{1.1}\setlength{\tabcolsep}{0pt}
\begin{tabularx}{\linewidth}{@{}p{47mm}@{\hspace{3mm}}>{\raggedright\arraybackslash}X@{}}
\textbf{Announced proof} & Of RH and of GRH for Dirichlet characters, in a four-page announcement with a
  Lean archive. This document maps its architecture; it does not assert that the Riemann hypothesis is
  proved.\\
\bptierbadge{\CutITier}\newline\bptierbadge{\CutIITier} & \textbf{Machine-checked implications}, with
  their remaining inputs as explicit premises (Section~\ref{sec:cuts}). Cut~I: \CutITitle{}
  (\MakeLowercase{\bptierlabel{\CutITier}}). Cut~II: \CutIITitle{} (\MakeLowercase{\bptierlabel{\CutIITier}}).\\
\bptierbadge{reported} & \textbf{Externally reported full replay.} As recorded by the reading project:
  ``\StatusExternalThird'' These are an external builder's results, not checks made for this document.
  The reading project adds: ``\StatusExternalOwn''\\
\bptierbadge{read} & \textbf{Source-read.} The producers of the arithmetic inputs were located at pinned
  source lines and read, not separately replayed.\\
\end{tabularx}
\end{tcolorbox}
\end{titlepage}
\restoregeometry

%% ------------------------------------------------------------------ contents
\pagestyle{plain}
{\sffamily\bfseries\fontsize{19}{23}\selectfont Contents\par}\vspace{3mm}
{\small\bptoc}
\clearpage
\pagestyle{fancy}

%% ------------------------------------------------------------------ 1. how to read
\section{How to read this blueprint}\label{sec:read}

This is a map of an \emph{announced} proof. Levent Alpöge's four-page announcement of 20~November~2026
claims proofs of the Riemann hypothesis and of the generalized Riemann hypothesis for Dirichlet
characters, accompanied by a Lean archive of \ModuleCount{} modules. The map divides the argument into
\NumComponents{} mathematical components in \NumStrataWord{} strata, joined by \NumInterfaces{} \emph{interfaces}: the
propositions through which the output of one component becomes the input of another. For each component
it gives a plain-language gist, a precise statement in standard notation, an explanation of the idea,
the pinned Lean declarations, a verification tier with its status, and what remains to be checked.

The statements are paraphrases, in standard notation, of the reading project's chapters and of the
archive's source lines they cite; nothing here is a new proof, and this repository does not re-check the
archive. Two implications have been machine-checked with their remaining inputs as explicit hypotheses
(Section~\ref{sec:cuts}). A full replay of the archive has been reported by an external builder; the
reading project records its own checks separately, and its status record of \StatusAsOf{} is quoted in
Section~\ref{sec:status}. Nothing in this document asserts that the Riemann hypothesis is proved.

\paragraph{Reading order.} The components are numbered as in the model, from the destination backwards:
the \NumExportsWord{} exported statements and the classical reductions to the half-plane theorem
(Section~\ref{stratum:dest}), the analytic engine that turns one-sided prime-sum bounds into zero-free
half-planes (Section~\ref{stratum:ana}), the estimate that produces those bounds
(Section~\ref{stratum:est}), and the arithmetic supply on which the estimate rests
(Section~\ref{stratum:arith}). The picture of Section~\ref{sec:picture} runs the other way, from the
arithmetic on the left to the endpoints on the right.

\paragraph{Links.} A monospace name such as \bpcode{Deduction.poleData} links to its source lines in the
public mirror \MirrorName{}, pinned at commit \MirrorCommitShort{}, whose tree is identical to the archive;
circled numbers and node titles link inside this document.

\subsection{Verification tiers}\label{sec:tiers}
Every component and every blueprint node carries one tier. The meanings are those of the model; node
tiers are computed from the receipts of the reading project.
% Generated by scripts/build_blueprint_tex.py from data/model.json -- do not edit by hand.
\begin{bptierlegend}
\bptierrow{two-kernels}{Accepted by two independent proof checkers in runs of the reading project: a Lean replay or verified-mode comparison, and the Rust kernel NanoDA. The ledger names the runs and their receipts.}
\bptierrow{checked}{Named declarations passed type, transitive-axiom (only \bpcode{propext}, \bpcode{Classical.\allowbreak{}choice}, \bpcode{Quot.\allowbreak{}sound}) and kernel-replay checks at Lean 4.33.1.}
\bptierrow{partial}{Some of the component's named declarations are kernel-checked; the producers of its remaining inputs are source-read.}
\bptierrow{covered}{Called (directly or transitively) by a kernel-checked declaration, so it was elaborated and kernel-checked as part of that closure and its axioms are within that declaration's report. For a definition this only means that it typechecks: a proposition it defines (for instance \bpcode{VWinAt}) can still be a hypothesis.}
\bptierrow{read}{Statement and proof term located at pinned source lines and read; not replayed as a named declaration by the reading project's audit. It is part of the complete proof, whose replays are listed in the ledger.}
\bptierrow{reported}{Checked in a full replay of the archive reported by an external builder; the reading project has not (yet) recorded a run of its own for it.}
\bptierrow{motivation}{A heuristic or historical idea from the announcement; no theorem is attached to it.}
\end{bptierlegend}


\subsection{The \NumStrataWord{} strata}
% Generated by scripts/build_blueprint_tex.py from data/model.json -- do not edit by hand.
\begin{bpstratalegend}
\bpstratumrow{arith}{Galois representations, theta lifts and Iwasawa-type cohomology bounds that are meant to produce the families on which a prime sum is small.}{\bpcompref{cohomology}, \bpcompref{primes}, \bpcompref{galois}, \bpcompref{inputs}}
\bpstratumrow{est}{A finite character comparison turns geometric control of a chain of theta lifts of length about \(\sqrt{x}\) into an upper bound for a weighted prime sum.}{\bpcompref{comparison}, \bpcompref{family}, \bpcompref{hosts}}
\bpstratumrow{ana}{A localized Landau oscillation argument converts one-sided bounds on whole power windows into a zero-free half-plane.}{\bpcompref{statistic}, \bpcompref{residues}, \bpcompref{criterion}, \bpcompref{choreography}}
\bpstratumrow{dest}{Euler factors, conductors and functional equations turn the half-plane theorem into the five exported statements, including Mathlib's \href{https://github.com/leanprover-community/mathlib4/blob/0df444a360eaa60ab8c11dca51a86af692955474/Mathlib/NumberTheory/LSeries/RiemannZeta.lean\#L182-L186}{\bpcode{RiemannHypothesis}}.}{\bpcompref{endpoints}, \bpcompref{reflection}, \bpcompref{descent}}
\end{bpstratalegend}


\subsection{Ways into the blueprint}
The reading project proposes \NumItinerariesWord{} itineraries through the blueprint graph; each step links to the
declaration table of Section~\ref{sec:table}.
% Generated by scripts/build_blueprint_tex.py from data/model.json -- do not edit by hand.
\begin{bpitineraries}
\bpitinerary{Read the last page of the proof first}{How does a right-half-plane estimate imply ordinary RH?}{\bpnoderef{main-rh} \bpstep \bpnoderef{rh-all} \bpstep \bpnoderef{rh-every} \bpstep \bpnoderef{rh-right}}
\bpitinerary{Find the mathematical estimate}{How wide must the intervals be, and where does the loss in the exponent come from?}{\bpnoderef{zero-top} \bpstep \bpnoderef{window-consumer} \bpstep \bpnoderef{prime-bound} \bpstep \bpnoderef{core-gm} \bpstep \bpnoderef{shape} \bpstep \bpnoderef{family-design}}
\bpitinerary{Find who supplies the arithmetic}{Which assertion has to produce the compatible Galois, theta and cohomology data?}{\bpnoderef{every} \bpstep \bpnoderef{from-row} \bpstep \bpnoderef{row-prod} \bpstep \bpnoderef{lit} \bpstep \bpnoderef{seven} \bpstep \bpnoderef{vwin-row} \bpstep \bpnoderef{four-produce}}
\bpitinerary{Follow a singular twist and an arbitrary finite prime sum}{Which bound controls a twist of unbounded abstract length, and how can one estimate work over every allowed finite set of coefficient primes?}{\bpnoderef{four-produce} \bpstep \bpnoderef{tower-fixed-growth} \bpstep \bpnoderef{four-use} \bpstep \bpnoderef{tower-window-from-lemma-k} \bpstep \bpnoderef{tower-bound-from-rank} \bpstep \bpnoderef{tower-operator-bound} \bpstep \bpnoderef{tower-window-constants} \bpstep \bpnoderef{counts-i} \bpstep \bpnoderef{counts-common-integer} \bpstep \bpnoderef{counts-window} \bpstep \bpnoderef{counts-joint} \bpstep \bpnoderef{counts-pointwise} \bpstep \bpnoderef{counts-inside} \bpstep \bpnoderef{counts-outside}}
\bpitinerary{Recover the centering constant in the ordinary zeta window}{Why is the linear correction (1+\(γ\))x?}{\bpnoderef{moment-target} \bpstep \bpnoderef{moment-theorem} \bpstep \bpnoderef{moment-calculation} \bpstep \bpnoderef{moment-majorant} \bpstep \bpnoderef{moment-dominated} \bpstep \bpnoderef{moment-laurent}}
\end{bpitineraries}


%% ------------------------------------------------------------------ 2. the picture
\clearpage
\newgeometry{margin=10mm}
\begin{landscape}
\thispagestyle{empty}
\refstepcounter{section}\phantomsection\label{sec:picture}%
\addcontentsline{toc}{section}{\protect\numberline{\thesection}The proof in one picture}%
\noindent\makebox[\linewidth][l]{% Generated by scripts/build_blueprint_tex.py from data/model.json -- do not edit by hand.
\begin{tikzpicture}[x=1mm,y=1mm,bpdiagram]
\useasboundingbox (0,0) rectangle (277.0,188.0);
\node[anchor=north west,inner sep=0pt] at (0,188.0) {\bpdiagramtitle};
\begin{pgfonlayer}{panels}
\fill[stratum-arith!4,rounded corners=2.5mm] (0.0,43.5) rectangle (52.0,168.0);
\fill[stratum-arith,rounded corners=2.5mm] (0.0,161.0) rectangle (52.0,168.0);
\node[bppanelhead] at (26.0,164.5) {\hyperref[stratum:arith]{Arithmetic supply}};
\draw[bpflow] (53.2,164.5) -- (73.8,164.5);
\fill[stratum-est!4,rounded corners=2.5mm] (75.0,43.5) rectangle (127.0,168.0);
\fill[stratum-est,rounded corners=2.5mm] (75.0,161.0) rectangle (127.0,168.0);
\node[bppanelhead] at (101.0,164.5) {\hyperref[stratum:est]{The estimate}};
\draw[bpflow] (128.2,164.5) -- (148.8,164.5);
\fill[stratum-ana!4,rounded corners=2.5mm] (150.0,43.5) rectangle (202.0,168.0);
\fill[stratum-ana,rounded corners=2.5mm] (150.0,161.0) rectangle (202.0,168.0);
\node[bppanelhead] at (176.0,164.5) {\hyperref[stratum:ana]{The analytic engine}};
\draw[bpflow] (203.2,164.5) -- (223.8,164.5);
\fill[stratum-dest!4,rounded corners=2.5mm] (225.0,43.5) rectangle (277.0,168.0);
\fill[stratum-dest,rounded corners=2.5mm] (225.0,161.0) rectangle (277.0,168.0);
\node[bppanelhead] at (251.0,164.5) {\hyperref[stratum:dest]{The destination}};
\end{pgfonlayer}
\begin{pgfonlayer}{cuts}
\path[bpcutII] (63.5,41.0) -- (63.5,120.5) -- (138.5,120.5) -- (138.5,171.0) -- (213.5,171.0) -- (213.5,108.5) -- (279.5,108.5) -- (279.5,41.0) -- cycle;
\path[bpcutI] (150.0,160.0) -- (202.0,160.0) -- (202.0,101.2) -- (277.0,101.2) -- (277.0,75.2) -- (150.0,75.2) -- cycle;
\end{pgfonlayer}
\node[bpcuttag=tier-checked,anchor=west] at (139.8,171.0) {\hyperref[cut:cut-dirichlet]{CUT II\enspace\bptiericonplain{checked}}};
\node[bpcuttag=tier-two-kernels,anchor=west] at (151.2,160.0) {\hyperref[cut:cut-zeta]{CUT I\enspace\bptiericonplain{two-kernels}}};
\node[bpcomp=stratum-dest,anchor=north] (c-endpoints) at (251.0,158.5) {%
  \bpnodehead{dest}{1}{\hyperref[comp:endpoints]{The five endpoints}}
  \bpnodekicker{\href{https://github.com/Deicyde/riemann-hypothesis/blob/40e56bab026552e089d7afcfeff08d2defb9c005/src/Main.lean}{\bpcode{Main.\allowbreak{}lean}}}
  \bpnodeexports{\href{https://github.com/Deicyde/riemann-hypothesis/blob/40e56bab026552e089d7afcfeff08d2defb9c005/src/Main.lean\#L4-L4}{\bpcode{riemannHypothesis}}\\\href{https://github.com/Deicyde/riemann-hypothesis/blob/40e56bab026552e089d7afcfeff08d2defb9c005/src/Main.lean\#L6-L6}{\bpcode{riemannZeta\_\allowbreak{}eq\_\allowbreak{}zero\_\allowbreak{}re}}\\\href{https://github.com/Deicyde/riemann-hypothesis/blob/40e56bab026552e089d7afcfeff08d2defb9c005/src/Main.lean\#L8-L8}{\bpcode{LFunction\_\allowbreak{}ne\_\allowbreak{}zero\_\allowbreak{}of\_\allowbreak{}half\_\allowbreak{}lt\_\allowbreak{}re}}\\\href{https://github.com/Deicyde/riemann-hypothesis/blob/40e56bab026552e089d7afcfeff08d2defb9c005/src/Main.lean\#L10-L10}{\bpcode{LFunction\_\allowbreak{}eq\_\allowbreak{}zero\_\allowbreak{}re}}\\\href{https://github.com/Deicyde/riemann-hypothesis/blob/40e56bab026552e089d7afcfeff08d2defb9c005/src/Main.lean\#L12-L12}{\bpcode{LFunction\_\allowbreak{}eq\_\allowbreak{}zero\_\allowbreak{}re\_\allowbreak{}of\_\allowbreak{}isPrimitive}}}
  \bpnodebadge{two-kernels}};
\node[bpcomp=stratum-dest] (c-reflection) at (251.0,88.0) {%
  \bpnodehead{dest}{2}{\hyperref[comp:reflection]{Reflection}}
  \bpnodekicker{\bpkickertext{functional equations}}
  \bpnodebadge{checked}};
\node[bpcomp=stratum-dest] (c-descent) at (251.0,58.5) {%
  \bpnodehead{dest}{3}{\hyperref[comp:descent]{Conductor descent}}
  \bpnodekicker{\bpkickertext{Euler factors}}
  \bpnodebadge{checked}};
\node[bpcomp=stratum-ana] (c-statistic) at (176.0,147.0) {%
  \bpnodehead{ana}{4}{\hyperref[comp:statistic]{The centered prime sum}}
  \bpnodekicker{\bpkickertext{what must be bounded}}
  \bpnodebadge{checked}};
\node[bpcomp=stratum-ana] (c-residues) at (176.0,117.5) {%
  \bpnodehead{ana}{5}{\hyperref[comp:residues]{Mellin transform and positive residues}}
  \bpnodekicker{\href{https://github.com/Deicyde/riemann-hypothesis/blob/40e56bab026552e089d7afcfeff08d2defb9c005/src/Development/AutomorphicArithmetic/DedekindLFunctions/LandauTheorems/poleData.lean\#L248-L259}{\bpcode{poleData}}}
  \bpnodebadge{checked}};
\node[bpcomp=stratum-ana] (c-criterion) at (176.0,88.0) {%
  \bpnodehead{ana}{6}{\hyperref[comp:criterion]{The window criterion}}
  \bpnodekicker{\bpkickertext{localized Landau}}
  \bpnodebadge{two-kernels}};
\node[bpcomp=stratum-ana] (c-choreography) at (176.0,58.5) {%
  \bpnodehead{ana}{7}{\hyperref[comp:choreography]{The order of choices}}
  \bpnodekicker{\bpkickertext{\(θ\) after the zero, \(ε\) after \(θ\)}}
  \bpnodebadge{checked}};
\node[bpcomp=stratum-est] (c-comparison) at (101.0,102.8) {%
  \bpnodehead{est}{8}{\hyperref[comp:comparison]{The finite character comparison}}
  \bpnodekicker{\href{https://github.com/Deicyde/riemann-hypothesis/blob/40e56bab026552e089d7afcfeff08d2defb9c005/src/Development/AutomorphicArithmetic/PrimitiveCosets/CoreGm.lean\#L12-L74}{\bpcode{core\_\allowbreak{}gm}}\bpkickersep \href{https://github.com/Deicyde/riemann-hypothesis/blob/40e56bab026552e089d7afcfeff08d2defb9c005/src/Development/AutomorphicArithmetic/DedekindLFunctions/LandauTheorems/ChCore.lean\#L343-L351}{\bpcode{ch\_\allowbreak{}core}}\bpkickersep \href{https://github.com/Deicyde/riemann-hypothesis/blob/40e56bab026552e089d7afcfeff08d2defb9c005/src/Development/AutomorphicArithmetic/DedekindLFunctions/LandauTheorems/ChTop.lean\#L37-L72}{\bpcode{ch\_\allowbreak{}top}}}
  \bpnodebadge{checked}};
\node[bpcomp=stratum-est] (c-family) at (101.0,132.2) {%
  \bpnodehead{est}{9}{\hyperref[comp:family]{The geometric family and the theta chain}}
  \bpnodekicker{\href{https://github.com/Deicyde/riemann-hypothesis/blob/40e56bab026552e089d7afcfeff08d2defb9c005/src/Development/AutomorphicArithmetic/CharacterBoxes/DesWin.lean\#L11-L31}{\bpcode{topDesWin}}\bpkickersep \href{https://github.com/Deicyde/riemann-hypothesis/blob/40e56bab026552e089d7afcfeff08d2defb9c005/src/Development/AbelianVarietyTheory/Definitions/GlobalTwisting/PairRowIffQdRows.lean\#L28-L32}{\bpcode{RankVal}}}
  \bpnodebadge{partial}};
\node[bpcomp=stratum-est] (c-hosts) at (101.0,73.2) {%
  \bpnodehead{est}{10}{\hyperref[comp:hosts]{Hosts for every character}}
  \bpnodekicker{\href{https://github.com/Deicyde/riemann-hypothesis/blob/40e56bab026552e089d7afcfeff08d2defb9c005/src/Development/AutomorphicArithmetic/CharacterBoxes/SupplyChi.lean\#L8-L16}{\bpcode{supplyChi}}}
  \bpnodebadge{checked}};
\node[bpcomp=stratum-arith] (c-cohomology) at (28.0,88.0) {%
  \bpnodehead{arith}{11}{\hyperref[comp:cohomology]{The cohomological window}}
  \bpnodekicker{\href{https://github.com/Deicyde/riemann-hypothesis/blob/40e56bab026552e089d7afcfeff08d2defb9c005/src/Development/AutomorphicArithmetic/CharacterBoxes/ExistsPageN.lean\#L44-L47}{\bpcode{VWinAt ψ 401}}}
  \bpnodebadge{partial}};
\node[bpcomp=stratum-arith] (c-primes) at (28.0,58.5) {%
  \bpnodehead{arith}{12}{\hyperref[comp:primes]{Auxiliary primes and Behrend towers}}
  \bpnodekicker{\href{https://github.com/Deicyde/riemann-hypothesis/blob/40e56bab026552e089d7afcfeff08d2defb9c005/src/Development/AutomorphicArithmetic/MeanSquareBounds/pairEx.lean\#L112-L169}{\bpcode{QH.\allowbreak{}pairEx}}\bpkickersep \href{https://github.com/Deicyde/riemann-hypothesis/blob/40e56bab026552e089d7afcfeff08d2defb9c005/src/Formalism/FrobeniusOrbits/FrobeniusOrbitsProofs/ThickExists.lean\#L181-L187}{\bpcode{hostThickExists}}}
  \bpnodebadge{checked}};
\node[bpcomp=stratum-arith] (c-galois) at (28.0,147.0) {%
  \bpnodehead{arith}{13}{\hyperref[comp:galois]{The Galois row}}
  \bpnodekicker{\href{https://github.com/Deicyde/riemann-hypothesis/blob/40e56bab026552e089d7afcfeff08d2defb9c005/src/Development/AutomorphicArithmetic/Foundations/GaloisRepresentations/TypeConditions/WeightsConjBis.lean\#L69-L73}{\bpcode{Row0b}}}
  \bpnodebadge{read}};
\node[bpcomp=stratum-arith] (c-inputs) at (28.0,117.5) {%
  \bpnodehead{arith}{14}{\hyperref[comp:inputs]{Seven inputs at the decisive call}}
  \bpnodekicker{\href{https://github.com/Deicyde/riemann-hypothesis/blob/40e56bab026552e089d7afcfeff08d2defb9c005/src/Development/AutomorphicArithmetic/ThetaCorrespondence/ThetaPairLifts/ThirdPlaceSeries/EveryCharacterOfLiteratureOfVWin.lean\#L810-L813}{\bpcode{every\_\allowbreak{}character\_\allowbreak{}of\_\allowbreak{}literature\_\allowbreak{}of\_\allowbreak{}vWin}}}
  \bpnodebadge{read}};
\draw[bpif=stratum-arith] ([xshift=-10mm]c-galois.south) -- ([xshift=-10mm]c-inputs.north) node[bplabel,right,midway,xshift=0.6mm] {\bpcode{Row0b} (two fields)};
\draw[bpif=stratum-arith] (c-galois.west) -- ++(-5.2,0) |- (c-cohomology.west) node[bplabel,pos=0.25,rotate=90,yshift=2.6mm] {\bpcode{Row0b} → window};
\draw[bpif=stratum-arith] ([xshift=-10mm]c-primes.north) -- ([xshift=-10mm]c-cohomology.south) node[bplabel,right,midway,xshift=0.6mm] {prime pairs, tower};
\draw[bpif=stratum-arith] (c-inputs.east) -- ++(5.5,0) |- (c-family.west) node[bplabel,pos=0.75,above,text width=17mm,align=center] {definite-side data};
\draw[bpif=stratum-arith] (c-cohomology.east) -- ++(5.5,0) |- (c-hosts.west) node[bplabel,pos=0.72,below=0.4mm] {\bpcode{VWinAt} \(ψ\) 401};
\draw[bpif=stratum-est] ([xshift=-12mm]c-family.south) -- ([xshift=-12mm]c-comparison.north) node[bplabel,right,pos=0.66,xshift=0.6mm] {\bpcode{TopG} (\(e_{D}\) = 2, \(e_{L}\) = 0)};
\draw[bpif=stratum-est] ([xshift=-12mm]c-hosts.north) -- ([xshift=-12mm]c-comparison.south) node[bplabel,right,midway,xshift=0.6mm] {hosts, width 3/2};
\draw[bpif=stratum-est] (c-comparison.east) to[out=0,in=180,looseness=0.9] node[bplabel,pos=0.47,text width=15mm,align=center] {\(G_{P}\) ≤ \(x^{1/2+9ε}\) on windows} (c-choreography.west);
\draw[bpif=stratum-ana] ([xshift=-12mm]c-statistic.south) -- ([xshift=-12mm]c-residues.north) node[bplabel,right,midway,xshift=0.6mm] {Mellin transform};
\draw[bpif=stratum-ana] ([xshift=-12mm]c-residues.south) -- ([xshift=-12mm]c-criterion.north) node[bplabel,right,midway,xshift=0.6mm] {pole data, S(\(ρ\)) > 0};
\draw[bpif=stratum-ana] ([xshift=-12mm]c-criterion.south) -- ([xshift=-12mm]c-choreography.north) node[bplabel,right,pos=0.66,xshift=0.6mm] {criterion at each level};
\draw[bpif=stratum-ana] (c-choreography.east) -- (c-descent.west) node[bplabel,midway,above=0.4mm,text width=16mm,align=center] {\bpcode{ZeroFreeAll} (primitive)};
\draw[bpif=stratum-dest] ([xshift=-12mm]c-descent.north) -- ([xshift=-12mm]c-reflection.south) node[bplabel,right,pos=0.38,xshift=0.6mm] {every \(χ\), Re s > ½};
\draw[bpif=stratum-dest] ([xshift=14mm]c-reflection.north) -- ([xshift=14mm]c-endpoints.south) node[bplabel,left,pos=0.62,xshift=-0.6mm] {critical line};
\fill[tier-checked] ([xshift=-12mm]c-family.south |- 0,120.5) circle (0.9);
\fill[tier-checked] (63.5,73.2) circle (0.9);
\draw[bpcutroute] (c-criterion.east) -- (c-reflection.west) node[bplabel,midway,above=0.3mm,text=tier-two-kernels!80!black] {\(\zeta(s)\neq0,\ \operatorname{Re}s>\tfrac12\)};
\draw[bpcutroute] (143.2,92.0) -- ([yshift=4mm]c-criterion.west) node[bplabel,pos=0.0,anchor=south west,yshift=0.4mm,xshift=-1mm,text=tier-two-kernels!80!black] {(W)};
\node[bplegendbox=tier-two-kernels,anchor=south west] at (0,8.0) {\bpcutlegend{cut-zeta}{I}{two-kernels}{One-sided prime windows imply RH}{\forall\theta\in(\tfrac12,1)\ \exists C\ \forall Y_0\ \exists Y\ge Y_0\ \forall y\in[Y,Y^{3/2}]:\ G(y)\le C\,y^{\theta}}{Mathlib's \href{https://github.com/leanprover-community/mathlib4/blob/0df444a360eaa60ab8c11dca51a86af692955474/Mathlib/NumberTheory/LSeries/RiemannZeta.lean\#L182-L186}{\bpcode{RiemannHypothesis}}}};
\node[bplegendbox=tier-checked,anchor=south west] at (99.5,8.0) {\bpcutlegend{cut-dirichlet}{II}{checked}{\bpcode{TopG} and the cohomological windows imply RH and Dirichlet GRH}{\mathrm{TopG}\ \wedge\ \forall\psi\ \text{primitive}:\ \mathrm{VWinAt}\ \psi\ 401}{RH ∧ half-plane nonvanishing ∧ critical-strip zeros, for every positive modulus}};
\node[bplegendbox=black!45,anchor=south east,text width=64mm] at (277.0,8.0) {\bptierlegendbox{\bptierlegendline{two-kernels}\\[0.35mm]\bptierlegendline{checked}\\[0.35mm]\bptierlegendline{partial}\\[0.35mm]\bptierlegendline{covered}\\[0.35mm]\bptierlegendline{read}\\[0.35mm]\bptierlegendline{reported}\\[0.35mm]\bptierlegendline{motivation}}};
\node[anchor=south west,inner sep=0pt,font=\bpfootfont] at (0,0) {\bpdiagramfootbox};
\end{tikzpicture}
}
\end{landscape}
\restoregeometry

%% ------------------------------------------------------------------ 3-6. strata and components
% Generated by scripts/build_blueprint_tex.py from data/model.json -- do not edit by hand.
\bpstratum{dest}{The destination}{Euler factors, conductors and functional equations turn the half-plane theorem into the five exported statements, including Mathlib's \href{https://github.com/leanprover-community/mathlib4/blob/0df444a360eaa60ab8c11dca51a86af692955474/Mathlib/NumberTheory/LSeries/RiemannZeta.lean\#L182-L186}{\bpcode{RiemannHypothesis}}.}{\bpstratcomp{endpoints}\bpstratcomp{reflection}\bpstratcomp{descent}}{\bpstratflow{choreography}{descent}{\bpcode{ZeroFreeAll} (primitive)}}
\bpcomponent{endpoints}{dest}{1}{The five endpoints}{two-kernels}
\bpkicker{\href{https://github.com/Deicyde/riemann-hypothesis/blob/40e56bab026552e089d7afcfeff08d2defb9c005/src/Main.lean}{\bpcode{Main.\allowbreak{}lean}}}
\bpgist{Ordinary RH in Mathlib's own formulation, and Dirichlet GRH for every character of every positive modulus.}
\begin{bpstatement}{dest}
\bpdisplay{\begin{gathered}
\zeta(s)=0,\ s\notin\{-2,-4,\dots\},\ s\neq1\ \Longrightarrow\ \operatorname{Re}s=\tfrac12; \\
L(s,\chi)=0,\ 0<\operatorname{Re}s<1\ \Longrightarrow\ \operatorname{Re}s=\tfrac12 .
\end{gathered}}
\end{bpstatement}
The first export has exactly Mathlib's type \href{https://github.com/leanprover-community/mathlib4/blob/0df444a360eaa60ab8c11dca51a86af692955474/Mathlib/NumberTheory/LSeries/RiemannZeta.lean\#L182-L186}{\bpcode{RiemannHypothesis}}. The other four state the zeta zeros in expanded form, nonvanishing of $L(s,\chi)$ for $\operatorname{Re}s>\tfrac12$ away from the principal pole, and the critical-line conclusion in the open strip for all characters and for primitive ones.

All five pass through one assertion: the half-plane theorem for every Dirichlet character. The extension to further $L$-functions mentioned in the announcement has no endpoint in the archive.

\bpfigure{strip}{0.62}{Critical strip, reflection and the half-plane theorem}{Schematic plane with the computed first zeta zeros on the critical line; a hypothetical off-line zero \(ρ\) and its mirror image 1 − conj(\(ρ\)) under the functional equation.}
\begin{bpleantable}{dest}
\bpleanrow{https://github.com/Deicyde/riemann-hypothesis/blob/40e56bab026552e089d7afcfeff08d2defb9c005/src/Main.lean\#L4-L4}{\bpcode{riemannHypothesis}}{Main.lean, L4}{type: \bpcode{RiemannHypothesis} (Mathlib)}
\bpleanrow{https://github.com/Deicyde/riemann-hypothesis/blob/40e56bab026552e089d7afcfeff08d2defb9c005/src/Main.lean\#L6-L6}{\bpcode{riemannZeta\_\allowbreak{}eq\_\allowbreak{}zero\_\allowbreak{}re}}{Main.lean, L6}{the same statement, quantifiers written out}
\bpleanrow{https://github.com/Deicyde/riemann-hypothesis/blob/40e56bab026552e089d7afcfeff08d2defb9c005/src/Main.lean\#L8-L8}{\bpcode{LFunction\_\allowbreak{}ne\_\allowbreak{}zero\_\allowbreak{}of\_\allowbreak{}half\_\allowbreak{}lt\_\allowbreak{}re}}{Main.lean, L8}{half-plane nonvanishing, all characters}
\bpleanrow{https://github.com/Deicyde/riemann-hypothesis/blob/40e56bab026552e089d7afcfeff08d2defb9c005/src/Main.lean\#L10-L10}{\bpcode{LFunction\_\allowbreak{}eq\_\allowbreak{}zero\_\allowbreak{}re}}{Main.lean, L10}{critical-strip zeros, all characters}
\bpleanrow{https://github.com/Deicyde/riemann-hypothesis/blob/40e56bab026552e089d7afcfeff08d2defb9c005/src/Main.lean\#L12-L12}{\bpcode{LFunction\_\allowbreak{}eq\_\allowbreak{}zero\_\allowbreak{}re\_\allowbreak{}of\_\allowbreak{}isPrimitive}}{Main.lean, L12}{critical-strip zeros, primitive characters}
\bpleanrow{https://github.com/Deicyde/riemann-hypothesis/blob/40e56bab026552e089d7afcfeff08d2defb9c005/src/Statements.lean\#L3-L16}{Statements (target types)}{Statements.lean, L3–16}{}
\bpleanrow{https://github.com/leanprover-community/mathlib4/blob/0df444a360eaa60ab8c11dca51a86af692955474/Mathlib/NumberTheory/LSeries/RiemannZeta.lean\#L182-L186}{Mathlib \bpcode{RiemannHypothesis}}{RiemannZeta.lean, L182–186 (Mathlib)}{Mathlib 0df444a3}
\end{bpleantable}
\bpinterfaces{\bpiface{reflection}{critical line}{\href{https://github.com/Deicyde/riemann-hypothesis/blob/40e56bab026552e089d7afcfeff08d2defb9c005/src/Development/AutomorphicArithmetic/CharacterBoxes/ClaimStrip.lean\#L75-L100}{\bpcode{Claim.\allowbreak{}zeta}} \textperiodcentered{} \href{https://github.com/Deicyde/riemann-hypothesis/blob/40e56bab026552e089d7afcfeff08d2defb9c005/src/Development/AutomorphicArithmetic/CharacterBoxes/ClaimStrip.lean\#L59-L73}{\bpcode{Claim.\allowbreak{}strip}}}}{}
\bpownednodes{13}{5\,\bptiericonplain{two-kernels}, 8\,\bptiericonplain{read}}{\bpnodeitem{main-rh}{two-kernels}{Mathlib RH}\bpnodeitem{main-zeta}{two-kernels}{Zeta in expanded form}\bpnodeitem{main-right}{two-kernels}{Nonvanishing right of one half}\bpnodeitem{main-strip}{two-kernels}{Dirichlet GRH in the open strip}\bpnodeitem{main-prim}{two-kernels}{Primitive critical-strip zeros}\bpnodeitem{stmt-zeta}{read}{Zeta target type}\bpnodeitem{stmt-right}{read}{Half-plane target type}\bpnodeitem{stmt-strip}{read}{Dirichlet strip target type}\bpnodeitem{stmt-prim}{read}{Primitive strip target type}\bpnodeitem{rh-all}{read}{The two zeta exports share a proof}\bpnodeitem{all-chi}{read}{All-character half-plane theorem}\bpnodeitem{crit-all}{read}{Primitive critical-line wrapper}\bpnodeitem{claim-all}{read}{Three endpoint propositions}}
\bpstatus{two-kernels}{Two independent kernels accepted the complete proof, and all five theorems matched the canonical Mathlib challenge. Andrew Sutherland and Simon \bpcode{DeDeo} each report independently building and replaying all 175,187 original modules with Lean 4.33.1. \bpcode{DeDeo} reports zero errors and crashes. Sutherland additionally records a con-leche comparison of the five targets and a lean4lean check of every original module. Five modules required higher reduction limits. The individual logs identify the names, flags and successful reruns. (Status recorded by the reading project at 2026-11-20 23:51 UTC.)}
\bpcheck{What machine checking cannot settle: that the five Lean statements say what is meant (they are compared with an independently written Mathlib challenge), the soundness of the checkers and of Mathlib's definitions, and the mathematical community's own reading of the argument.}{\bpaudit{main-rh}{Exact endpoint replay}{Compile the five unmodified exports at the archive pins, verify their exact types, inspect their transitive axioms, and replay the selected kernel checker.}}
\bpannouncement{\bptlitem{20}{11/16/26}{Dirichlet GRH}{Endpoint statement matched}}

\bpcomponent{reflection}{dest}{2}{Reflection}{checked}
\bpkicker{\bpkickertext{functional equations}}
\bpgist{A zero left of the critical line reflects to a zero right of it, so the right half-plane is all one has to clear.}
\begin{bpstatement}{dest}
\bpdisplay{\begin{gathered}
\zeta(s)=2(2\pi)^{-w}\,\Gamma(w)\cos\!\big(\tfrac{\pi w}{2}\big)\,\zeta(w),\quad w=1-s; \\
\Lambda(s,\chi)=\varepsilon(\chi)\,\Lambda(1-s,\chi^{-1})\ \ (\chi\ \text{primitive}).
\end{gathered}}
\end{bpstatement}
For zeta, put $w=1-s$ with $\operatorname{Re}w>\tfrac12$. The power and Gamma factors do not vanish, and a zero of the cosine would make $s$ an even integer, i.e. $0$ or a trivial zero. Hence a nontrivial zero on the left gives a zero on the right.

For a primitive character the completed functional equation sends a strip zero at $s$ to a zero of $L(\cdot,\chi^{-1})$ at $1-s$; the Gamma factors are regular and nonzero for $0<\operatorname{Re}s<1$. Taking $\chi^{-1}$ matters when $\chi$ is complex.

\bpfigref{strip}{Critical strip, reflection and the half-plane theorem}
\begin{bpleantable}{dest}
\bpleanrow{https://github.com/Deicyde/riemann-hypothesis/blob/40e56bab026552e089d7afcfeff08d2defb9c005/src/Formalism/GammaFactorNonvanishing/RhOfRight.lean\#L5-L62}{\bpcode{Formalism.\allowbreak{}VW.\allowbreak{}rh\_\allowbreak{}of\_\allowbreak{}right}}{RhOfRight.lean, L5–62}{reflect a zeta zero}
\bpleanrow{https://github.com/Deicyde/riemann-hypothesis/blob/40e56bab026552e089d7afcfeff08d2defb9c005/src/Development/AutomorphicArithmetic/CharacterBoxes/ClaimStrip.lean\#L42-L57}{\bpcode{Claim.\allowbreak{}prim\_\allowbreak{}reflect}}{ClaimStrip.lean, L42–57}{transfer a primitive zero to 1 − s}
\bpleanrow{https://github.com/Deicyde/riemann-hypothesis/blob/40e56bab026552e089d7afcfeff08d2defb9c005/src/Development/AutomorphicArithmetic/CharacterBoxes/ClaimStrip.lean\#L59-L73}{\bpcode{Claim.\allowbreak{}strip}}{ClaimStrip.lean, L59–73}{half-plane ⇒ critical-strip zeros on the line}
\bpleanrow{https://github.com/Deicyde/riemann-hypothesis/blob/40e56bab026552e089d7afcfeff08d2defb9c005/src/Development/AutomorphicArithmetic/CharacterBoxes/ClaimStrip.lean\#L75-L100}{\bpcode{Claim.\allowbreak{}zeta}}{ClaimStrip.lean, L75–100}{half-plane ⇒ the zeta statement}
\bpleanrow{https://github.com/Deicyde/riemann-hypothesis/blob/40e56bab026552e089d7afcfeff08d2defb9c005/src/Formalism/GammaFactorNonvanishing/CritLine.lean\#L6-L9}{\bpcode{Formalism.\allowbreak{}VW.\allowbreak{}gammaFactor\_\allowbreak{}ne\_\allowbreak{}zero}}{CritLine.lean, L6–9}{Gamma factors do not vanish in the strip}
\end{bpleantable}
\bpinterfaces{\bpiface{descent}{every \(χ\), Re s > ½}{\href{https://github.com/Deicyde/riemann-hypothesis/blob/40e56bab026552e089d7afcfeff08d2defb9c005/src/Development/AutomorphicArithmetic/CharacterBoxes/Imprimitive.lean\#L5-L35}{\bpcode{UpperTheory.\allowbreak{}Top.\allowbreak{}imprimitive}}}}{\bpiface{endpoints}{critical line}{\href{https://github.com/Deicyde/riemann-hypothesis/blob/40e56bab026552e089d7afcfeff08d2defb9c005/src/Development/AutomorphicArithmetic/CharacterBoxes/ClaimStrip.lean\#L75-L100}{\bpcode{Claim.\allowbreak{}zeta}} \textperiodcentered{} \href{https://github.com/Deicyde/riemann-hypothesis/blob/40e56bab026552e089d7afcfeff08d2defb9c005/src/Development/AutomorphicArithmetic/CharacterBoxes/ClaimStrip.lean\#L59-L73}{\bpcode{Claim.\allowbreak{}strip}}}}
\bpownednodes{8}{3\,\bptiericonplain{checked}, 1\,\bptiericonplain{covered}, 4\,\bptiericonplain{read}}{\bpnodeitem{rh-every}{read}{Specialize at modulus one}\bpnodeitem{rh-right}{checked}{Reflect a zeta zero}\bpnodeitem{vw-crit}{read}{Reflect a primitive character}\bpnodeitem{gamma}{read}{Gamma factors in the strip}\bpnodeitem{strip}{checked}{Reflect an arbitrary character}\bpnodeitem{prim-reflect}{covered}{Transfer a primitive zero to 1−s}\bpnodeitem{bundle-zeta}{checked}{Zeta component of the conjunction}\bpnodeitem{alternate-rh}{read}{A second zeta deduction}}
\bpstatus{checked}{\href{https://github.com/Deicyde/riemann-hypothesis/blob/40e56bab026552e089d7afcfeff08d2defb9c005/src/Development/AutomorphicArithmetic/CharacterBoxes/ClaimStrip.lean\#L59-L73}{\bpcode{Claim.\allowbreak{}strip}}, \href{https://github.com/Deicyde/riemann-hypothesis/blob/40e56bab026552e089d7afcfeff08d2defb9c005/src/Development/AutomorphicArithmetic/CharacterBoxes/ClaimStrip.lean\#L75-L100}{\bpcode{Claim.\allowbreak{}zeta}} and \href{https://github.com/Deicyde/riemann-hypothesis/blob/40e56bab026552e089d7afcfeff08d2defb9c005/src/Formalism/GammaFactorNonvanishing/RhOfRight.lean\#L5-L62}{\bpcode{Formalism.\allowbreak{}VW.\allowbreak{}rh\_\allowbreak{}of\_\allowbreak{}right}} passed type, axiom and kernel checks at Lean 4.33.1 (Helios job 26639692).}
\bpcheck{None for the deduction itself; it is classical and machine-checked.}{}

\bpcomponent{descent}{dest}{3}{Conductor descent}{checked}
\bpkicker{\bpkickertext{Euler factors}}
\bpgist{An imprimitive character differs from its primitive inducing character by finitely many Euler factors, none of which vanishes for $\operatorname{Re}s>0$.}
\begin{bpstatement}{dest}
\bpdisplay{\begin{gathered}
L(s,\chi)=L(s,\chi^{*})\prod_{p\mid N}\bigl(1-\chi^{*}(p)p^{-s}\bigr), \\
\bigl|\chi^{*}(p)p^{-s}\bigr|\le p^{-\sigma}<1\quad(\sigma>0).
\end{gathered}}
\end{bpstatement}
Primes dividing the conductor have $\chi^*(p)=0$; for the others the second term of the factor has modulus less than one. So nonvanishing passes from primitive characters to all characters, with the single exception of the principal character's pole at $s=1$.

The strict left boundary of the strip matters: $1-2^{-s}$ vanishes on $\operatorname{Re}s=0$, so $L(s,\mathbf 1_2)=(1-2^{-s})\zeta(s)$ has extra zeros on the imaginary axis.

\begin{bpleantable}{dest}
\bpleanrow{https://github.com/Deicyde/riemann-hypothesis/blob/40e56bab026552e089d7afcfeff08d2defb9c005/src/Development/AutomorphicArithmetic/CharacterBoxes/Imprimitive.lean\#L5-L35}{\bpcode{UpperTheory.\allowbreak{}Top.\allowbreak{}imprimitive}}{Imprimitive.lean, L5–35}{primitive ⇒ all characters}
\bpleanrow{https://github.com/Deicyde/riemann-hypothesis/blob/40e56bab026552e089d7afcfeff08d2defb9c005/src/Development/AutomorphicArithmetic/CharacterBoxes/ClaimStrip.lean\#L23-L33}{\bpcode{Claim.\allowbreak{}prim\_\allowbreak{}zero}}{ClaimStrip.lean, L23–33}{remove the imprimitive Euler factors}
\bpleanrow{https://github.com/Deicyde/riemann-hypothesis/blob/40e56bab026552e089d7afcfeff08d2defb9c005/src/Development/AutomorphicArithmetic/CharacterBoxes/ClaimStrip.lean\#L7-L21}{\bpcode{Claim.\allowbreak{}factor\_\allowbreak{}ne\_\allowbreak{}zero}}{ClaimStrip.lean, L7–21}{a removed Euler factor is nonzero}
\end{bpleantable}
\bpinterfaces{\bpiface{choreography}{\bpcode{ZeroFreeAll} (primitive)}{\href{https://github.com/Deicyde/riemann-hypothesis/blob/40e56bab026552e089d7afcfeff08d2defb9c005/src/Development/AutomorphicArithmetic/CharacterBoxes/ZeroFreeAllOfTopG.lean\#L26-L34}{\bpcode{zeroFreeAll\_\allowbreak{}of\_\allowbreak{}topG}}}}{\bpiface{reflection}{every \(χ\), Re s > ½}{\href{https://github.com/Deicyde/riemann-hypothesis/blob/40e56bab026552e089d7afcfeff08d2defb9c005/src/Development/AutomorphicArithmetic/CharacterBoxes/Imprimitive.lean\#L5-L35}{\bpcode{UpperTheory.\allowbreak{}Top.\allowbreak{}imprimitive}}}}
\bpownednodes{4}{1\,\bptiericonplain{checked}, 2\,\bptiericonplain{covered}, 1\,\bptiericonplain{read}}{\bpnodeitem{prim-zero}{covered}{Remove imprimitive Euler factors}\bpnodeitem{factor}{covered}{Nonzero removed Euler factor}\bpnodeitem{every-green}{read}{Pass from primitive to all characters}\bpnodeitem{imprimitive}{checked}{Restore all characters}}
\bpstatus{checked}{\href{https://github.com/Deicyde/riemann-hypothesis/blob/40e56bab026552e089d7afcfeff08d2defb9c005/src/Development/AutomorphicArithmetic/CharacterBoxes/Imprimitive.lean\#L5-L35}{\bpcode{UpperTheory.\allowbreak{}Top.\allowbreak{}imprimitive}} passed its type, axiom and defining-module checks in the 1,421-module Dirichlet closure (Helios job 26647685).}
\bpcheck{None for the deduction; classical and machine-checked.}{}

\bpstratum{ana}{The analytic engine}{A localized Landau oscillation argument converts one-sided bounds on whole power windows into a zero-free half-plane.}{\bpstratcomp{statistic}\bpstratcomp{residues}\bpstratcomp{criterion}\bpstratcomp{choreography}}{\bpstratflow{comparison}{choreography}{\(G_{P}\) ≤ \(x^{1/2+9ε}\) on windows}\bpstratflow{choreography}{descent}{\bpcode{ZeroFreeAll} (primitive)}}
\bpcomponent{statistic}{ana}{4}{The centered prime sum}{checked}
\bpkicker{\bpkickertext{what must be bounded}}
\bpgist{A ramp-weighted von Mangoldt sum, centered so that the main terms cancel exactly, and its positive combinations over a cyclic family of characters.}
\begin{bpstatement}{ana}
\bpdisplay{\begin{gathered}
G(y)=\sum_{n\le y}\Lambda(n)\Bigl(\frac yn-1\Bigr)-y\log y+(1+\gamma)\,y, \\
G_P(x)=G(x)+\frac{1}{3o+1}\sum_{\chi\in P\setminus\{1\}}\operatorname{Re} g_\chi(x).
\end{gathered}}
\end{bpstatement}
Here $g_\chi(y)=\sum_{n\le y}\Lambda(n)\chi(n)(y/n-1)+y\,(L'/L)(1,\chi)$ for nonprincipal primitive $\chi$, and $P$ consists of the primitive characters inducing the cyclic group generated by a fixed primitive $\psi$ (order $o$). At the trivial level $G_P=G$.

By the explicit formula $G(y)=\log 2\pi-\sum_\rho y^\rho/(\rho(\rho-1))+O(y^{-2})$, so under RH $|G(y)-\log 2\pi|\le(2+\gamma-\log 4\pi)\sqrt y$. The proof needs much less: an \emph{upper} bound $G_P(x)\le x^{1/2+O(\varepsilon)}$ on whole windows.

\bpfigure{explicit-formula}{0.86}{Prime sum versus the sum over zeros}{Computed: normalized G(y) from primes (amber) and the explicit formula with the first 100 zeta zeros (indigo), 100 ≤ y ≤ 2e6; dashed: the term a hypothetical zero at 0.62+40i would add (it grows like \(y^{0}\).12).}
\begin{bpleantable}{ana}
\bpleanrow{https://github.com/Deicyde/riemann-hypothesis/blob/40e56bab026552e089d7afcfeff08d2defb9c005/src/Formalism/StatementOC.lean\#L13-L13}{\bpcode{Formalism.\allowbreak{}Gz}}{StatementOC.lean, L13}{}
\bpleanrow{https://github.com/Deicyde/riemann-hypothesis/blob/40e56bab026552e089d7afcfeff08d2defb9c005/src/Development/AutomorphicArithmetic/DedekindLFunctions/CharFamily.lean\#L10-L45}{\bpcode{Deduction.\allowbreak{}G} / Gchi (\bpcode{CharFamily})}{CharFamily.lean, L10–45}{}
\bpleanrow{https://github.com/Deicyde/riemann-hypothesis/blob/40e56bab026552e089d7afcfeff08d2defb9c005/src/Formalism/SharpRestriction/RamificationFrobenius/TotientGaloisClasses/ExplicitFormulaErrorOfHostConditions.lean\#L38-L38}{\bpcode{Formalism.\allowbreak{}HC.\allowbreak{}Level.\allowbreak{}Gm}}{ExplicitFormulaErrorOfHostConditions.lean, L38}{}
\bpleanrow{https://github.com/Deicyde/riemann-hypothesis/blob/40e56bab026552e089d7afcfeff08d2defb9c005/src/Development/AutomorphicArithmetic/PrimitiveCosets/GzEq.lean\#L6-L47}{\bpcode{UpperTheory.\allowbreak{}Comparison₂.\allowbreak{}Gz\_\allowbreak{}eq}}{GzEq.lean, L6–47}{}
\end{bpleantable}
\bpinterfaces{}{\bpiface{residues}{Mellin transform}{\href{https://github.com/Deicyde/riemann-hypothesis/blob/40e56bab026552e089d7afcfeff08d2defb9c005/src/Development/AutomorphicArithmetic/DedekindLFunctions/LandauTheorems/MellinRamp.lean\#L10-L32}{\bpcode{MellinRamp}}}}
\bpownednodes{8}{1\,\bptiericonplain{checked}, 7\,\bptiericonplain{covered}}{\bpnodeitem{g-m}{covered}{The weighted prime sum}\bpnodeitem{window-def}{covered}{Uniform families on power windows}\bpnodeitem{moment-target}{covered}{The first moment of the Chebyshev error}\bpnodeitem{moment-theorem}{checked}{Evaluate the first moment}\bpnodeitem{moment-calculation}{covered}{Match the two limits at one}\bpnodeitem{moment-laurent}{covered}{Read Euler’s constant from the logarithmic derivative}\bpnodeitem{moment-majorant}{covered}{Dominate the centered character sum}\bpnodeitem{moment-dominated}{covered}{Pass the Mellin integral to \(σ\) = 1}}
\bpstatus{checked}{\href{https://github.com/Deicyde/riemann-hypothesis/blob/40e56bab026552e089d7afcfeff08d2defb9c005/src/Development/AutomorphicArithmetic/PrimitiveCosets/GzEq.lean\#L6-L47}{\bpcode{Gz\_\allowbreak{}eq}} (\href{https://github.com/Deicyde/riemann-hypothesis/blob/40e56bab026552e089d7afcfeff08d2defb9c005/src/Formalism/StatementOC.lean\#L13-L13}{\bpcode{Formalism.\allowbreak{}Gz}} = \bpcode{An.\allowbreak{}G} = \href{https://github.com/Deicyde/riemann-hypothesis/blob/40e56bab026552e089d7afcfeff08d2defb9c005/src/Development/AutomorphicArithmetic/DedekindLFunctions/CharFamily.lean\#L10-L45}{\bpcode{Deduction.\allowbreak{}G}} pointwise) passed its checks in a 130-module union; the definitions are inside every checked closure below.}
\bpcheck{None: a definition and a checked normalization identity.}{}
\bpannouncement{\bptlitem{5}{10/20/26}{The smoothed von Mangoldt sum}{Formula matched}}

\bpcomponent{residues}{ana}{5}{Mellin transform and positive residues}{checked}
\bpkicker{\href{https://github.com/Deicyde/riemann-hypothesis/blob/40e56bab026552e089d7afcfeff08d2defb9c005/src/Development/AutomorphicArithmetic/DedekindLFunctions/LandauTheorems/poleData.lean\#L248-L259}{\bpcode{poleData}}}
\bpgist{The Mellin transform of a ramp is $n^{-s}/(s(s-1))$; a zero of any family member becomes a pole whose residue cannot be cancelled, because the weights are positive.}
\begin{bpstatement}{ana}
\bpdisplay{\begin{gathered}
\int_1^\infty\Bigl(\frac yn-1\Bigr)_{\!+}y^{-s-1}\,dy=\frac{n^{-s}}{s(s-1)}, \\
\operatorname*{Res}_{s=\rho}=-\frac{S(\rho)}{\rho(\rho-1)},\quad S(\rho)=\tfrac12\sum_i a_i\bigl(m_{\chi_i}(\rho)+m_{\chi_i^{-1}}(\rho)\bigr)>0 .
\end{gathered}}
\end{bpstatement}
Summing ramps against $\Lambda(n)\chi(n)$ produces $-L'/L$; the compensating linear term removes the pole at $s=1$, so $\int_1^\infty G_\chi(y)y^{-s-1}dy=-[(L'/L)(s,\chi)-s(L'/L)(1,\chi)]/(s(s-1))$.

Taking real parts averages $\chi$ and $\chi^{-1}$, so every member vanishing at $\rho$ contributes a positive multiple of the same complex number. For a real zero $0<\beta<1$ the residue is even positive, which lets the one-sided bound exclude exceptional real zeros too.

\bpfigure{residue}{1.0}{Positive weights keep the residue alive}{Schematic, exact in direction: at a common zero \(ρ\) every member of a positive family contributes a positive multiple of the same complex number -1/(\(ρ\)(\(ρ\)-1)), so the contributions add; with signed weights two contributions could cancel.}
\begin{bpleantable}{ana}
\bpleanrow{https://github.com/Deicyde/riemann-hypothesis/blob/40e56bab026552e089d7afcfeff08d2defb9c005/src/Development/AutomorphicArithmetic/DedekindLFunctions/LandauTheorems/MellinRamp.lean\#L10-L32}{\bpcode{MellinRamp}}{MellinRamp.lean, L10–32}{}
\bpleanrow{https://github.com/Deicyde/riemann-hypothesis/blob/40e56bab026552e089d7afcfeff08d2defb9c005/src/Development/AutomorphicArithmetic/DedekindLFunctions/LandauTheorems/mellinGchiOfNeOne.lean\#L116-L125}{\bpcode{mellinGchiOfNeOne}}{mellinGchiOfNeOne.lean, L116–125}{}
\bpleanrow{https://github.com/Deicyde/riemann-hypothesis/blob/40e56bab026552e089d7afcfeff08d2defb9c005/src/Development/AutomorphicArithmetic/DedekindLFunctions/LandauTheorems/poleData.lean\#L248-L259}{\bpcode{Deduction.\allowbreak{}poleData}}{poleData.lean, L248–259}{}
\bpleanrow{https://github.com/Deicyde/riemann-hypothesis/blob/40e56bab026552e089d7afcfeff08d2defb9c005/src/Development/AutomorphicArithmetic/DedekindLFunctions/LandauTheorems/poleData.lean\#L336-L383}{positivity of S(\(ρ\))}{poleData.lean, L336–383}{}
\end{bpleantable}
\bpinterfaces{\bpiface{statistic}{Mellin transform}{\href{https://github.com/Deicyde/riemann-hypothesis/blob/40e56bab026552e089d7afcfeff08d2defb9c005/src/Development/AutomorphicArithmetic/DedekindLFunctions/LandauTheorems/MellinRamp.lean\#L10-L32}{\bpcode{MellinRamp}}}}{\bpiface{criterion}{pole data, S(\(ρ\)) > 0}{\href{https://github.com/Deicyde/riemann-hypothesis/blob/40e56bab026552e089d7afcfeff08d2defb9c005/src/Development/AutomorphicArithmetic/DedekindLFunctions/LandauTheorems/poleData.lean\#L248-L259}{\bpcode{Deduction.\allowbreak{}poleData}}}}
\bpownednodes{1}{1\,\bptiericonplain{checked}}{\bpnodeitem{pole}{checked}{Poles of the finite-family transform}}
\bpstatus{checked}{\href{https://github.com/Deicyde/riemann-hypothesis/blob/40e56bab026552e089d7afcfeff08d2defb9c005/src/Development/AutomorphicArithmetic/DedekindLFunctions/LandauTheorems/poleData.lean\#L248-L259}{\bpcode{Deduction.\allowbreak{}poleData}} passed type, axiom and kernel checks (38-module landmark closure, Helios job 26639692).}
\bpcheck{None: an analytic lemma, machine-checked.}{}

\bpcomponent{criterion}{ana}{6}{The window criterion}{two-kernels}
\bpkicker{\bpkickertext{localized Landau}}
\bpgist{If a positive family's prime sum stays below $Cy^\theta$ throughout arbitrarily late windows $[Y,Y^{1+\delta}]$, none of its $L$-functions vanishes to the right of $\theta$.}
\begin{bpstatement}{ana}
\bpdisplay{\exists C\ \forall Y_0\ \exists Y\ge Y_0\ \forall y\in[Y,Y^{1+\delta}]:\ G_F(y)\le C\,y^{\theta}\ \Longrightarrow\ L(s,\chi_i)\neq0\ \text{ for }\operatorname{Re}s>\theta .}
\end{bpstatement}
With $u=\log y$ and $g(u)=G_F(e^u)$, the moments $A_k(z)=\int g(u)\,\frac{u^k e^{-\sigma u}}{k!}e^{-itu}\,du$ are $(-1)^k\Phi^{(k)}(z)/k!$ for the Laplace transform $\Phi$. The kernel peaks at $u=k/\sigma$, so $k=\lceil(1+\delta/2)\sigma U\rceil$ aims it at the window $[U,(1+\delta)U]$; the tails are exponentially small.

The one-sided trick: $|g|+g=2\max(g,0)\le 2R$ for any majorant $R\ge g$, which controls $|A_k(\sigma+it)|+\operatorname{Re}A_k(\sigma)$. A pole at distance $d$ contributes about $d^{-k}$ while the window bound gives $e^{\eta U}r^{-k}$ with $r>d$; since $k\ge U$ and $\eta$ is arbitrary, the pole cannot exist.

\bpfigure{spotlight}{0.86}{The spotlight of high derivatives}{Computed: kernels \(u^{k}\) \(e^{-u}\)/k! (normalized to peak 1) with k = ceil(1.25 U) for U = 20, 40, 80, 160 and the windows [U, 1.5 U] in u = log y; labels give the fraction of each kernel's mass inside its window.}
\begin{bpleantable}{ana}
\bpleanrow{https://github.com/Deicyde/riemann-hypothesis/blob/40e56bab026552e089d7afcfeff08d2defb9c005/src/Development/AutomorphicArithmetic/DedekindLFunctions/LandauTheorems/ZeroFreeOfRpowWindow.lean\#L11-L14}{\bpcode{Deduction.\allowbreak{}zeroFree\_\allowbreak{}of\_\allowbreak{}rpow\_\allowbreak{}window}}{ZeroFreeOfRpowWindow.lean, L11–14}{the criterion (statement)}
\bpleanrow{https://github.com/Deicyde/riemann-hypothesis/blob/40e56bab026552e089d7afcfeff08d2defb9c005/src/Development/AutomorphicArithmetic/DedekindLFunctions/LandauTheorems/landauWindow.lean\#L83-L95}{\bpcode{Deduction.\allowbreak{}landau\_\allowbreak{}window}}{landauWindow.lean, L83–95}{rule out a pole right of \(θ\)}
\bpleanrow{https://github.com/Deicyde/riemann-hypothesis/blob/40e56bab026552e089d7afcfeff08d2defb9c005/src/Development/AutomorphicArithmetic/DedekindLFunctions/LandauTheorems/landauWindow.lean\#L83-L116}{isolation of an offending pole (proof)}{landauWindow.lean, L83–116}{}
\bpleanrow{https://github.com/Deicyde/riemann-hypothesis/blob/40e56bab026552e089d7afcfeff08d2defb9c005/src/Development/AutomorphicArithmetic/DedekindLFunctions/LandauTheorems/landauWindowIneq.lean\#L57-L99}{one-sided moment inequality}{landauWindowIneq.lean, L57–99}{}
\bpleanrow{https://github.com/Deicyde/riemann-hypothesis/blob/40e56bab026552e089d7afcfeff08d2defb9c005/src/Development/AutomorphicArithmetic/DedekindLFunctions/LandauTheorems/landauTails.lean\#L80-L102}{tail estimate}{landauTails.lean, L80–102}{}
\end{bpleantable}
\bpinterfaces{\bpiface{residues}{pole data, S(\(ρ\)) > 0}{\href{https://github.com/Deicyde/riemann-hypothesis/blob/40e56bab026552e089d7afcfeff08d2defb9c005/src/Development/AutomorphicArithmetic/DedekindLFunctions/LandauTheorems/poleData.lean\#L248-L259}{\bpcode{Deduction.\allowbreak{}poleData}}}}{\bpiface{choreography}{criterion at each level}{\href{https://github.com/Deicyde/riemann-hypothesis/blob/40e56bab026552e089d7afcfeff08d2defb9c005/src/Development/AutomorphicArithmetic/DedekindLFunctions/LandauTheorems/ZeroFreeOfRpowWindow.lean\#L11-L14}{\bpcode{zeroFree\_\allowbreak{}of\_\allowbreak{}rpow\_\allowbreak{}window}}}}
\bpownednodes{4}{1\,\bptiericonplain{checked}, 3\,\bptiericonplain{covered}}{\bpnodeitem{form-b}{covered}{The analytic criterion for a level}\bpnodeitem{landau}{checked}{The character-family power-window theorem}\bpnodeitem{landau-lemma}{covered}{Apply the pole obstruction}\bpnodeitem{oscillation}{covered}{Rule out a pole right of \(θ\)}}
\bpstatus{two-kernels}{\href{https://github.com/Deicyde/riemann-hypothesis/blob/40e56bab026552e089d7afcfeff08d2defb9c005/src/Development/AutomorphicArithmetic/DedekindLFunctions/LandauTheorems/ZeroFreeOfRpowWindow.lean\#L11-L14}{\bpcode{zeroFree\_\allowbreak{}of\_\allowbreak{}rpow\_\allowbreak{}window}} passed at the archive pin; its specialization to zeta (\bpcode{riemannHypothesis\_\allowbreak{}of\_\allowbreak{}prime\_\allowbreak{}windows}) passed a fresh Lean replay over 10,530 modules and the Rust kernel NanoDA on 61,697 exported declarations.}
\bpcheck{None: machine-checked twice.}{}
\bpannouncement{\bptlitem{6}{10/8/26}{Two-support positivity test}{Historical calculation}}

\bpcomponent{choreography}{ana}{7}{The order of choices}{checked}
\bpkicker{\bpkickertext{\(θ\) after the zero, \(ε\) after \(θ\)}}
\bpgist{The loss exponents are fixed first; given a zero with $\beta>\tfrac12$ choose $\theta\in(\tfrac12,\beta)$, then $\varepsilon$ with $\tfrac12+9\varepsilon\le\theta$. No bound at $\varepsilon=0$ is ever needed.}
\begin{bpstatement}{ana}
\bpdisplay{\begin{gathered}
\theta=\tfrac12+\tfrac12\min\bigl(\beta-\tfrac12,\tfrac12\bigr), \\
\tfrac12+(e_D+e_L+7)\,\varepsilon\le\theta\quad(e_D=2,\ e_L=0).
\end{gathered}}
\end{bpstatement}
This is where the arithmetic and the analysis meet. The arithmetic must deliver, for every small $\varepsilon$, families with $G_P(x)\le x^{1/2+(e_D+e_L+7)\varepsilon}$ on arbitrarily late windows $[Y,Y^{3/2}]$, with $e_D,e_L$ independent of $\varepsilon$. The families themselves may depend on $\varepsilon$.

Applying the criterion at each primitive character's cyclic level gives \bpcode{ZeroFreeAll}: $L(s,\psi)\neq0$ for $\operatorname{Re}s>\tfrac12$ and every primitive $\psi$.

\bpfigure{choreography}{1.0}{The order of choices (schematic)}{Schematic: the loss exponents are fixed first; a candidate zero \(β\) > 1/2 is given; \(θ\) is chosen between 1/2 and \(β\); \(ε\) is chosen so that 1/2 + (\(e_{D}\)+\(e_{L}\)+7) \(ε\) ≤ \(θ\); the arithmetic then supplies windows for that \(ε\); the analytic criterion excludes the zero.}
\begin{bpleantable}{ana}
\bpleanrow{https://github.com/Deicyde/riemann-hypothesis/blob/40e56bab026552e089d7afcfeff08d2defb9c005/src/Formalism/SharpRestriction/RamificationFrobenius/TotientGaloisClasses/TopWin.lean\#L68-L81}{\bpcode{Formalism.\allowbreak{}HC.\allowbreak{}level\_\allowbreak{}of\_\allowbreak{}tiedWin}}{TopWin.lean, L68–81}{choice of \(θ\) and \(ε\) after a candidate zero}
\bpleanrow{https://github.com/Deicyde/riemann-hypothesis/blob/40e56bab026552e089d7afcfeff08d2defb9c005/src/Development/AutomorphicArithmetic/CharacterBoxes/ZeroFreeAllOfTopG.lean\#L26-L34}{\bpcode{UpperTheory.\allowbreak{}Top.\allowbreak{}zeroFreeAll\_\allowbreak{}of\_\allowbreak{}topG}}{ZeroFreeAllOfTopG.lean, L26–34}{\bpcode{TopG} and hosts ⇒ primitive zero-free half-plane}
\bpleanrow{https://github.com/Deicyde/riemann-hypothesis/blob/40e56bab026552e089d7afcfeff08d2defb9c005/src/Development/AutomorphicArithmetic/ThetaCorrespondence/ThetaPairLifts/ThirdPlaceSeries/TopPlug.lean\#L12-L15}{\bpcode{UpperTheory.\allowbreak{}Top.\allowbreak{}zero\_\allowbreak{}free\_\allowbreak{}all\_\allowbreak{}of\_\allowbreak{}green\_\allowbreak{}v\_\allowbreak{}win}}{TopPlug.lean, L12–15}{selects the window width 3/2}
\end{bpleantable}
\bpinterfaces{\bpiface{comparison}{\(G_{P}\) ≤ \(x^{1/2+9ε}\) on windows}{\href{https://github.com/Deicyde/riemann-hypothesis/blob/40e56bab026552e089d7afcfeff08d2defb9c005/src/Development/AutomorphicArithmetic/PrimitiveCosets/LevelGmLeOnFamilyScales.lean\#L28-L30}{\bpcode{level\_\allowbreak{}gm\_\allowbreak{}le\_\allowbreak{}on\_\allowbreak{}family\_\allowbreak{}scales}}}\bpiface{criterion}{criterion at each level}{\href{https://github.com/Deicyde/riemann-hypothesis/blob/40e56bab026552e089d7afcfeff08d2defb9c005/src/Development/AutomorphicArithmetic/DedekindLFunctions/LandauTheorems/ZeroFreeOfRpowWindow.lean\#L11-L14}{\bpcode{zeroFree\_\allowbreak{}of\_\allowbreak{}rpow\_\allowbreak{}window}}}}{\bpiface{descent}{\bpcode{ZeroFreeAll} (primitive)}{\href{https://github.com/Deicyde/riemann-hypothesis/blob/40e56bab026552e089d7afcfeff08d2defb9c005/src/Development/AutomorphicArithmetic/CharacterBoxes/ZeroFreeAllOfTopG.lean\#L26-L34}{\bpcode{zeroFreeAll\_\allowbreak{}of\_\allowbreak{}topG}}}}
\bpownednodes{5}{1\,\bptiericonplain{checked}, 3\,\bptiericonplain{covered}, 1\,\bptiericonplain{read}}{\bpnodeitem{zero-green}{read}{Select the window width 3/2}\bpnodeitem{zero-top}{checked}{Produce primitive zero-free levels}\bpnodeitem{tied-cyc}{covered}{Assemble the cyclic levels}\bpnodeitem{tied}{covered}{Pass from levels to characters}\bpnodeitem{window-consumer}{covered}{Choose epsilon after a candidate zero}}
\bpstatus{checked}{\href{https://github.com/Deicyde/riemann-hypothesis/blob/40e56bab026552e089d7afcfeff08d2defb9c005/src/Development/AutomorphicArithmetic/CharacterBoxes/ZeroFreeAllOfTopG.lean\#L26-L34}{\bpcode{zeroFreeAll\_\allowbreak{}of\_\allowbreak{}topG}} passed its checks in the 1,421-module Dirichlet closure.}
\bpcheck{None for the logic; the premise \bpcode{TopG} is the arithmetic input.}{\bpaudit{window-consumer}{Separate the analytic implication}{Trace the finite positive-weight character family through \href{https://github.com/Deicyde/riemann-hypothesis/blob/40e56bab026552e089d7afcfeff08d2defb9c005/src/Development/AutomorphicArithmetic/DedekindLFunctions/LandauTheorems/poleData.lean\#L248-L259}{\bpcode{poleData}}, \href{https://github.com/Deicyde/riemann-hypothesis/blob/40e56bab026552e089d7afcfeff08d2defb9c005/src/Development/AutomorphicArithmetic/DedekindLFunctions/LandauTheorems/landauWindow.lean\#L83-L95}{\bpcode{landau\_\allowbreak{}window}} and the uniform window bound. Check the pole treatment and the order of quantifiers.}}

\bpstratum{est}{The estimate}{A finite character comparison turns geometric control of a chain of theta lifts of length about \(\sqrt{x}\) into an upper bound for a weighted prime sum.}{\bpstratcomp{comparison}\bpstratcomp{family}\bpstratcomp{hosts}}{\bpstratflow{inputs}{family}{definite-side data}\bpstratflow{cohomology}{hosts}{\bpcode{VWinAt} \(ψ\) 401}\bpstratflow{comparison}{choreography}{\(G_{P}\) ≤ \(x^{1/2+9ε}\) on windows}}
\bpcomponent{comparison}{est}{8}{The finite character comparison}{checked}
\bpkicker{\href{https://github.com/Deicyde/riemann-hypothesis/blob/40e56bab026552e089d7afcfeff08d2defb9c005/src/Development/AutomorphicArithmetic/PrimitiveCosets/CoreGm.lean\#L12-L74}{\bpcode{core\_\allowbreak{}gm}}\bpkickersep \href{https://github.com/Deicyde/riemann-hypothesis/blob/40e56bab026552e089d7afcfeff08d2defb9c005/src/Development/AutomorphicArithmetic/DedekindLFunctions/LandauTheorems/ChCore.lean\#L343-L351}{\bpcode{ch\_\allowbreak{}core}}\bpkickersep \href{https://github.com/Deicyde/riemann-hypothesis/blob/40e56bab026552e089d7afcfeff08d2defb9c005/src/Development/AutomorphicArithmetic/DedekindLFunctions/LandauTheorems/ChTop.lean\#L37-L72}{\bpcode{ch\_\allowbreak{}top}}}
\bpgist{Finite Fourier orthogonality on a subgroup and an exact main-term identity convert a geometric discrepancy bound into the prime-sum bound, with exponent $\tfrac12+(e_D+e_L+7)\varepsilon$.}
\begin{bpstatement}{est}
\bpdisplay{Q\,\tfrac{\delta}{2}\Bigl(G_\Lambda(x)-\Sigma(x)-4x^{1/2+\varepsilon}\Bigr)\le E+\tfrac{2Q}{x}\quad\Longrightarrow\quad G_\Lambda(x)\le x^{1/2+(e_D+e_L+7)\varepsilon}.}
\end{bpstatement}
The geometric discrepancy contains a cost proportional to $x/m$ and another proportional to $m$. With $m=x^a$ the larger exponent is $\max(1-a,a)$, smallest at $a=\tfrac12$: this is the square root. The worst of the four budget terms is $4x^{1/2+13\varepsilon/8}$.

Summing a character over the subgroup annihilates it unless it is trivial there; the survivors that are not retained ($\Sigma\le 4o\,x^{1/2+\varepsilon}$) need a separate square-root input (\bpcode{EAtH.\allowbreak{}bdG}). The losses add up to exactly $(e_D+e_L+7)\varepsilon$.

\bpfigure{balance}{0.86}{Where the square root and each \(ε\) are paid}{Left: exponents 1-a of x/m and a of m when m = \(x^{a}\); the larger is smallest at a = 1/2. Right: exact staircase of the losses in units of \(ε\), from the finite comparison to 1/2 + (\(e_{D}\)+\(e_{L}\)+7) \(ε\) = 1/2 + 9 \(ε\) for the source's choice \(e_{D}\) = 2, \(e_{L}\) = 0.}
\begin{bpleantable}{est}
\bpleanrow{https://github.com/Deicyde/riemann-hypothesis/blob/40e56bab026552e089d7afcfeff08d2defb9c005/src/Development/AutomorphicArithmetic/PrimitiveCosets/CoreGm.lean\#L12-L74}{\bpcode{UpperTheory.\allowbreak{}Comparison₂.\allowbreak{}core\_\allowbreak{}gm}}{CoreGm.lean, L12–74}{from the height discrepancy to the prime sum}
\bpleanrow{https://github.com/Deicyde/riemann-hypothesis/blob/40e56bab026552e089d7afcfeff08d2defb9c005/src/Development/AutomorphicArithmetic/DedekindLFunctions/LandauTheorems/ChCore.lean\#L343-L351}{\bpcode{Deduction.\allowbreak{}ch\_\allowbreak{}core}}{ChCore.lean, L343–351}{average over characters}
\bpleanrow{https://github.com/Deicyde/riemann-hypothesis/blob/40e56bab026552e089d7afcfeff08d2defb9c005/src/Development/AutomorphicArithmetic/DedekindLFunctions/LandauTheorems/ChCore.lean\#L16-L54}{subgroup identity (in \bpcode{ChCore.\allowbreak{}lean})}{ChCore.lean, L16–54}{}
\bpleanrow{https://github.com/Deicyde/riemann-hypothesis/blob/40e56bab026552e089d7afcfeff08d2defb9c005/src/Development/AutomorphicArithmetic/DedekindLFunctions/LandauTheorems/ChTop.lean\#L37-L72}{\bpcode{Deduction.\allowbreak{}ch\_\allowbreak{}top}}{ChTop.lean, L37–72}{bound the discrepancy per index}
\bpleanrow{https://github.com/Deicyde/riemann-hypothesis/blob/40e56bab026552e089d7afcfeff08d2defb9c005/src/Development/AutomorphicArithmetic/PrimitiveCosets/WindowInequalityOfArtinInputs.lean\#L17-L25}{\bpcode{UpperTheory.\allowbreak{}Comparison₂.\allowbreak{}windowInequality\_\allowbreak{}of\_\allowbreak{}artinInputs}}{WindowInequalityOfArtinInputs.lean, L17–25}{the estimate on a whole window}
\end{bpleantable}
\bpinterfaces{\bpiface{family}{\bpcode{TopG} (\(e_{D}\) = 2, \(e_{L}\) = 0)}{\bpcode{UpperTheory.\allowbreak{}Top.\allowbreak{}TopG}}\bpiface{hosts}{hosts, width 3/2}{\href{https://github.com/Deicyde/riemann-hypothesis/blob/40e56bab026552e089d7afcfeff08d2defb9c005/src/Development/AutomorphicArithmetic/CharacterBoxes/SupplyChi.lean\#L8-L16}{\bpcode{UpperTheory.\allowbreak{}Top.\allowbreak{}supplyChi}}}}{\bpiface{choreography}{\(G_{P}\) ≤ \(x^{1/2+9ε}\) on windows}{\href{https://github.com/Deicyde/riemann-hypothesis/blob/40e56bab026552e089d7afcfeff08d2defb9c005/src/Development/AutomorphicArithmetic/PrimitiveCosets/LevelGmLeOnFamilyScales.lean\#L28-L30}{\bpcode{level\_\allowbreak{}gm\_\allowbreak{}le\_\allowbreak{}on\_\allowbreak{}family\_\allowbreak{}scales}}}}
\bpownednodes{5}{4\,\bptiericonplain{checked}, 1\,\bptiericonplain{covered}}{\bpnodeitem{prime-bound}{covered}{The prime-sum inequality on tied families}\bpnodeitem{artin-window}{checked}{Estimate the weighted level on a window}\bpnodeitem{core-gm}{checked}{From the height discrepancy to the prime sum}\bpnodeitem{ch-top}{checked}{Bound the discrepancy per index}\bpnodeitem{ch-core}{checked}{Average over characters}}
\bpstatus{checked}{\href{https://github.com/Deicyde/riemann-hypothesis/blob/40e56bab026552e089d7afcfeff08d2defb9c005/src/Development/AutomorphicArithmetic/DedekindLFunctions/LandauTheorems/ChTop.lean\#L37-L72}{\bpcode{ch\_\allowbreak{}top}}, \href{https://github.com/Deicyde/riemann-hypothesis/blob/40e56bab026552e089d7afcfeff08d2defb9c005/src/Development/AutomorphicArithmetic/DedekindLFunctions/LandauTheorems/ChCore.lean\#L343-L351}{\bpcode{ch\_\allowbreak{}core}}, \href{https://github.com/Deicyde/riemann-hypothesis/blob/40e56bab026552e089d7afcfeff08d2defb9c005/src/Development/AutomorphicArithmetic/PrimitiveCosets/CoreGm.lean\#L12-L74}{\bpcode{core\_\allowbreak{}gm}} and \href{https://github.com/Deicyde/riemann-hypothesis/blob/40e56bab026552e089d7afcfeff08d2defb9c005/src/Development/AutomorphicArithmetic/PrimitiveCosets/WindowInequalityOfArtinInputs.lean\#L17-L25}{\bpcode{windowInequality\_\allowbreak{}of\_\allowbreak{}artinInputs}} passed at the archive pin (316 modules), with the comparison hypotheses \bpcode{EAtH} and \bpcode{CoreAtH} explicit.}
\bpcheck{The hypotheses: existence of the hosts, frames and witnesses at every scale of the window, and the square-root bound for discarded characters.}{}
\bpannouncement{\bptlitem{14}{11/15/26}{The provisional exponent two-thirds}{Intermediate strategy}\bptlitem{16}{10/14/26}{Period pairings and approximate orthogonality}{Supporting comparison}}

\bpcomponent{family}{est}{9}{The geometric family and the theta chain}{partial}
\bpkicker{\href{https://github.com/Deicyde/riemann-hypothesis/blob/40e56bab026552e089d7afcfeff08d2defb9c005/src/Development/AutomorphicArithmetic/CharacterBoxes/DesWin.lean\#L11-L31}{\bpcode{topDesWin}}\bpkickersep \href{https://github.com/Deicyde/riemann-hypothesis/blob/40e56bab026552e089d7afcfeff08d2defb9c005/src/Development/AbelianVarietyTheory/Definitions/GlobalTwisting/PairRowIffQdRows.lean\#L28-L32}{\bpcode{RankVal}}}
\bpgist{A chain of about $\sqrt x$ theta lifts, with divisibility, slope and covolume control, is chosen so that both loss exponents are fixed numbers: $e_D=2$, $e_L=0$.}
\begin{bpstatement}{est}
\bpdisplay{\begin{gathered}
\sqrt{x\delta_1}\le m\le2\sqrt{x\delta_1},\ \ \delta_1=x^{-\varepsilon/8}\ \Rightarrow\ m\approx x^{1/2-\varepsilon/16}; \\
\mu(A)\le\mu_{\max}(B)+\log C; \\
w^2=p\ \Rightarrow\ w\mid W_cV_0-W_0V_c .
\end{gathered}}
\end{bpstatement}
\bpcode{TopShape} calls \href{https://github.com/Deicyde/riemann-hypothesis/blob/40e56bab026552e089d7afcfeff08d2defb9c005/src/Development/AutomorphicArithmetic/CharacterBoxes/DesWin.lean\#L11-L31}{\bpcode{topDesWin}} with $C_{\rm abs}=2$, $C_0=3$, $e_D=2$, $e_L=0$, which gives the exponent $\tfrac12+9\varepsilon$. \href{https://github.com/Deicyde/riemann-hypothesis/blob/40e56bab026552e089d7afcfeff08d2defb9c005/src/Development/AbelianVarietyTheory/Definitions/GlobalTwisting/PairRowIffQdRows.lean\#L28-L32}{\bpcode{RankVal}} fixes the chain length; the announcement's “length of chain of theta lifts $\sqrt X$” is this parameter.

Three local tools are visible in isolation: the denominators of an iterated kernel clear by $\prod_j D_j$ (\bpcode{ChainInt}); a $2\times2$ minor is divisible by $w$ when $w^2=p$ (\href{https://github.com/Deicyde/riemann-hypothesis/blob/40e56bab026552e089d7afcfeff08d2defb9c005/src/Formalism/SharpRestriction/GrowthBounds/CoarseMultiplicity/KeyMinorDvd.lean\#L132-L155}{\bpcode{key\_\allowbreak{}minor\_\allowbreak{}dvd}}); and Bost's slope inequality bounds the slope of an injectively mapped lattice. The global gain requires these to hold simultaneously on one family.

\bpfigure{chain}{1.0}{The chain of theta lifts (schematic)}{Schematic: unitary groups U(X0) → U(X1) → ... → U(Xm) linked by theta lifts; each lift carries a denominator \(D_{j}\); the archive's \href{https://github.com/Deicyde/riemann-hypothesis/blob/40e56bab026552e089d7afcfeff08d2defb9c005/src/Development/AbelianVarietyTheory/Definitions/GlobalTwisting/PairRowIffQdRows.lean\#L28-L32}{\bpcode{RankVal}} puts m between \(\sqrt{x δ₁}\) and 2 \(\sqrt{x δ₁}\), \(δ\)₁ = \(x^{−ε/8}\).}
\begin{bpleantable}{est}
\bpleanrow{https://github.com/Deicyde/riemann-hypothesis/blob/40e56bab026552e089d7afcfeff08d2defb9c005/src/Development/AutomorphicArithmetic/CharacterBoxes/DesWin.lean\#L11-L31}{\bpcode{UpperTheory.\allowbreak{}Top.\allowbreak{}topDesWin}}{DesWin.lean, L11–31}{combine divisibility, slopes and windows}
\bpleanrow{https://github.com/Deicyde/riemann-hypothesis/blob/40e56bab026552e089d7afcfeff08d2defb9c005/src/Development/AutomorphicArithmetic/ThetaCorrespondence/ThetaPairLifts/ThirdPlaceSeries/TopShape.lean\#L33-L55}{\bpcode{UpperTheory.\allowbreak{}Top.\allowbreak{}shape\_\allowbreak{}closed\_\allowbreak{}link\_\allowbreak{}n\_\allowbreak{}of\_\allowbreak{}green}}{TopShape.lean, L33–55}{calls \bpcode{topDesWin} with \(e_{D}\) = 2, \(e_{L}\) = 0}
\bpleanrow{https://github.com/Deicyde/riemann-hypothesis/blob/40e56bab026552e089d7afcfeff08d2defb9c005/src/Development/AbelianVarietyTheory/Definitions/GlobalTwisting/PairRowIffQdRows.lean\#L28-L32}{\bpcode{GlobalTheory.\allowbreak{}EX.\allowbreak{}RankVal}}{PairRowIffQdRows.lean, L28–32}{the chain length and the other ranks}
\bpleanrow{https://github.com/Deicyde/riemann-hypothesis/blob/40e56bab026552e089d7afcfeff08d2defb9c005/src/Formalism/SharpRestriction/GrowthBounds/CoarseMultiplicity/KeyMinorDvd.lean\#L132-L155}{\bpcode{Formalism.\allowbreak{}MainNotions.\allowbreak{}ColumnSystems.\allowbreak{}Pf.\allowbreak{}key\_\allowbreak{}minor\_\allowbreak{}dvd}}{KeyMinorDvd.lean, L132–155}{a local source of divisibility}
\bpleanrow{https://github.com/Deicyde/riemann-hypothesis/blob/40e56bab026552e089d7afcfeff08d2defb9c005/src/Formalism/LatticeSlopes/LatticeSlopesProofs/SlopeInequality.lean\#L111-L125}{\bpcode{Formalism.\allowbreak{}F.\allowbreak{}slopeInequality\_\allowbreak{}of}}{SlopeInequality.lean, L111–125}{the lattice slope inequality}
\end{bpleantable}
\bpinterfaces{\bpiface{inputs}{definite-side data}{\href{https://github.com/Deicyde/riemann-hypothesis/blob/40e56bab026552e089d7afcfeff08d2defb9c005/src/Development/AutomorphicArithmetic/ThetaCorrespondence/ThetaPairLifts/ThirdPlaceSeries/SiegelEisensteinSecondTermHypotheses.lean\#L69-L73}{\bpcode{every\_\allowbreak{}character\_\allowbreak{}of\_\allowbreak{}theta\_\allowbreak{}integrals\_\allowbreak{}local\_\allowbreak{}matching}}}}{\bpiface{comparison}{\bpcode{TopG} (\(e_{D}\) = 2, \(e_{L}\) = 0)}{\bpcode{UpperTheory.\allowbreak{}Top.\allowbreak{}TopG}}}
\bpownednodes{11}{1\,\bptiericonplain{checked}, 9\,\bptiericonplain{read}, 1\,\bptiericonplain{motivation}}{\bpnodeitem{goal-shape}{read}{Apply the shape to the same chosen data}\bpnodeitem{shape-fields}{read}{Supply the actual fields of the shape}\bpnodeitem{shape}{read}{Choose a common geometric family}\bpnodeitem{family-design}{read}{Combine divisibility, slopes and windows}\bpnodeitem{divin}{read}{Control the local divisibility loss}\bpnodeitem{vv-sum}{read}{Sum the local losses}\bpnodeitem{rank}{read}{The geometric ranks}\bpnodeitem{minor}{checked}{A local source of divisibility}\bpnodeitem{empty-data}{read}{An existence check on a period identity}\bpnodeitem{slope}{read}{The lattice slope inequality}\bpnodeitem{paper-chain}{motivation}{The square-root multiplicity analogy}}
\bpstatus{partial}{\href{https://github.com/Deicyde/riemann-hypothesis/blob/40e56bab026552e089d7afcfeff08d2defb9c005/src/Formalism/SharpRestriction/GrowthBounds/CoarseMultiplicity/KeyMinorDvd.lean\#L132-L155}{\bpcode{key\_\allowbreak{}minor\_\allowbreak{}dvd}} passed in its 509-module closure. \href{https://github.com/Deicyde/riemann-hypothesis/blob/40e56bab026552e089d7afcfeff08d2defb9c005/src/Development/AutomorphicArithmetic/CharacterBoxes/DesWin.lean\#L11-L31}{\bpcode{topDesWin}} and the family construction are source-read here and included in the externally reported full build.}
\bpcheck{Nonempty compatible data: follow the witnesses of one geometric family through the test-vector grids, linked bundles and definite-side hypotheses.}{}
\bpannouncement{\bptlitem{1}{10/6/26}{Stepanov determinant}{Historical motivation}\bptlitem{2}{10/9/26}{Invariant differentials on a CM curve}{Historical calculation}\bptlitem{8}{11/15/26}{A chain with square-root length}{Quantitative comparison}\bptlitem{9}{11/15/26}{Norms, levels and arithmetic periods}{Supporting comparison}\bptlitem{10}{11/15/26}{Divisibility of repeated lifts}{Supporting comparison}\bptlitem{13}{10/6/26}{The slope inequality}{Mathematical statement matched}\bptlitem{15}{11/15/26}{A common integral structure across characters}{Supporting comparison}}

\bpcomponent{hosts}{est}{10}{Hosts for every character}{checked}
\bpkicker{\href{https://github.com/Deicyde/riemann-hypothesis/blob/40e56bab026552e089d7afcfeff08d2defb9c005/src/Development/AutomorphicArithmetic/CharacterBoxes/SupplyChi.lean\#L8-L16}{\bpcode{supplyChi}}}
\bpgist{For each primitive character $\psi$, the cohomological window produces arithmetic hosts of width $3/2$ at all large scales of a system of scales.}
\begin{bpstatement}{est}
\bpdisplay{\mathrm{VWinAt}\ \psi\ 401\ \Longrightarrow\ \text{hosts of width }3/2\text{ satisfying }(G_\chi\,\psi).\mathrm{req}.}
\end{bpstatement}
The supplier is applied separately to each character; a common host for all characters would be a stronger quantifier order and is not needed. The character requirement combines the cyclotomic-level condition \bpcode{EAtH} with the cohomological inequality for the host field.

\begin{bpleantable}{est}
\bpleanrow{https://github.com/Deicyde/riemann-hypothesis/blob/40e56bab026552e089d7afcfeff08d2defb9c005/src/Development/AutomorphicArithmetic/CharacterBoxes/SupplyChi.lean\#L8-L16}{\bpcode{UpperTheory.\allowbreak{}Top.\allowbreak{}supplyChi}}{SupplyChi.lean, L8–16}{\bpcode{VWinAt} \(ψ\) 401 ⇒ hosts of width 3/2}
\end{bpleantable}
\bpinterfaces{\bpiface{cohomology}{\bpcode{VWinAt} \(ψ\) 401}{\href{https://github.com/Deicyde/riemann-hypothesis/blob/40e56bab026552e089d7afcfeff08d2defb9c005/src/Development/AutomorphicArithmetic/CharacterBoxes/ExistsPageN.lean\#L44-L47}{\bpcode{UpperTheory.\allowbreak{}Top.\allowbreak{}VWinAt}}}}{\bpiface{comparison}{hosts, width 3/2}{\href{https://github.com/Deicyde/riemann-hypothesis/blob/40e56bab026552e089d7afcfeff08d2defb9c005/src/Development/AutomorphicArithmetic/CharacterBoxes/SupplyChi.lean\#L8-L16}{\bpcode{UpperTheory.\allowbreak{}Top.\allowbreak{}supplyChi}}}}
\bpownednodes{1}{1\,\bptiericonplain{checked}}{\bpnodeitem{supply}{checked}{Produce cyclic character hosts}}
\bpstatus{checked}{\href{https://github.com/Deicyde/riemann-hypothesis/blob/40e56bab026552e089d7afcfeff08d2defb9c005/src/Development/AutomorphicArithmetic/CharacterBoxes/SupplyChi.lean\#L8-L16}{\bpcode{supplyChi}} passed in the Dirichlet closure, with \href{https://github.com/Deicyde/riemann-hypothesis/blob/40e56bab026552e089d7afcfeff08d2defb9c005/src/Development/AutomorphicArithmetic/CharacterBoxes/ExistsPageN.lean\#L44-L47}{\bpcode{VWinAt}} \(ψ\) 401 as its hypothesis.}
\bpcheck{Its premise \href{https://github.com/Deicyde/riemann-hypothesis/blob/40e56bab026552e089d7afcfeff08d2defb9c005/src/Development/AutomorphicArithmetic/CharacterBoxes/ExistsPageN.lean\#L44-L47}{\bpcode{VWinAt}} \(ψ\) 401, for every primitive \(ψ\).}{}

\bpstratum{arith}{Arithmetic supply}{Galois representations, theta lifts and Iwasawa-type cohomology bounds that are meant to produce the families on which a prime sum is small.}{\bpstratcomp{cohomology}\bpstratcomp{primes}\bpstratcomp{galois}\bpstratcomp{inputs}}{\bpstratflow{inputs}{family}{definite-side data}\bpstratflow{cohomology}{hosts}{\bpcode{VWinAt} \(ψ\) 401}}
\bpcomponent{cohomology}{arith}{11}{The cohomological window}{partial}
\bpkicker{\href{https://github.com/Deicyde/riemann-hypothesis/blob/40e56bab026552e089d7afcfeff08d2defb9c005/src/Development/AutomorphicArithmetic/CharacterBoxes/ExistsPageN.lean\#L44-L47}{\bpcode{VWinAt ψ 401}}}
\bpgist{On a whole interval $[Y,Y^2]$, weighted lengths of locally trivial Galois cohomology of twisted modules are at most $\sqrt{\exp(x^{\varepsilon})}$, uniformly over the relevant fields.}
\begin{bpstatement}{arith}
\bpdisplay{\begin{gathered}
\sum_{\ell\in S}\log\ell\sum_{1\le h<\lfloor x^{401}\rfloor}s_{\ell,h}(E)\ \le\ \sqrt{\exp(x^{\varepsilon})} \\
\text{for all }x\in[Y,Y^2],\ \text{all finite }S,\ \text{all admissible }E .
\end{gathered}}
\end{bpstatement}
An $\ell$-adic determinant estimate gives at most $R$ centers; each center accounts for at most one exceptional twist $h\le Z$ (ultrametric separation). Finitely many exceptional twists forbid at most $R$ logarithmic intervals $[a,Ba)$.

A real-variable induction (\href{https://github.com/Deicyde/riemann-hypothesis/blob/40e56bab026552e089d7afcfeff08d2defb9c005/src/Development/AutomorphicArithmetic/AutomorphicVarieties/LayeredIsomorphisms/Window1.lean\#L6-L58}{\bpcode{window1}}) finds $[y,2y]$ avoiding all of them with $y\ge U/(2^{R+2}B^R)$; exponentiating gives $[Y,Y^2]$ with $X^{1/J}\le Y$, $J=4(2B)^R$. The bound on the sums comes from four counting cases, a Ferrero–Washington input and descent.

\bpfigure{windows}{1.0}{One interval survives all exceptional twists}{Computed with the induction of \href{https://github.com/Deicyde/riemann-hypothesis/blob/40e56bab026552e089d7afcfeff08d2defb9c005/src/Development/AutomorphicArithmetic/AutomorphicVarieties/LayeredIsomorphisms/Window1.lean\#L6-L58}{\bpcode{UpperTheory.\allowbreak{}V.\allowbreak{}Win.\allowbreak{}window1}} on the book's example: forbidden logarithmic intervals [a, B a) for a = 8, 90 with B = 2 inside [0, 128]; the recursion lowers the ceiling 128 → 32 and returns [16, 32], i.e. [Y, \(Y^{2}\)] = [\(2^{16}\), \(2^{32}\)].}
\begin{bpleantable}{arith}
\bpleanrow{https://github.com/Deicyde/riemann-hypothesis/blob/40e56bab026552e089d7afcfeff08d2defb9c005/src/Development/AutomorphicArithmetic/CharacterBoxes/ExistsPageN.lean\#L44-L47}{\bpcode{UpperTheory.\allowbreak{}Top.\allowbreak{}VWinAt}}{ExistsPageN.lean, L44–47}{the cohomological window (definition)}
\bpleanrow{https://github.com/Deicyde/riemann-hypothesis/blob/40e56bab026552e089d7afcfeff08d2defb9c005/src/Development/AutomorphicArithmetic/CharacterBoxes/ExistsPageN.lean\#L12-L17}{\bpcode{UpperTheory.\allowbreak{}Top.\allowbreak{}VVle}}{ExistsPageN.lean, L12–17}{weighted lengths of local–global kernels (definition)}
\bpleanrow{https://github.com/Deicyde/riemann-hypothesis/blob/40e56bab026552e089d7afcfeff08d2defb9c005/src/Development/AutomorphicArithmetic/AutomorphicVarieties/LayeredIsomorphisms/Window1.lean\#L6-L58}{\bpcode{UpperTheory.\allowbreak{}V.\allowbreak{}Win.\allowbreak{}window1}}{Window1.lean, L6–58}{}
\bpleanrow{https://github.com/Deicyde/riemann-hypothesis/blob/40e56bab026552e089d7afcfeff08d2defb9c005/src/Development/AutomorphicArithmetic/AutomorphicVarieties/LayeredIsomorphisms/Window.lean\#L7-L43}{\bpcode{UpperTheory.\allowbreak{}V.\allowbreak{}Win.\allowbreak{}window}}{Window.lean, L7–43}{}
\bpleanrow{https://github.com/Deicyde/riemann-hypothesis/blob/40e56bab026552e089d7afcfeff08d2defb9c005/src/Development/AutomorphicArithmetic/AutomorphicVarieties/ModularityArgument/ThresholdArguments/DenseThresholdCases/VwinOfRow.lean\#L4-L5}{\bpcode{UpperTheory.\allowbreak{}One.\allowbreak{}vwin\_\allowbreak{}of\_\allowbreak{}row}}{VwinOfRow.lean, L4–5}{the window from the Galois row}
\bpleanrow{https://github.com/Deicyde/riemann-hypothesis/blob/40e56bab026552e089d7afcfeff08d2defb9c005/src/Development/AutomorphicArithmetic/AutomorphicVarieties/ModularityArgument/ThresholdArguments/DenseThresholdCases/HVCore.lean\#L18-L22}{\bpcode{UpperTheory.\allowbreak{}V.\allowbreak{}HV.\allowbreak{}Core.\allowbreak{}of\_\allowbreak{}four\_\allowbreak{}alt}}{HVCore.lean, L18–22}{combine the cohomology counts}
\bpleanrow{https://github.com/Deicyde/riemann-hypothesis/blob/40e56bab026552e089d7afcfeff08d2defb9c005/src/Development/AutomorphicArithmetic/AutomorphicVarieties/LayeredIsomorphisms/ThmBabOf.lean\#L11-L103}{\bpcode{UpperTheory.\allowbreak{}V.\allowbreak{}Win.\allowbreak{}thmBab\_\allowbreak{}of}}{ThmBabOf.lean, L11–103}{Choose one set of tower constants}
\bpleanrow{https://github.com/Deicyde/riemann-hypothesis/blob/40e56bab026552e089d7afcfeff08d2defb9c005/src/Development/AutomorphicArithmetic/AutomorphicVarieties/AdditiveValuationNorms/ABoundOfIsTower.lean\#L7-L11}{\bpcode{UpperTheory.\allowbreak{}V.\allowbreak{}PropA.\allowbreak{}aBound\_\allowbreak{}of\_\allowbreak{}isTower}}{ABoundOfIsTower.lean, L7–11}{From a finitely generated tower to two length bounds}
\bpleanrow{https://github.com/Deicyde/riemann-hypothesis/blob/40e56bab026552e089d7afcfeff08d2defb9c005/src/Development/AutomorphicArithmetic/AutomorphicVarieties/LayeredIsomorphisms/ThmBabOfLemmaK.lean\#L9-L12}{\bpcode{UpperTheory.\allowbreak{}V.\allowbreak{}Win.\allowbreak{}thmBab\_\allowbreak{}of\_\allowbreak{}lemmaK}}{ThmBabOfLemmaK.lean, L9–12}{Assemble the arithmetic tower window}
\end{bpleantable}
\bpinterfaces{\bpiface{galois}{\bpcode{Row0b} → window}{\href{https://github.com/Deicyde/riemann-hypothesis/blob/40e56bab026552e089d7afcfeff08d2defb9c005/src/Development/AutomorphicArithmetic/AutomorphicVarieties/ModularityArgument/ThresholdArguments/DenseThresholdCases/VwinOfRow.lean\#L4-L5}{\bpcode{vwin\_\allowbreak{}of\_\allowbreak{}row}}}\bpiface{primes}{prime pairs, tower}{\bpcode{WinPair} \textperiodcentered{} \href{https://github.com/Deicyde/riemann-hypothesis/blob/40e56bab026552e089d7afcfeff08d2defb9c005/src/Formalism/FrobeniusOrbits/FrobeniusOrbitsProofs/ThickExists.lean\#L181-L187}{\bpcode{hostThickExists}}}}{\bpiface{hosts}{\bpcode{VWinAt} \(ψ\) 401}{\href{https://github.com/Deicyde/riemann-hypothesis/blob/40e56bab026552e089d7afcfeff08d2defb9c005/src/Development/AutomorphicArithmetic/CharacterBoxes/ExistsPageN.lean\#L44-L47}{\bpcode{UpperTheory.\allowbreak{}Top.\allowbreak{}VWinAt}}}}
\bpownednodes{25}{3\,\bptiericonplain{checked}, 2\,\bptiericonplain{covered}, 20\,\bptiericonplain{read}}{\bpnodeitem{vv}{read}{Weighted lengths of local-global kernels}\bpnodeitem{vwin-def}{covered}{The cohomological window}\bpnodeitem{vwin-row}{read}{Cohomological window from the row}\bpnodeitem{window-inputs}{read}{Supply the automorphic window inputs}\bpnodeitem{depth-window}{read}{Build the rank-N package}\bpnodeitem{lit-window}{read}{Window from rank-N automorphy}\bpnodeitem{patch-window}{read}{Pass through the descent hypotheses}\bpnodeitem{four-produce}{read}{Produce the four counting bounds}\bpnodeitem{four-use}{read}{Combine the cohomology counts}\bpnodeitem{tower-window-from-lemma-k}{checked}{Assemble the arithmetic tower window}\bpnodeitem{tower-window-constants}{checked}{Choose one set of tower constants}\bpnodeitem{tower-bound-from-rank}{checked}{From a finitely generated tower to two length bounds}\bpnodeitem{tower-operator-bound}{covered}{Control exceptional twists by the fixed layer}\bpnodeitem{tower-class-rank}{read}{Uniform ℓ-rank of class groups in the tower}\bpnodeitem{tower-fixed-growth}{read}{Supply the absolute bound for the reference layer}\bpnodeitem{counts-i}{read}{Case I gives one bound across coefficient primes}\bpnodeitem{counts-ii-iii}{read}{Cases II and III from the rank-four symplectic argument}\bpnodeitem{counts-v}{read}{Case V at twist one}\bpnodeitem{counts-common-integer}{read}{A nonzero integer controls one character orbit}\bpnodeitem{counts-window}{read}{Convert the joint window into \bpcode{VWin}}\bpnodeitem{counts-joint}{read}{Reuse the tower interval for the positive part}\bpnodeitem{counts-height}{read}{Evaluate the pointwise estimate on admissible windows}\bpnodeitem{counts-pointwise}{read}{Combine the two prime sums at a point}\bpnodeitem{counts-inside}{read}{Add contributions inside the exceptional set}\bpnodeitem{counts-outside}{read}{Bound the sum over arbitrary finite sets of other primes}}
\bpstatus{partial}{\href{https://github.com/Deicyde/riemann-hypothesis/blob/40e56bab026552e089d7afcfeff08d2defb9c005/src/Development/AutomorphicArithmetic/AutomorphicVarieties/LayeredIsomorphisms/Window1.lean\#L6-L58}{\bpcode{window1}} and window (the interval selection) passed at the archive pin, and so did the tower deductions \href{https://github.com/Deicyde/riemann-hypothesis/blob/40e56bab026552e089d7afcfeff08d2defb9c005/src/Development/AutomorphicArithmetic/AutomorphicVarieties/LayeredIsomorphisms/ThmBabOf.lean\#L11-L103}{\bpcode{thmBab\_\allowbreak{}of}} (one record of window constants for every cutoff Ct ≥ 1 from \bpcode{LemmaK}, \bpcode{FixedIn} and \bpcode{TwoIn}), \href{https://github.com/Deicyde/riemann-hypothesis/blob/40e56bab026552e089d7afcfeff08d2defb9c005/src/Development/AutomorphicArithmetic/AutomorphicVarieties/AdditiveValuationNorms/ABoundOfIsTower.lean\#L7-L11}{\bpcode{aBound\_\allowbreak{}of\_\allowbreak{}isTower}} (a tower cohomology bound from duality and finiteness premises) and \href{https://github.com/Deicyde/riemann-hypothesis/blob/40e56bab026552e089d7afcfeff08d2defb9c005/src/Development/AutomorphicArithmetic/AutomorphicVarieties/LayeredIsomorphisms/ThmBabOfLemmaK.lean\#L9-L12}{\bpcode{thmBab\_\allowbreak{}of\_\allowbreak{}lemmaK}}, each with its premises explicit. Those premises (Poitou–Tate and local Tate duality, local finiteness, Ferrero–Washington, \bpcode{LemmaK}) and the counting cases are source-read here.}
\bpcheck{Uniform cohomological bound: track finiteness and absolute constants from the four counting cases to the finite prime sums on the claimed windows.}{\bpaudit{four-use}{Uniform cohomological bound}{Track finiteness and the absolute constants from the four counting cases to the finite prime sums on the claimed windows.}}
\bpannouncement{\bptlitem{11}{10/6/26}{The tower of fields}{Supporting comparison}\bptlitem{17}{11/15/26}{Cohomology and a weakening of Vandiver}{Quantitative comparison}}

\bpcomponent{primes}{arith}{12}{Auxiliary primes and Behrend towers}{checked}
\bpkicker{\href{https://github.com/Deicyde/riemann-hypothesis/blob/40e56bab026552e089d7afcfeff08d2defb9c005/src/Development/AutomorphicArithmetic/MeanSquareBounds/pairEx.lean\#L112-L169}{\bpcode{QH.\allowbreak{}pairEx}}\bpkickersep \href{https://github.com/Deicyde/riemann-hypothesis/blob/40e56bab026552e089d7afcfeff08d2defb9c005/src/Formalism/FrobeniusOrbits/FrobeniusOrbitsProofs/ThickExists.lean\#L181-L187}{\bpcode{hostThickExists}}}
\bpgist{Two auxiliary primes with prescribed congruences and square-root prime-sum bounds for their quadratic twists, and an additive-combinatorial (Behrend) choice of the tower.}
\begin{bpstatement}{arith}
\bpdisplay{\begin{gathered}
q\equiv3\ (8),\ \ d\equiv7\ (16),\ \ q\mid d+1, \\
|g_\chi(y)|\le y^{1/2+\varepsilon/2}\ \text{on }[Y,Y^A]; \\
|S|\ge\frac nL\,e^{-C\sqrt{\log n\,\log L}} .
\end{gathered}}
\end{bpstatement}
With $z=Y^{\varepsilon/5}$ the primes satisfy $z/2<\log q\le z$ and $z(\log Y)^2/2<\log d\le z(\log Y)^2$; $d+1$ is divisible by $q$ and by every odd prime up to $z$. They are chosen sequentially from positive shares of candidates, and may change with $Y$.

Behrend: integer vectors of fixed norm, encoded in a large base, avoid short affine relations because a line meets a sphere at most twice.

\bpfigure{behrend}{0.86}{A Behrend digit sphere}{Computed: the vectors of \{0,1,2\}\(^{3}\) with squared norm 5 are the six permutations of (0,1,2); they lie on the circle where the sphere meets the plane x+y+z = 3. Their base-24 encodings avoid every relation j(s2-s1) = i(s3-s1) with 1 ≤ |i|,|j| ≤ 2, i ≠ j: 1,440 checks.}
\begin{bpleantable}{arith}
\bpleanrow{https://github.com/Deicyde/riemann-hypothesis/blob/40e56bab026552e089d7afcfeff08d2defb9c005/src/Development/AutomorphicArithmetic/MeanSquareBounds/pairEx.lean\#L112-L169}{\bpcode{QH.\allowbreak{}pairEx}}{pairEx.lean, L112–169}{}
\bpleanrow{https://github.com/Deicyde/riemann-hypothesis/blob/40e56bab026552e089d7afcfeff08d2defb9c005/src/Development/AutomorphicArithmetic/MeanSquareBounds/PairExists.lean\#L9-L48}{two successive choices}{PairExists.lean, L9–48}{}
\bpleanrow{https://github.com/Deicyde/riemann-hypothesis/blob/40e56bab026552e089d7afcfeff08d2defb9c005/src/Formalism/FrobeniusOrbits/FrobeniusOrbitsProofs/ThickExists.lean\#L85-L95}{\bpcode{Formalism.\allowbreak{}FrobeniusOrbits.\allowbreak{}ThickExists.\allowbreak{}noRel\_\allowbreak{}image}}{ThickExists.lean, L85–95}{digit spheres avoid short relations}
\bpleanrow{https://github.com/Deicyde/riemann-hypothesis/blob/40e56bab026552e089d7afcfeff08d2defb9c005/src/Formalism/FrobeniusOrbits/FrobeniusOrbitsProofs/ThickExists.lean\#L181-L187}{\bpcode{hostThickExists}}{ThickExists.lean, L181–187}{}
\end{bpleantable}
\bpinterfaces{}{\bpiface{cohomology}{prime pairs, tower}{\bpcode{WinPair} \textperiodcentered{} \href{https://github.com/Deicyde/riemann-hypothesis/blob/40e56bab026552e089d7afcfeff08d2defb9c005/src/Formalism/FrobeniusOrbits/FrobeniusOrbitsProofs/ThickExists.lean\#L181-L187}{\bpcode{hostThickExists}}}}
\bpownednodes{2}{1\,\bptiericonplain{checked}, 1\,\bptiericonplain{motivation}}{\bpnodeitem{behrend}{checked}{Behrend’s collision-free indices}\bpnodeitem{paper-behrend}{motivation}{The tower selection in the timeline}}
\bpstatus{checked}{\href{https://github.com/Deicyde/riemann-hypothesis/blob/40e56bab026552e089d7afcfeff08d2defb9c005/src/Development/AutomorphicArithmetic/MeanSquareBounds/pairEx.lean\#L112-L169}{\bpcode{QH.\allowbreak{}pairEx}} passed in its 293-module closure; the Behrend landmarks (\bpcode{sphere\_\allowbreak{}rel}, \href{https://github.com/Deicyde/riemann-hypothesis/blob/40e56bab026552e089d7afcfeff08d2defb9c005/src/Formalism/FrobeniusOrbits/FrobeniusOrbitsProofs/ThickExists.lean\#L85-L95}{\bpcode{noRel\_\allowbreak{}image}}, \href{https://github.com/Deicyde/riemann-hypothesis/blob/40e56bab026552e089d7afcfeff08d2defb9c005/src/Formalism/FrobeniusOrbits/FrobeniusOrbitsProofs/ThickExists.lean\#L181-L187}{\bpcode{hostThickExists}}) passed in the 38-module landmark closure.}
\bpcheck{How the chosen primes feed the window for a fixed target character.}{}
\bpannouncement{\bptlitem{19}{10/19/26}{A Behrend set in the tower construction}{Supporting comparison}}

\bpcomponent{galois}{arith}{13}{The Galois row}{read}
\bpkicker{\href{https://github.com/Deicyde/riemann-hypothesis/blob/40e56bab026552e089d7afcfeff08d2defb9c005/src/Development/AutomorphicArithmetic/Foundations/GaloisRepresentations/TypeConditions/WeightsConjBis.lean\#L69-L73}{\bpcode{Row0b}}}
\bpgist{For regular algebraic polarized cuspidal $\pi$ on $\mathrm{GL}_a$ over a CM field, a continuous semisimple $\ell$-adic representation whose Frobenius polynomials match the Satake parameters of $\pi$.}
\begin{bpstatement}{arith}
\bpdisplay{\begin{gathered}
\iota\bigl(\operatorname{charpoly}\rho(\sigma^{-1})\bigr)=\prod_{\alpha\in B_v}\bigl(X-q^{(a-1)/2}\alpha\bigr)\quad\text{for all but finitely many }v, \\
\rho:G_E\to\mathrm{GL}_a(\overline{\mathbb Q}_\ell).
\end{gathered}}
\end{bpstatement}
The construction starts from a geometric seed and a high-rank theta lift (\href{https://github.com/Deicyde/riemann-hypothesis/blob/40e56bab026552e089d7afcfeff08d2defb9c005/src/Development/AutomorphicArithmetic/Foundations/GaloisRepresentations/ThetaParseval/ThetaKernels/TLiftK.lean\#L4-L5}{\bpcode{tLiftK}}), treats a restricted case, removes the prescribed polarization by twisting, and descends from infinitely many quadratic CM extensions $K=E(\sqrt{-p})$ by index-two patching of trace identities.

The exported row also records local–global compatibility away from $\ell$ and crystalline/de Rham comparisons above $\ell$. A rank-one example shows that the hypotheses are inhabited.

\bpfigure{galois}{0.86}{The Galois row by quadratic descent (schematic)}{Schematic: representations \(r_{K}\) over infinitely many quadratic CM extensions K = E(\(\sqrt{−p}\)) are patched to one \(ρ\) over E whose Frobenius polynomials match the Satake parameters of \(π\).}
\begin{bpleantable}{arith}
\bpleanrow{https://github.com/Deicyde/riemann-hypothesis/blob/40e56bab026552e089d7afcfeff08d2defb9c005/src/Development/AutomorphicArithmetic/Foundations/GaloisRepresentations/TypeConditions/WeightsConjBis.lean\#L69-L73}{\bpcode{UpperTheory.\allowbreak{}GalStar.\allowbreak{}Row0b}}{WeightsConjBis.lean, L69–73}{the local–global Galois row (definition)}
\bpleanrow{https://github.com/Deicyde/riemann-hypothesis/blob/40e56bab026552e089d7afcfeff08d2defb9c005/src/Development/AutomorphicArithmetic/Foundations/GaloisRepresentations/IntegralStructures/RowSrc.lean\#L6-L6}{\bpcode{UpperTheory.\allowbreak{}Gal.\allowbreak{}Int.\allowbreak{}galoisLocalGlobalStrong\_\allowbreak{}recParameter}}{RowSrc.lean, L6}{producer of the row}
\bpleanrow{https://github.com/Deicyde/riemann-hypothesis/blob/40e56bab026552e089d7afcfeff08d2defb9c005/src/Development/AutomorphicArithmetic/Foundations/GaloisRepresentations/IntegralStructures/GaloisRepresentationsOfThetaLift.lean\#L6-L6}{\bpcode{UpperTheory.\allowbreak{}Gal.\allowbreak{}Int.\allowbreak{}galoisRepresentations\_\allowbreak{}of\_\allowbreak{}thetaLift}}{GaloisRepresentationsOfThetaLift.lean, L6}{Galois representations from theta lifting}
\bpleanrow{https://github.com/Deicyde/riemann-hypothesis/blob/40e56bab026552e089d7afcfeff08d2defb9c005/src/Development/AutomorphicArithmetic/Foundations/GaloisRepresentations/ThetaParseval/ThetaKernels/TLiftK.lean\#L4-L5}{\bpcode{UpperTheory.\allowbreak{}TP.\allowbreak{}K.\allowbreak{}tLiftK}}{TLiftK.lean, L4–5}{theta-lift existence at high rank (theorem)}
\bpleanrow{https://github.com/Deicyde/riemann-hypothesis/blob/40e56bab026552e089d7afcfeff08d2defb9c005/src/Development/AutomorphicArithmetic/Foundations/GaloisRepresentations/ThetaParseval/HalfIndexLattices/TLiftK.lean\#L12-L23}{\bpcode{UpperTheory.\allowbreak{}TP.\allowbreak{}TLiftK}}{TLiftK.lean, L12–23}{quantifiers of the theta lift (definition)}
\bpleanrow{https://github.com/Deicyde/riemann-hypothesis/blob/40e56bab026552e089d7afcfeff08d2defb9c005/src/Development/AutomorphicArithmetic/Foundations/GaloisRepresentations/IntegralStructures/GeometricFlagBases/WeightOneGeometricRealization.lean\#L11-L11}{\bpcode{UpperTheory.\allowbreak{}Gal.\allowbreak{}Int.\allowbreak{}Geo.\allowbreak{}geoOne\_\allowbreak{}holds}}{WeightOneGeometricRealization.lean, L11}{weight-one geometric realization}
\end{bpleantable}
\bpinterfaces{}{\bpiface{inputs}{\bpcode{Row0b} (two fields)}{\href{https://github.com/Deicyde/riemann-hypothesis/blob/40e56bab026552e089d7afcfeff08d2defb9c005/src/Development/AutomorphicArithmetic/CharacterBoxes/MainConclusion.lean\#L7-L8}{\bpcode{everyCharacter\_\allowbreak{}of\_\allowbreak{}row}}}\bpiface{cohomology}{\bpcode{Row0b} → window}{\href{https://github.com/Deicyde/riemann-hypothesis/blob/40e56bab026552e089d7afcfeff08d2defb9c005/src/Development/AutomorphicArithmetic/AutomorphicVarieties/ModularityArgument/ThresholdArguments/DenseThresholdCases/VwinOfRow.lean\#L4-L5}{\bpcode{vwin\_\allowbreak{}of\_\allowbreak{}row}}}}
\bpownednodes{6}{6\,\bptiericonplain{read}}{\bpnodeitem{row-def}{read}{The local–global Galois row}\bpnodeitem{row-prod}{read}{Producer of the Galois row}\bpnodeitem{row-geo}{read}{Galois representations from theta lifting}\bpnodeitem{tlift}{read}{Theta-lift existence at high rank}\bpnodeitem{geo-one}{read}{Weight-one geometric realization}\bpnodeitem{tlift-def}{read}{Quantifiers of the theta lift}}
\bpstatus{read}{Source-read against 62 passages in 36 modules; included in the externally reported full build; not separately replayed by this audit.}
\bpcheck{Unfold the row, the local normalization convention and the theta-lift producer together; check all quantified places, twists and weight restrictions.}{\bpaudit{row-prod}{Source comparison of the Galois input}{Unfold the row, the local-normalization convention and the theta-lift producer together. Check all quantified places, twists and weight restrictions.}}

\bpcomponent{inputs}{arith}{14}{Seven inputs at the decisive call}{read}
\bpkicker{\href{https://github.com/Deicyde/riemann-hypothesis/blob/40e56bab026552e089d7afcfeff08d2defb9c005/src/Development/AutomorphicArithmetic/ThetaCorrespondence/ThetaPairLifts/ThirdPlaceSeries/EveryCharacterOfLiteratureOfVWin.lean\#L810-L813}{\bpcode{every\_\allowbreak{}character\_\allowbreak{}of\_\allowbreak{}literature\_\allowbreak{}of\_\allowbreak{}vWin}}}
\bpgist{Unitary-group and local theta packages, agreement of theta integrals, local matching, a covolume bound, local Langlands parameters and the window: the seven arguments of the arithmetic theorem.}
\begin{bpstatement}{arith}
\bpdisplay{P,\ F,\ h_A,\ h_M,\ h_X,\ h_L,\ h_V\ \Longrightarrow\ \text{nonempty definite-side data}\ +\ \text{Green input}\ \Longrightarrow\ \mathrm{TopG}.}
\end{bpstatement}
\href{https://github.com/Deicyde/riemann-hypothesis/blob/40e56bab026552e089d7afcfeff08d2defb9c005/src/Development/AutomorphicArithmetic/CharacterBoxes/MainConclusion.lean\#L7-L8}{\bpcode{everyCharacter\_\allowbreak{}of\_\allowbreak{}row}} fills two fields of the literature record with the same Galois row, adds the existence of local Langlands parameters, and obtains the cohomological window from the same row. The next line supplies the row by name, so the half-plane theorem has no explicit premise.

A name such as \href{https://github.com/Deicyde/riemann-hypothesis/blob/40e56bab026552e089d7afcfeff08d2defb9c005/src/Development/AutomorphicArithmetic/ThetaCorrespondence/ThetaPairLifts/ThirdPlaceSeries/EveryCharacterOfLiteratureOfVWin.lean\#L666-L672}{\bpcode{published\_\allowbreak{}results\_\allowbreak{}eigen}} is where to start verifying what the term establishes, not an explanation of it.

\begin{bpleantable}{arith}
\bpleanrow{https://github.com/Deicyde/riemann-hypothesis/blob/40e56bab026552e089d7afcfeff08d2defb9c005/src/Development/AutomorphicArithmetic/CharacterBoxes/MainConclusion.lean\#L7-L8}{\bpcode{UpperTheory.\allowbreak{}FinalTheorems.\allowbreak{}everyCharacter\_\allowbreak{}of\_\allowbreak{}row}}{MainConclusion.lean, L7–8}{the half-plane theorem from one row}
\bpleanrow{https://github.com/Deicyde/riemann-hypothesis/blob/40e56bab026552e089d7afcfeff08d2defb9c005/src/Development/AutomorphicArithmetic/CharacterBoxes/MainConclusion.lean\#L10-L11}{\bpcode{UpperTheory.\allowbreak{}FinalTheorems.\allowbreak{}everyCharacter}}{MainConclusion.lean, L10–11}{closed half-plane assembly}
\bpleanrow{https://github.com/Deicyde/riemann-hypothesis/blob/40e56bab026552e089d7afcfeff08d2defb9c005/src/Development/AutomorphicArithmetic/ThetaCorrespondence/ThetaPairLifts/ThirdPlaceSeries/EveryCharacterOfLiteratureOfVWin.lean\#L810-L813}{\bpcode{UpperTheory.\allowbreak{}LiteratureTheorems.\allowbreak{}every\_\allowbreak{}character\_\allowbreak{}of\_\allowbreak{}literature\_\allowbreak{}of\_\allowbreak{}vWin}}{EveryCharacterOfLiteratureOfVWin.lean, L810–813}{assemble the automorphic inputs}
\bpleanrow{https://github.com/Deicyde/riemann-hypothesis/blob/40e56bab026552e089d7afcfeff08d2defb9c005/src/Development/AutomorphicArithmetic/ThetaCorrespondence/ThetaPairLifts/ThirdPlaceSeries/SiegelEisensteinSecondTermHypotheses.lean\#L69-L73}{\bpcode{UpperTheory.\allowbreak{}Lit.\allowbreak{}every\_\allowbreak{}character\_\allowbreak{}of\_\allowbreak{}theta\_\allowbreak{}integrals\_\allowbreak{}local\_\allowbreak{}matching}}{SiegelEisensteinSecondTermHypotheses.lean, L69–73}{seven inputs for every character}
\end{bpleantable}
\bpinterfaces{\bpiface{galois}{\bpcode{Row0b} (two fields)}{\href{https://github.com/Deicyde/riemann-hypothesis/blob/40e56bab026552e089d7afcfeff08d2defb9c005/src/Development/AutomorphicArithmetic/CharacterBoxes/MainConclusion.lean\#L7-L8}{\bpcode{everyCharacter\_\allowbreak{}of\_\allowbreak{}row}}}}{\bpiface{family}{definite-side data}{\href{https://github.com/Deicyde/riemann-hypothesis/blob/40e56bab026552e089d7afcfeff08d2defb9c005/src/Development/AutomorphicArithmetic/ThetaCorrespondence/ThetaPairLifts/ThirdPlaceSeries/SiegelEisensteinSecondTermHypotheses.lean\#L69-L73}{\bpcode{every\_\allowbreak{}character\_\allowbreak{}of\_\allowbreak{}theta\_\allowbreak{}integrals\_\allowbreak{}local\_\allowbreak{}matching}}}}
\bpownednodes{16}{16\,\bptiericonplain{read}}{\bpnodeitem{every}{read}{Closed half-plane assembly}\bpnodeitem{from-row}{read}{The half-plane theorem from one row}\bpnodeitem{rec}{read}{Existence of local Langlands parameters}\bpnodeitem{lit}{read}{Assemble the automorphic inputs}\bpnodeitem{open-record}{read}{Three fields in the literature input}\bpnodeitem{eigen-package}{read}{The unitary and eigenvalue package}\bpnodeitem{theta-package}{read}{The local theta package}\bpnodeitem{hodge-package}{read}{The Hodge-bundle package}\bpnodeitem{local-rec}{read}{The local parameter input at the call site}\bpnodeitem{local-rec-proj}{read}{Read the parameter field}\bpnodeitem{green}{read}{Regular characters yield the Green input}\bpnodeitem{linked-two-three}{read}{Bind the estimates at 2 and 3}\bpnodeitem{theta-match}{read}{Match global integrals and local data}\bpnodeitem{covol}{read}{Produce the covolume bound}\bpnodeitem{seven}{read}{Seven inputs for every character}\bpnodeitem{definite}{read}{Build a nonempty definite-side package}}
\bpstatus{read}{Source-read; included in the externally reported full build; not separately replayed by this audit.}
\bpcheck{Source comparison of each named input with the mathematics it is meant to encode.}{\bpaudit{definite}{Nonempty compatible data}{Follow the witnesses for the same geometric family through the test-vector grids, linked bundles and nonempty definite-side hypotheses.}}
\bpannouncement{\bptlitem{7}{10/14/26}{Interlacing weights and periods}{Supporting comparison}\bptlitem{12}{11/15/26}{Seesaw and Bessel periods}{Supporting comparison}\bptlitem{18}{11/15/26}{Covolumes, lattice indices and congruences}{Supporting comparison}}


%% ------------------------------------------------------------------ 7. cuts and skeleton
\clearpage
\section{The verified cuts and the Lean skeleton}\label{sec:cuts}

Two implications have been machine-checked with the remaining inputs as explicit hypotheses. Each is a
Lean theorem, written by the reading project against the archive, whose only open inputs are the premise
displayed; the declarations it calls, directly or transitively, were elaborated and kernel-checked with
it (tier \bptiericonplain{covered} in Section~\ref{sec:table}). Together they cut the proof into the part
downstream of the premises, which is machine-checked, and the arithmetic supply, which is partly checked
and partly source-read. In the picture of Section~\ref{sec:picture} they are the dashed frontiers.
% Generated by scripts/build_blueprint_tex.py from data/model.json -- do not edit by hand.
\begin{bpcut}{cut-zeta}{I}{two-kernels}{One-sided prime windows imply RH}
\bpcutpremise{\forall\theta\in(\tfrac12,1)\ \exists C\ \forall Y_0\ \exists Y\ge Y_0\ \forall y\in[Y,Y^{3/2}]:\ G(y)\le C\,y^{\theta}}
\bpcutconclusion{Mathlib's \href{https://github.com/leanprover-community/mathlib4/blob/0df444a360eaa60ab8c11dca51a86af692955474/Mathlib/NumberTheory/LSeries/RiemannZeta.lean\#L182-L186}{\bpcode{RiemannHypothesis}}}
\bpcutmeta{\bpcode{Reading.\allowbreak{}WindowBridge.\allowbreak{}riemannHypothesis\_\allowbreak{}of\_\allowbreak{}prime\_\allowbreak{}windows}}{references/archive-bridges/WindowBridge.lean}{Fresh Lean replay over 10,530 modules (766,938 constants) and the independent Rust kernel NanoDA on 61,697 exported declarations; 31 original modules; axioms \bpcode{propext}, \bpcode{Classical.\allowbreak{}choice}, \bpcode{Quot.\allowbreak{}sound}.}{\bpcompref{reflection}, \bpcompref{statistic}, \bpcompref{residues}, \bpcompref{criterion}}
\end{bpcut}
\bpquotecaption{references/archive-bridges/WindowBridge.lean}{\bpcode{OneSidedWindows}, \bpcode{zetaFamily}, \bpcode{riemannHypothesis\_\allowbreak{}of\_\allowbreak{}prime\_\allowbreak{}windows}}
\begin{bpleanquote}{references/archive-bridges/WindowBridge.lean}{8}{23}
\begin{bplean}[firstnumber=8]
\LCom{/-- The constant may depend on the exponent. A chosen window must work at every real point. -/}
\LKw{def} \LDecl{OneSidedWindows} (g : ℝ → ℝ) : \LSort{Prop} :=
  ∀ θ : ℝ, 1 / 2 < θ → θ < 1 → ∃ C : ℝ,
    ∀ Y₀ : ℝ, ∃ Y : ℝ, Y₀ ≤ Y ∧
      ∀ y ∈ Set.Icc Y (Y ^ (3 / 2 : ℝ)), g y ≤ C * y ^ θ

\LCom{/-- The family has one index, the character modulo one and a positive coefficient of one. -/}
\LKw{def} \LDecl{zetaFamily} : Deduction.CharFamily \LKw{where}
  ι := Unit
  fin := inferInstance
  N := \LKw{fun} _ => 1
  nz := \LKw{fun} _ => inferInstance
  χ := \LKw{fun} _ => (1 : DirichletCharacter ℂ 1)
  a := \LKw{fun} _ => 1
  a_pos := \LKw{fun} _ => zero_lt_one
  adm := \LKw{fun} _ => Or.inl \LTac{rfl}
\end{bplean}
\end{bpleanquote}
\begin{bpleanquote}{references/archive-bridges/WindowBridge.lean}{37}{53}
\begin{bplean}[firstnumber=37]
\LCom{/-- A one-sided estimate on arbitrarily late three-halves windows implies Mathlib's RH. -/}
\LKw{theorem} \LDecl{riemannHypothesis_of_prime_windows} (h : OneSidedWindows Deduction.G) : RiemannHypothesis := \LKw{by}
  \LTac{apply} Formalism.VW.rh_of_right
  \LTac{intro} s hs _
  \LKw{have} hhalf : (1 / 2 : ℝ) < min s.re 1 := lt_min hs (\LKw{by} \LTac{norm_num})
  \LTac{obtain} ⟨θ, hθhalf, hθmin⟩ := \LTac{exists_between} hhalf
  \LKw{have} hθs : θ < s.re := lt_of_lt_of_le hθmin (min_le_left _ _)
  \LKw{have} hθ1 : θ < 1 := lt_of_lt_of_le hθmin (min_le_right _ _)
  \LTac{obtain} ⟨C, hC⟩ := h θ hθhalf hθ1
  \LKw{have} hscale : (1 : ℝ) + 1 / 2 = 3 / 2 := \LKw{by} \LTac{norm_num}
  \LKw{have} hfree : Deduction.ZeroFree zetaFamily θ :=
    Deduction.zeroFree_of_rpow_window zetaFamily (θ := θ) (δ := 1 / 2) (C := C)
      (lt_trans (\LKw{by} \LTac{norm_num} : (0 : ℝ) < 1 / 2) hθhalf) hθ1 (\LKw{by} \LTac{norm_num})
      (\LKw{by} \LTac{simpa} \LKw{only} [hscale, gm_zetaFamily] \LKw{using} hC)
  \LKw{have} hz : DirichletCharacter.LFunction (1 : DirichletCharacter ℂ 1) s ≠ 0 :=
    hfree () s hθs
  \LTac{rwa} [DirichletCharacter.LFunction_modOne_eq] \LKw{at} hz
\end{bplean}
\end{bpleanquote}

\begin{bpcut}{cut-dirichlet}{II}{checked}{\bpcode{TopG} and the cohomological windows imply RH and Dirichlet GRH}
\bpcutpremise{\mathrm{TopG}\ \wedge\ \forall\psi\ \text{primitive}:\ \mathrm{VWinAt}\ \psi\ 401}
\bpcutconclusion{RH ∧ half-plane nonvanishing ∧ critical-strip zeros, for every positive modulus}
\bpcutmeta{\bpcode{Reading.\allowbreak{}DirichletBridge.\allowbreak{}endpoints\_\allowbreak{}of\_\allowbreak{}topG\_\allowbreak{}vWin401}}{references/archive-bridges/DirichletBridge.lean}{1,421 original modules at Lean 4.33.1; seven type/axiom reports; six defining modules replayed (Helios job 26647685).}{\bpcompref{reflection}, \bpcompref{descent}, \bpcompref{statistic}, \bpcompref{residues}, \bpcompref{criterion}, \bpcompref{choreography}, \bpcompref{comparison}, \bpcompref{hosts}}
\end{bpcut}
\bpquotecaption{references/archive-bridges/DirichletBridge.lean}{\bpcode{primitiveZeroFree\_\allowbreak{}of\_\allowbreak{}topG\_\allowbreak{}vWin401}, \bpcode{endpoints\_\allowbreak{}of\_\allowbreak{}topG\_\allowbreak{}vWin401}}
\begin{bpleanquote}{references/archive-bridges/DirichletBridge.lean}{11}{26}
\begin{bplean}[firstnumber=11]
\LCom{/-- The cohomology cutoff is 401 for each primitive character. The supplied}
\LCom{host width is the same real number 3/2 for all these characters. -/}
\LKw{theorem} \LDecl{primitiveZeroFree_of_topG_vWin401} (hG : UpperTheory.Top.TopG)
    (hV : ∀ ψ : PChar, VWinAt ψ 401) : ZeroFreeAll := \LKw{by}
  \LTac{apply} UpperTheory.Top.zeroFreeAll_of_topG (la := (3 / 2 : ℝ)) (\LKw{by} \LTac{norm_num}) hG
  \LTac{intro} ψ
  \LTac{exact} UpperTheory.Top.supplyChi ψ (hV ψ)

\LCom{/-- Pass to arbitrary positive moduli, then use the archived critical-strip}
\LCom{and zeta reductions. The right-half-plane predicate retains its pole exception. -/}
\LKw{theorem} \LDecl{endpoints_of_topG_vWin401} (hG : UpperTheory.Top.TopG)
    (hV : ∀ ψ : PChar, VWinAt ψ 401) :
    RiemannHypothesis ∧ DirichletLNonvanishingRightOfHalf ∧ DirichletLZerosOnLine := \LKw{by}
  \LKw{have} hχ : DirichletLNonvanishingRightOfHalf :=
    UpperTheory.Top.imprimitive (primitiveZeroFree_of_topG_vWin401 hG hV)
  \LTac{exact} ⟨Claim.zeta hχ, hχ, Claim.strip hχ⟩
\end{bplean}
\end{bpleanquote}


\subsection{The architecture as a checked skeleton}\label{sec:skeleton}
This repository's own Lean library (\SkeletonLean, no imports) restates the architecture over abstract
types: each component becomes a hypothesis, a field of a structure, and the theorems prove only how the
components compose. It is a map, not a proof of the Riemann hypothesis. A toy model at the end of the
file (\bpcode{toyArchitecture}, not quoted) satisfies every hypothesis with zeros present, all on the
line, so the skeleton is not vacuous; it says nothing about which model is the real one. The build
\SkeletonBuild{} with \SkeletonWarnings{} warnings; the axioms of all \SkeletonDecls{} audited
declarations are listed at the end of this section.

\subsubsection{The objects} Points, characters and zeros, with the predicates the endpoints talk about.
% Generated by scripts/build_blueprint_tex.py from data/model.json -- do not edit by hand.
\begin{bpleanquote}{lean/Atlas/Architecture.lean}{31}{60}
\begin{bplean}[firstnumber=31]
\LCom{/-- The objects the endpoints talk about. `Point` stands for complex numbers `s`, `Char` for}
\LCom{Dirichlet characters of all positive moduli, `Zero χ s` for `L(s, χ) = 0`. Lean's `L`-functions}
\LCom{are total functions, so a value exists even at the principal pole. -/}
\LKw{structure} \LDecl{Plane} \LKw{where}
  Point : \LSort{Type}
  Char : \LSort{Type}
  Zero : Char → Point → \LSort{Prop}
  \LCom{/-- `Re s > 1/2`. -/}
  Right : Point → \LSort{Prop}
  \LCom{/-- `Re s = 1/2`. -/}
  Line : Point → \LSort{Prop}
  \LCom{/-- `Re s < 1/2`. -/}
  Left : Point → \LSort{Prop}
  \LCom{/-- `0 < Re s < 1`. -/}
  Strip : Point → \LSort{Prop}
  \LCom{/-- `(χ, s) = (1_N, 1)`: the principal character at the pole. -/}
  Pole : Char → Point → \LSort{Prop}
  Primitive : Char → \LSort{Prop}
  \LCom{/-- The primitive character inducing `χ`. -/}
  prim : Char → Char
  \LCom{/-- `χ ↦ χ⁻¹`. -/}
  inv : Char → Char
  \LCom{/-- `s ↦ 1 - s`. -/}
  refl : Point → Point
  \LCom{/-- The character of modulus one, whose `L`-function is `ζ`. -/}
  zeta : Char
  \LCom{/-- `s ∈ \{-2, -4, -6, …\}`. -/}
  Trivial : Point → \LSort{Prop}
  \LCom{/-- `s = 1`. -/}
  One : Point → \LSort{Prop}
\end{bplean}
\end{bpleanquote}

\subsubsection{The target propositions} RH as in Mathlib, the half-plane theorem, its primitive form and
Dirichlet GRH in the strip.
% Generated by scripts/build_blueprint_tex.py from data/model.json -- do not edit by hand.
\begin{bpleanquote}{lean/Atlas/Architecture.lean}{64}{82}
\begin{bplean}[firstnumber=64]
\LCom{/-- Mathlib's `RiemannHypothesis`, transported to the abstract plane: every zero of `ζ` that is}
\LCom{neither trivial nor at `s = 1` lies on the critical line. -/}
\LKw{def} \LDecl{RH} : \LSort{Prop} :=
  ∀ s, P.Zero P.zeta s → ¬ P.Trivial s → ¬ P.One s → P.Line s

\LCom{/-- The half-plane theorem (archive: `DirichletLNonvanishingRightOfHalf`): no `L(s, χ)` vanishes}
\LCom{for `Re s > 1/2`, the principal pole excepted. -/}
\LKw{def} \LDecl{HalfPlane} : \LSort{Prop} :=
  ∀ χ s, P.Right s → ¬ P.Pole χ s → ¬ P.Zero χ s

\LCom{/-- The primitive half-plane statement (archive: `ZeroFreeAll`). It quantifies over primitive}
\LCom{characters only and has no pole exception, because Lean's totalized `ζ` takes a nonzero value}
\LCom{at `s = 1`. -/}
\LKw{def} \LDecl{PrimitiveHalfPlane} : \LSort{Prop} :=
  ∀ ψ s, P.Primitive ψ → P.Right s → ¬ P.Zero ψ s

\LCom{/-- Dirichlet GRH in the open strip (archive: `DirichletLZerosOnLine`). -/}
\LKw{def} \LDecl{StripGRH} : \LSort{Prop} :=
  ∀ χ s, P.Strip s → P.Zero χ s → P.Line s
\end{bplean}
\end{bpleanquote}

\subsubsection{The destination as hypotheses} Reflection (functional equations) and conductor descent
(Euler factors), components \bpcompref{reflection} and \bpcompref{descent}.
% Generated by scripts/build_blueprint_tex.py from data/model.json -- do not edit by hand.
\begin{bpleanquote}{lean/Atlas/Architecture.lean}{86}{118}
\begin{bplean}[firstnumber=86]
\LCom{/-- Classical facts used by the closing deductions. All of them are proved in the archive and}
\LCom{were kernel-checked at Lean 4.33.1 by the reading project's audit (`rh_of_right`, `Claim.strip`,}
\LCom{`Claim.zeta`). -/}
\LKw{structure} \LDecl{Reflection} \LKw{where}
  \LCom{/-- Every point is left of, on, or right of the critical line. -/}
  trichotomy : ∀ s, P.Left s ∨ P.Line s ∨ P.Right s
  \LCom{/-- `s ↦ 1 - s` maps the open strip to itself. -/}
  refl_strip : ∀ s, P.Strip s → P.Strip (P.refl s)
  \LCom{/-- `s ↦ 1 - s` sends the left half of the strip to the right half. -/}
  refl_right : ∀ s, P.Strip s → P.Left s → P.Right (P.refl s)
  \LCom{/-- The completed functional equation of a primitive character reflects a strip zero of `χ`}
\LCom{  to a zero of `χ⁻¹` at `1 - s`; the Gamma factors do not vanish in the strip. -/}
  functional_equation :
    ∀ ψ s, P.Primitive ψ → P.Strip s → P.Zero ψ s → P.Zero (P.inv ψ) (P.refl s)
  \LCom{/-- The open strip contains no pole. -/}
  strip_no_pole : ∀ χ s, P.Strip s → ¬ P.Pole χ s
  \LCom{/-- `ζ` is the `L`-function of the primitive character of modulus one, which is its own inverse. -/}
  zeta_primitive : P.Primitive P.zeta
  zeta_inv : P.inv P.zeta = P.zeta
  \LCom{/-- Outside the open strip `ζ` has only the trivial zeros (functional equation) and no zeros}
\LCom{  for `Re s ≥ 1` (in Mathlib); `s = 1` is the pole. Stated constructively: a zeta zero that is}
\LCom{  neither trivial nor at `1` lies in the strip. -/}
  zeta_strip : ∀ s, P.Zero P.zeta s → ¬ P.Trivial s → ¬ P.One s → P.Strip s

\LCom{/-- Conductor descent (archive: `Claim.prim_zero`, `UpperTheory.Top.imprimitive`):}
\LCom{`L(s,χ) = L(s,χ*) ∏_\{p ∣ N\} (1 - χ*(p) p^\{-s\})`, and no removed factor vanishes for `Re s > 0`}
\LCom{because `|χ*(p) p^\{-s\}| ≤ p^\{-Re s\} < 1`. -/}
\LKw{structure} \LDecl{Descent} \LKw{where}
  prim_primitive : ∀ χ, P.Primitive (P.prim χ)
  \LCom{/-- A zero of `L(s, χ)` with `Re s > 1/2`, away from the pole, is a zero of `L(s, χ*)`. -/}
  euler_right : ∀ χ s, P.Right s → ¬ P.Pole χ s → P.Zero χ s → P.Zero (P.prim χ) s
  \LCom{/-- The same in the open strip. -/}
  euler_strip : ∀ χ s, P.Strip s → P.Zero χ s → P.Zero (P.prim χ) s
\end{bplean}
\end{bpleanquote}

\subsubsection{The closing deductions} Primitive half-plane \(\Rightarrow\) half-plane \(\Rightarrow\)
strip and RH.
% Generated by scripts/build_blueprint_tex.py from data/model.json -- do not edit by hand.
\begin{bpleanquote}{lean/Atlas/Architecture.lean}{122}{143}
\begin{bplean}[firstnumber=122]
\LCom{/-- Primitive half-plane ⇒ half-plane for every character (mirrors `UpperTheory.Top.imprimitive`). -/}
\LKw{theorem} \LDecl{halfPlane_of_primitive} (D : Descent P) (h : PrimitiveHalfPlane P) : HalfPlane P :=
  \LKw{fun} χ s hs hpole hz => h (P.prim χ) s (D.prim_primitive χ) hs (D.euler_right χ s hs hpole hz)

\LCom{/-- Half-plane ⇒ Dirichlet GRH in the strip (mirrors `Claim.strip`). A zero left of the line}
\LCom{descends to the primitive character and reflects to a zero of its inverse on the right; a zero}
\LCom{on the right contradicts the hypothesis directly. -/}
\LKw{theorem} \LDecl{strip_of_halfPlane} (R : Reflection P) (D : Descent P) (h : HalfPlane P) : StripGRH P := \LKw{by}
  \LTac{intro} χ s hstrip hz
  \LTac{rcases} R.trichotomy s \LKw{with} hl | hline | hr
  · \LKw{have} hp : P.Zero (P.prim χ) s := D.euler_strip χ s hstrip hz
    \LKw{have} hzr : P.Zero (P.inv (P.prim χ)) (P.refl s) :=
      R.functional_equation _ s (D.prim_primitive χ) hstrip hp
    \LTac{exact} absurd hzr
      (h _ _ (R.refl_right s hstrip hl) (R.strip_no_pole _ _ (R.refl_strip s hstrip)))
  · \LTac{exact} hline
  · \LTac{exact} absurd hz (h χ s hr (R.strip_no_pole χ s hstrip))

\LCom{/-- Half-plane ⇒ RH (mirrors `Claim.zeta`): a nontrivial zeta zero lies in the strip, where the}
\LCom{previous theorem applies at modulus one. -/}
\LKw{theorem} \LDecl{rh_of_halfPlane} (R : Reflection P) (D : Descent P) (h : HalfPlane P) : RH P :=
  \LKw{fun} s hz htriv hone => strip_of_halfPlane R D h P.zeta s (R.zeta_strip s hz htriv hone) hz
\end{bplean}
\end{bpleanquote}

\subsubsection{The engine and the supply} The analytic criterion as a field, and the arithmetic supply
with the order of choices (\(\theta\) after the zero, \(\varepsilon\) after \(\theta\)).
% Generated by scripts/build_blueprint_tex.py from data/model.json -- do not edit by hand.
\begin{bpleanquote}{lean/Atlas/Architecture.lean}{147}{184}
\begin{bplean}[firstnumber=147]
\LCom{/-- The analytic engine over an abstract scale of exponents.}

\LCom{* `Exp` carries the exponents `θ`, with `Le θ θ'` for `θ ≤ θ'`;}
\LCom{* `RightOf s θ` is `Re s > θ`;}
\LCom{* `Family` carries the finite positive character families and `Mem χ F` is membership;}
\LCom{* `Windows F θ` is the one-sided estimate: one constant `C` and arbitrarily late windows}
\LCom{  `[Y, Y^(3/2)]` on which the family's centered prime sum is at most `C·y^θ` at every point. -/}
\LKw{structure} \LDecl{Engine} (P : Plane) \LKw{where}
  Exp : \LSort{Type}
  Le : Exp → Exp → \LSort{Prop}
  RightOf : P.Point → Exp → \LSort{Prop}
  Family : \LSort{Type}
  Mem : P.Char → Family → \LSort{Prop}
  Windows : Family → Exp → \LSort{Prop}
  \LCom{/-- A weaker exponent is easier to satisfy. -/}
  windows_mono : ∀ F θ θ', Windows F θ → Le θ θ' → Windows F θ'
  \LCom{/-- `Deduction.zeroFree_of_rpow_window`: a one-sided bound on arbitrarily late power windows}
\LCom{  forbids zeros of every member to the right of the exponent. Proof in the archive: the Mellin}
\LCom{  transform of the ramp, positive residues (`poleData`), and high derivatives aimed at the}
\LCom{  windows (`landau_window`). -/}
  criterion : ∀ F θ, Windows F θ → ∀ χ, Mem χ F → ∀ s, RightOf s θ → ¬ P.Zero χ s

\LCom{/-- The arithmetic supply, in the form consumed by `zeroFreeAll_of_topG`. For each primitive `ψ`}
\LCom{and each admissible small `ε` there is a family containing `ψ` whose prime sum satisfies the}
\LCom{window bound with exponent `loss ε = 1/2 + (e_D + e_L + 7)·ε`, where `e_D = 2`, `e_L = 0` were}
\LCom{fixed **before** `ε`. In the archive this is produced from `TopG` and `VWinAt ψ 401` (through}
\LCom{`supplyChi`, `core_gm` and `topDesWin`). It is the one non-classical input. -/}
\LKw{structure} \LDecl{Supply} (P : Plane) (E : Engine P) \LKw{where}
  Eps : \LSort{Type}
  Small : Eps → \LSort{Prop}
  loss : Eps → E.Exp
  family : P.Char → Eps → E.Family
  mem_family : ∀ ψ ε, E.Mem ψ (family ψ ε)
  windows : ∀ ψ ε, P.Primitive ψ → Small ε → E.Windows (family ψ ε) (loss ε)
  \LCom{/-- θ after the zero, ε after θ (`level_of_tiedWin`): for `Re s > 1/2` take}
\LCom{  `θ = 1/2 + min(Re s - 1/2, 1/2)/2`, so that `Re s > θ`, then a small `ε` with `loss ε ≤ θ`.}
\LCom{  This is possible because `loss ε → 1/2` as `ε → 0` with the losses fixed. -/}
  choose : ∀ s, P.Right s → ∃ θ, E.RightOf s θ ∧ ∃ ε, Small ε ∧ E.Le (loss ε) θ
\end{bplean}
\end{bpleanquote}

\subsubsection{The cuts in the skeleton}
% Generated by scripts/build_blueprint_tex.py from data/model.json -- do not edit by hand.
\begin{bpleanquote}{lean/Atlas/Architecture.lean}{186}{218}
\begin{bplean}[firstnumber=186]
\LCom{/-- The supply and the engine give the primitive half-plane (mirrors `zeroFreeAll_of_topG`). -/}
\LKw{theorem} \LDecl{primitiveHalfPlane_of_supply} (E : Engine P) (S : Supply P E) : PrimitiveHalfPlane P := \LKw{by}
  \LTac{intro} ψ s hψ hs hz
  \LTac{obtain} ⟨θ, hθ, ε, hε, hle⟩ := S.choose s hs
  \LTac{exact} E.criterion (S.family ψ ε) θ (E.windows_mono _ _ _ (S.windows ψ ε hψ hε) hle)
    ψ (S.mem_family ψ ε) s hθ hz

\LCom{/-- **The second verified cut** (mirrors `Reading.DirichletBridge.endpoints_of_topG_vWin401`,}
\LCom{kernel-checked against the archive in a 1,421-module closure): the arithmetic supply implies}
\LCom{RH, the half-plane theorem for every character and Dirichlet GRH in the strip. -/}
\LKw{theorem} \LDecl{endpoints_of_supply} (R : Reflection P) (D : Descent P) (E : Engine P) (S : Supply P E) :
    RH P ∧ HalfPlane P ∧ StripGRH P :=
  \LKw{have} hhalf := halfPlane_of_primitive D (primitiveHalfPlane_of_supply E S)
  ⟨rh_of_halfPlane R D hhalf, hhalf, strip_of_halfPlane R D hhalf⟩

\LCom{/-- **The first verified cut** (mirrors `Reading.WindowBridge.riemannHypothesis_of_prime_windows`,}
\LCom{checked by the Lean kernel and by the Rust kernel NanoDA): one family `\{ζ\}`, and the window bound}
\LCom{at every exponent `θ ∈ (1/2, 1)`, give RH. No other character and no descent is involved. -/}
\LKw{theorem} \LDecl{rh_of_prime_windows} (R : Reflection P) (E : Engine P) (Fz : E.Family)
    (hmem : E.Mem P.zeta Fz) (Above : E.Exp → \LSort{Prop})
    (choose : ∀ s, P.Right s → ∃ θ, Above θ ∧ E.RightOf s θ)
    (h : ∀ θ, Above θ → E.Windows Fz θ) : RH P := \LKw{by}
  \LTac{intro} s hz htriv hone
  \LKw{have} hstrip : P.Strip s := R.zeta_strip s hz htriv hone
  \LTac{rcases} R.trichotomy s \LKw{with} hl | hline | hr
  · \LKw{have} hz' : P.Zero (P.inv P.zeta) (P.refl s) :=
      R.functional_equation _ s R.zeta_primitive hstrip hz
    \LTac{rw} [R.zeta_inv] \LKw{at} hz'
    \LTac{obtain} ⟨θ, hθ, hright⟩ := choose (P.refl s) (R.refl_right s hstrip hl)
    \LTac{exact} absurd hz' (E.criterion Fz θ (h θ hθ) P.zeta hmem _ hright)
  · \LTac{exact} hline
  · \LTac{obtain} ⟨θ, hθ, hright⟩ := choose s hr
    \LTac{exact} absurd hz (E.criterion Fz θ (h θ hθ) P.zeta hmem s hright)
\end{bplean}
\end{bpleanquote}

\subsubsection{The whole architecture in one record}
% Generated by scripts/build_blueprint_tex.py from data/model.json -- do not edit by hand.
\begin{bpleanquote}{lean/Atlas/Architecture.lean}{222}{234}
\begin{bplean}[firstnumber=222]
\LCom{/-- The four layers of the argument bundled together. The first three are classical or}
\LCom{machine-checked; the fourth, `supply`, is where the arithmetic lives. -/}
\LKw{structure} \LDecl{Architecture} \LKw{where}
  plane : Plane
  reflection : Reflection plane
  descent : Descent plane
  engine : Engine plane
  supply : Supply plane engine

\LCom{/-- The five endpoints, as conjunctions on the abstract plane. -/}
\LKw{theorem} \LDecl{Architecture.endpoints} (A : Architecture) :
    RH A.plane ∧ HalfPlane A.plane ∧ StripGRH A.plane :=
  endpoints_of_supply A.reflection A.descent A.engine A.supply
\end{bplean}
\end{bpleanquote}

\subsubsection{A checked lemma from the cohomological window} The integer core of the ultrametric
separation used in component \bpcompref{cohomology}: each centre accounts for at most one exceptional
twist.
% Generated by scripts/build_blueprint_tex.py from data/model.json -- do not edit by hand.
\begin{bpleanquote}{lean/Atlas/Separation.lean}{31}{42}
\begin{bplean}[firstnumber=31]
\LCom{/-- Two twists in `[1, Z]` that are both `m`-adically close to the same centre `s`, with `m > Z`,}
\LCom{coincide: each centre accounts for at most one exceptional twist. -/}
\LKw{theorem} \LDecl{at_most_one_per_centre} \{a b s m Z : Int\} (hm : Z < m)
    (ha1 : 1 ≤ a) (haZ : a ≤ Z) (hb1 : 1 ≤ b) (hbZ : b ≤ Z)
    (hsa : m ∣ a - s) (hsb : m ∣ b - s) : a = b := \LKw{by}
  \LTac{obtain} ⟨c₁, h₁⟩ := hsa
  \LTac{obtain} ⟨c₂, h₂⟩ := hsb
  \LTac{refine} eq_of_dvd_sub hm ha1 haZ hb1 hbZ ⟨c₁ - c₂, ?_⟩
  \LTac{rw} [Int.mul_sub]
  \LTac{generalize} m * c₁ = t₁ \LKw{at} h₁ ⊢
  \LTac{generalize} m * c₂ = t₂ \LKw{at} h₂ ⊢
  \LTac{omega}
\end{bplean}
\end{bpleanquote}


\subsubsection{Axiom report} Generated \SkeletonGenerated{} from the build log of
\bppath{lean/Atlas/Audit.lean} (\bppath{verification/logs/axioms.json}); \SkeletonAxiomFree{} of
\SkeletonDecls{} declarations use no axiom at all, the others only Lean's standard axioms.
% Generated by scripts/build_blueprint_tex.py from data/model.json -- do not edit by hand.
\begin{bpaxiomtable}
\bpaxiomrow{\bpcode{Atlas.\allowbreak{}halfPlane\_\allowbreak{}of\_\allowbreak{}primitive}}{\bpnoaxioms}{1}
\bpaxiomrow{\bpcode{Atlas.\allowbreak{}strip\_\allowbreak{}of\_\allowbreak{}halfPlane}}{\bpnoaxioms}{1}
\bpaxiomrow{\bpcode{Atlas.\allowbreak{}rh\_\allowbreak{}of\_\allowbreak{}halfPlane}}{\bpnoaxioms}{1}
\bpaxiomrow{\bpcode{Atlas.\allowbreak{}primitiveHalfPlane\_\allowbreak{}of\_\allowbreak{}supply}}{\bpnoaxioms}{1}
\bpaxiomrow{\bpcode{Atlas.\allowbreak{}endpoints\_\allowbreak{}of\_\allowbreak{}supply}}{\bpnoaxioms}{1}
\bpaxiomrow{\bpcode{Atlas.\allowbreak{}rh\_\allowbreak{}of\_\allowbreak{}prime\_\allowbreak{}windows}}{\bpnoaxioms}{1}
\bpaxiomrow{\bpcode{Atlas.\allowbreak{}Architecture.\allowbreak{}endpoints}}{\bpnoaxioms}{1}
\bpaxiomrow{\bpcode{Atlas.\allowbreak{}toyArchitecture}}{\bpnoaxioms}{1}
\bpaxiomrow{\bpcode{Atlas.\allowbreak{}Separation.\allowbreak{}eq\_\allowbreak{}of\_\allowbreak{}dvd\_\allowbreak{}sub}}{\bpcode{propext}, \bpcode{Quot.\allowbreak{}sound}}{0}
\bpaxiomrow{\bpcode{Atlas.\allowbreak{}Separation.\allowbreak{}at\_\allowbreak{}most\_\allowbreak{}one\_\allowbreak{}per\_\allowbreak{}centre}}{\bpcode{propext}, \bpcode{Quot.\allowbreak{}sound}}{0}
\bpaxiomrow{\bpcode{Reading.\allowbreak{}excludeZeroByWindows}}{\bpnoaxioms}{1}
\bpaxiomrow{\bpcode{Reading.\allowbreak{}criticalLineByReflection}}{\bpnoaxioms}{1}
\bpaxiomrow{\bpcode{BehrendReading.\allowbreak{}normSq\_\allowbreak{}nonneg}}{\bpcode{propext}}{0}
\bpaxiomrow{\bpcode{BehrendReading.\allowbreak{}coordinates\_\allowbreak{}zero\_\allowbreak{}of\_\allowbreak{}normSq\_\allowbreak{}zero}}{\bpcode{propext}, \bpcode{Quot.\allowbreak{}sound}}{0}
\bpaxiomrow{\bpcode{BehrendReading.\allowbreak{}distance\_\allowbreak{}pos\_\allowbreak{}of\_\allowbreak{}ne}}{\bpcode{propext}, \bpcode{Quot.\allowbreak{}sound}}{0}
\bpaxiomrow{\bpcode{BehrendReading.\allowbreak{}chord\_\allowbreak{}identity}}{\bpcode{propext}, \bpcode{Classical.\allowbreak{}choice}, \bpcode{Quot.\allowbreak{}sound}}{0}
\bpaxiomrow{\bpcode{BehrendReading.\allowbreak{}sphere\_\allowbreak{}rel}}{\bpcode{propext}, \bpcode{Classical.\allowbreak{}choice}, \bpcode{Quot.\allowbreak{}sound}}{0}
\bpaxiomrow{\bpcode{BehrendReading.\allowbreak{}sphere\_\allowbreak{}midpoint}}{\bpcode{propext}, \bpcode{Classical.\allowbreak{}choice}, \bpcode{Quot.\allowbreak{}sound}}{0}
\end{bpaxiomtable}


%% ------------------------------------------------------------------ 8. blueprint table
\clearpage
\section{The blueprint: \NumDecls{} declarations}\label{sec:table}

The reading project's blueprint graph has \NumNodes{} nodes, \NumDecls{} Lean declarations and
\NumIdeas{} ideas from the announcement without a theorem, joined by \NumEdges{} edges; the model assigns
every node to exactly one component. The table lists them component by component, in the reading order, with
the declaration linked to its pinned source lines and its contract in one sentence. Node tiers are
computed from receipts: \bptiericonplain{checked} a named declaration with a kernel-check receipt,
\bptiericonplain{covered} called by one, \bptiericonplain{reported} the five exports,
\bptiericonplain{read} located and read, \bptiericonplain{motivation} motivation only. Figure~\ref{fig:graph}
draws the whole graph.
% Generated by scripts/build_blueprint_tex.py from data/model.json -- do not edit by hand.
\begin{bpnodetable}
\bpnodegroup{endpoints}{dest}{1}{The five endpoints}{13}
\bpnoderow{main-rh}{two-kernels}{Mathlib RH}{\href{https://github.com/Deicyde/riemann-hypothesis/blob/40e56bab026552e089d7afcfeff08d2defb9c005/src/Main.lean\#L4-L4}{\bpcode{riemannHypothesis}}}{Main.lean, L4}{Every nontrivial zeta zero distinct from the pole has real part one half.}
\bpnoderow{main-zeta}{two-kernels}{Zeta in expanded form}{\href{https://github.com/Deicyde/riemann-hypothesis/blob/40e56bab026552e089d7afcfeff08d2defb9c005/src/Main.lean\#L6-L6}{\bpcode{riemannZeta\_\allowbreak{}eq\_\allowbreak{}zero\_\allowbreak{}re}}}{Main.lean, L6}{Every zero other than the negative even integers and the point one has real part one half.}
\bpnoderow{main-right}{two-kernels}{Nonvanishing right of one half}{\href{https://github.com/Deicyde/riemann-hypothesis/blob/40e56bab026552e089d7afcfeff08d2defb9c005/src/Main.lean\#L8-L8}{\bpcode{LFunction\_\allowbreak{}ne\_\allowbreak{}zero\_\allowbreak{}of\_\allowbreak{}half\_\allowbreak{}lt\_\allowbreak{}re}}}{Main.lean, L8}{For all positive moduli and all characters, L(s,\(χ\)) is nonzero for Re(s)>1/2, excluding the principal pole.}
\bpnoderow{main-strip}{two-kernels}{Dirichlet GRH in the open strip}{\href{https://github.com/Deicyde/riemann-hypothesis/blob/40e56bab026552e089d7afcfeff08d2defb9c005/src/Main.lean\#L10-L10}{\bpcode{LFunction\_\allowbreak{}eq\_\allowbreak{}zero\_\allowbreak{}re}}}{Main.lean, L10}{For every character of positive modulus, any zero with 0<Re(s)<1 satisfies Re(s)=1/2.}
\bpnoderow{main-prim}{two-kernels}{Primitive critical-strip zeros}{\href{https://github.com/Deicyde/riemann-hypothesis/blob/40e56bab026552e089d7afcfeff08d2defb9c005/src/Main.lean\#L12-L12}{\bpcode{LFunction\_\allowbreak{}eq\_\allowbreak{}zero\_\allowbreak{}re\_\allowbreak{}of\_\allowbreak{}isPrimitive}}}{Main.lean, L12}{For every primitive character of positive modulus, all zeros in the open critical strip have real part one half.}
\bpnoderow{stmt-zeta}{read}{Zeta target type}{\href{https://github.com/Deicyde/riemann-hypothesis/blob/40e56bab026552e089d7afcfeff08d2defb9c005/src/Statements.lean\#L3-L4}{\bpcode{RiemannZetaZerosOnLine}}}{Statements.lean, L3–4}{∀ s∈ℂ, \(ζ\)(s)=0, s outside \{−2(n+1)\} for n∈ℕ and s≠1 imply Re(s)=1/2.}
\bpnoderow{stmt-right}{read}{Half-plane target type}{\href{https://github.com/Deicyde/riemann-hypothesis/blob/40e56bab026552e089d7afcfeff08d2defb9c005/src/Statements.lean\#L6-L8}{\bpcode{DirichletLNonvanishingRightOfHalf}}}{Statements.lean, L6–8}{∀ N≠0, \(χ\) mod N, s∈ℂ, 1/2<Re(s) and (\(χ\)≠1 or s≠1) imply L(s,\(χ\))≠0.}
\bpnoderow{stmt-strip}{read}{Dirichlet strip target type}{\href{https://github.com/Deicyde/riemann-hypothesis/blob/40e56bab026552e089d7afcfeff08d2defb9c005/src/Statements.lean\#L10-L12}{\bpcode{DirichletLZerosOnLine}}}{Statements.lean, L10–12}{∀ N≠0, \(χ\) mod N, s∈ℂ, 0<Re(s)<1 and L(s,\(χ\))=0 imply Re(s)=1/2.}
\bpnoderow{stmt-prim}{read}{Primitive strip target type}{\href{https://github.com/Deicyde/riemann-hypothesis/blob/40e56bab026552e089d7afcfeff08d2defb9c005/src/Statements.lean\#L14-L16}{\bpcode{PrimitiveDirichletLZerosOnLine}}}{Statements.lean, L14–16}{The preceding strip implication for each \(χ\) satisfying \bpcode{IsPrimitive}.}
\bpnoderow{rh-all}{read}{The two zeta exports share a proof}{\href{https://github.com/Deicyde/riemann-hypothesis/blob/40e56bab026552e089d7afcfeff08d2defb9c005/src/Supplements/Kernels/MainTheorems.lean\#L9-L9}{\bpcode{rhOfAll}}}{MainTheorems.lean, L9}{\bpcode{RiemannHypothesis}.}
\bpnoderow{all-chi}{read}{All-character half-plane theorem}{\href{https://github.com/Deicyde/riemann-hypothesis/blob/40e56bab026552e089d7afcfeff08d2defb9c005/src/Supplements/Kernels/MainTheorems.lean\#L7-L7}{\bpcode{allChi}}}{MainTheorems.lean, L7}{\bpcode{UpperTheory.\allowbreak{}EveryCharacter}.}
\bpnoderow{crit-all}{read}{Primitive critical-line wrapper}{\href{https://github.com/Deicyde/riemann-hypothesis/blob/40e56bab026552e089d7afcfeff08d2defb9c005/src/Supplements/Kernels/MainTheorems.lean\#L11-L13}{\bpcode{critAll}}}{MainTheorems.lean, L11–13}{For primitive \(χ\), every open-strip zero has real part one half.}
\bpnoderow{claim-all}{read}{Three endpoint propositions}{\href{https://github.com/Deicyde/riemann-hypothesis/blob/40e56bab026552e089d7afcfeff08d2defb9c005/src/Supplements/Kernels/MainTheorems.lean\#L15-L16}{\bpcode{claimAll}}}{MainTheorems.lean, L15–16}{\bpcode{HalfPlane} ∧ Strip ∧ Zeta.}
\bpnodegroup{reflection}{dest}{2}{Reflection}{8}
\bpnoderow{rh-every}{read}{Specialize at modulus one}{\href{https://github.com/Deicyde/riemann-hypothesis/blob/40e56bab026552e089d7afcfeff08d2defb9c005/src/Formalism/GammaFactorNonvanishing/CritLine.lean\#L29-L32}{\bpcode{Formalism.\allowbreak{}VW.\allowbreak{}rh\_\allowbreak{}of\_\allowbreak{}every}}}{CritLine.lean, L29–32}{\bpcode{EveryCharacter} implies \bpcode{RiemannHypothesis}.}
\bpnoderow{rh-right}{checked}{Reflect a zeta zero}{\href{https://github.com/Deicyde/riemann-hypothesis/blob/40e56bab026552e089d7afcfeff08d2defb9c005/src/Formalism/GammaFactorNonvanishing/RhOfRight.lean\#L5-L62}{\bpcode{Formalism.\allowbreak{}VW.\allowbreak{}rh\_\allowbreak{}of\_\allowbreak{}right}}}{RhOfRight.lean, L5–62}{Right-half-plane zeta nonvanishing away from one implies Mathlib RH.}
\bpnoderow{vw-crit}{read}{Reflect a primitive character}{\href{https://github.com/Deicyde/riemann-hypothesis/blob/40e56bab026552e089d7afcfeff08d2defb9c005/src/Formalism/GammaFactorNonvanishing/CritLine.lean\#L11-L27}{\bpcode{Formalism.\allowbreak{}VW.\allowbreak{}critLine}}}{CritLine.lean, L11–27}{\bpcode{EveryCharacter} implies the primitive open-strip conclusion.}
\bpnoderow{gamma}{read}{Gamma factors in the strip}{\href{https://github.com/Deicyde/riemann-hypothesis/blob/40e56bab026552e089d7afcfeff08d2defb9c005/src/Formalism/GammaFactorNonvanishing/CritLine.lean\#L6-L9}{\bpcode{Formalism.\allowbreak{}VW.\allowbreak{}gammaFactor\_\allowbreak{}ne\_\allowbreak{}zero}}}{CritLine.lean, L6–9}{For every character \(χ\) and Re(s)>0 its gamma factor at s is nonzero.}
\bpnoderow{strip}{checked}{Reflect an arbitrary character}{\href{https://github.com/Deicyde/riemann-hypothesis/blob/40e56bab026552e089d7afcfeff08d2defb9c005/src/Development/AutomorphicArithmetic/CharacterBoxes/ClaimStrip.lean\#L59-L73}{\bpcode{Claim.\allowbreak{}strip}}}{ClaimStrip.lean, L59–73}{\bpcode{Claim.\allowbreak{}HalfPlane} implies \bpcode{Claim.\allowbreak{}Strip}.}
\bpnoderow{prim-reflect}{covered}{Transfer a primitive zero to 1−s}{\href{https://github.com/Deicyde/riemann-hypothesis/blob/40e56bab026552e089d7afcfeff08d2defb9c005/src/Development/AutomorphicArithmetic/CharacterBoxes/ClaimStrip.lean\#L42-L57}{\bpcode{Claim.\allowbreak{}prim\_\allowbreak{}reflect}}}{ClaimStrip.lean, L42–57}{At an open-strip zero of a primitive \(ψ\), its inverse has a zero at 1−s.}
\bpnoderow{bundle-zeta}{checked}{Zeta component of the conjunction}{\href{https://github.com/Deicyde/riemann-hypothesis/blob/40e56bab026552e089d7afcfeff08d2defb9c005/src/Development/AutomorphicArithmetic/CharacterBoxes/ClaimStrip.lean\#L75-L100}{\bpcode{Claim.\allowbreak{}zeta}}}{ClaimStrip.lean, L75–100}{\bpcode{Claim.\allowbreak{}HalfPlane} implies \bpcode{Claim.\allowbreak{}Zeta}.}
\bpnoderow{alternate-rh}{read}{A second zeta deduction}{\href{https://github.com/Deicyde/riemann-hypothesis/blob/40e56bab026552e089d7afcfeff08d2defb9c005/src/Development/AutomorphicArithmetic/CharacterBoxes/RiemannHypothesisOfEveryCharacter.lean\#L215-L231}{\bpcode{UpperTheory.\allowbreak{}riemannHypothesis\_\allowbreak{}of\_\allowbreak{}everyCharacter}}}{RiemannHypothesisOfEveryCharacter.lean, L215–231}{\bpcode{EveryCharacter} implies Mathlib RH.}
\bpnodegroup{descent}{dest}{3}{Conductor descent}{4}
\bpnoderow{prim-zero}{covered}{Remove imprimitive Euler factors}{\href{https://github.com/Deicyde/riemann-hypothesis/blob/40e56bab026552e089d7afcfeff08d2defb9c005/src/Development/AutomorphicArithmetic/CharacterBoxes/ClaimStrip.lean\#L23-L33}{\bpcode{Claim.\allowbreak{}prim\_\allowbreak{}zero}}}{ClaimStrip.lean, L23–33}{An imprimitive character zero in 0<Re(s)<1 gives a zero of its primitive inducing character.}
\bpnoderow{factor}{covered}{Nonzero removed Euler factor}{\href{https://github.com/Deicyde/riemann-hypothesis/blob/40e56bab026552e089d7afcfeff08d2defb9c005/src/Development/AutomorphicArithmetic/CharacterBoxes/ClaimStrip.lean\#L7-L21}{\bpcode{Claim.\allowbreak{}factor\_\allowbreak{}ne\_\allowbreak{}zero}}}{ClaimStrip.lean, L7–21}{For prime p and Re(s)>0, 1−\(ψ\)(p)\(p^{−s}\)≠0.}
\bpnoderow{every-green}{read}{Pass from primitive to all characters}{\href{https://github.com/Deicyde/riemann-hypothesis/blob/40e56bab026552e089d7afcfeff08d2defb9c005/src/Development/AutomorphicArithmetic/ThetaCorrespondence/ThetaPairLifts/ThirdPlaceSeries/TopPlug.lean\#L16-L18}{\bpcode{UpperTheory.\allowbreak{}Top.\allowbreak{}every\_\allowbreak{}character\_\allowbreak{}of\_\allowbreak{}green\_\allowbreak{}v\_\allowbreak{}win}}}{TopPlug.lean, L16–18}{Green ∧ \bpcode{VWin} ∧ nonempty definite-side hypotheses imply \bpcode{EveryCharacter}, written as three binders.}
\bpnoderow{imprimitive}{checked}{Restore all characters}{\href{https://github.com/Deicyde/riemann-hypothesis/blob/40e56bab026552e089d7afcfeff08d2defb9c005/src/Development/AutomorphicArithmetic/CharacterBoxes/Imprimitive.lean\#L5-L35}{\bpcode{UpperTheory.\allowbreak{}Top.\allowbreak{}imprimitive}}}{Imprimitive.lean, L5–35}{\bpcode{ZeroFreeAll} for primitive characters implies the all-moduli right-half-plane assertion with the principal pole exception.}
\bpnodegroup{statistic}{ana}{4}{The centered prime sum}{8}
\bpnoderow{g-m}{covered}{The weighted prime sum}{\href{https://github.com/Deicyde/riemann-hypothesis/blob/40e56bab026552e089d7afcfeff08d2defb9c005/src/Formalism/SharpRestriction/RamificationFrobenius/TotientGaloisClasses/ExplicitFormulaErrorOfHostConditions.lean\#L38-L38}{\bpcode{Formalism.\allowbreak{}HC.\allowbreak{}Level.\allowbreak{}Gm}}}{ExplicitFormulaErrorOfHostConditions.lean, L38}{\(G_{P}\)(x)=G\(ζ\)(x) plus a finite weighted sum of real parts of twisted smoothed errors.}
\bpnoderow{window-def}{covered}{Uniform families on power windows}{\href{https://github.com/Deicyde/riemann-hypothesis/blob/40e56bab026552e089d7afcfeff08d2defb9c005/src/Formalism/SharpRestriction/RamificationFrobenius/TotientGaloisClasses/TopWin.lean\#L30-L32}{\bpcode{Formalism.\allowbreak{}HC.\allowbreak{}ChebyshevInequalityAtWindowScale}}}{TopWin.lean, L30–32}{There exist \bpcode{eD},\bpcode{eL} and \(ε\)₁>0 such that every 0<\(ε\)<\(ε\)₁ admits a tied family with whole windows of width \(λ\).}
\bpnoderow{moment-target}{covered}{The first moment of the Chebyshev error}{\href{https://github.com/Deicyde/riemann-hypothesis/blob/40e56bab026552e089d7afcfeff08d2defb9c005/src/Formalism/GL3MeanSquares/PrimeCountAsymptoticInProgressions.lean\#L18-L20}{\bpcode{Formalism.\allowbreak{}HermitianLattices.\allowbreak{}GL2a2.\allowbreak{}PsiZeroConst}}}{PrimeCountAsymptoticInProgressions.lean, L18–20}{The truncated integral of (\(ψ\)(t)−t)/t² from 1 to U tends to −1−\(γ\) as U tends to infinity.}
\bpnoderow{moment-theorem}{checked}{Evaluate the first moment}{\href{https://github.com/Deicyde/riemann-hypothesis/blob/40e56bab026552e089d7afcfeff08d2defb9c005/src/Formalism/GL3MeanSquares/HookFormulaTorsion/PsiZeroIntegral.lean\#L530-L531}{\bpcode{Formalism.\allowbreak{}HermitianLattices.\allowbreak{}GL2a2.\allowbreak{}psi\_\allowbreak{}zero\_\allowbreak{}integral\_\allowbreak{}constant}}}{PsiZeroIntegral.lean, L530–531}{The archive supplies its first-moment limit without an RH assumption.}
\bpnoderow{moment-calculation}{covered}{Match the two limits at one}{\href{https://github.com/Deicyde/riemann-hypothesis/blob/40e56bab026552e089d7afcfeff08d2defb9c005/src/Formalism/GL3MeanSquares/HookFormulaTorsion/PsiZeroIntegral.lean\#L423-L521}{\bpcode{psiZeroConst\_\allowbreak{}main}}}{PsiZeroIntegral.lean, L423–521}{For the trivial character modulo one the Chebyshev remainder has first moment −1−\(γ\).}
\bpnoderow{moment-laurent}{covered}{Read Euler’s constant from the logarithmic derivative}{\href{https://github.com/Deicyde/riemann-hypothesis/blob/40e56bab026552e089d7afcfeff08d2defb9c005/src/Formalism/GL3MeanSquares/HookFormulaTorsion/PsiZeroIntegral.lean\#L216-L277}{\bpcode{zeta\_\allowbreak{}logderiv\_\allowbreak{}limit}}}{PsiZeroIntegral.lean, L216–277}{As complex s tends puncturedly to 1, \(ζ\)′(s)/\(ζ\)(s)+1/(s−1) tends to \(γ\).}
\bpnoderow{moment-majorant}{covered}{Dominate the centered character sum}{\href{https://github.com/Deicyde/riemann-hypothesis/blob/40e56bab026552e089d7afcfeff08d2defb9c005/src/Formalism/GL3MeanSquares/HookFormulaTorsion/PsiZeroIntegral.lean\#L279-L301}{\bpcode{sw\_\allowbreak{}bound}}}{PsiZeroIntegral.lean, L279–301}{For each fixed character, the centered von Mangoldt sum is eventually bounded in norm by Kt/log²t.}
\bpnoderow{moment-dominated}{covered}{Pass the Mellin integral to \(σ\) = 1}{\href{https://github.com/Deicyde/riemann-hypothesis/blob/40e56bab026552e089d7afcfeff08d2defb9c005/src/Formalism/GL3MeanSquares/HookFormulaTorsion/PsiZeroIntegral.lean\#L157-L176}{\bpcode{core\_\allowbreak{}tendsto}}}{PsiZeroIntegral.lean, L157–176}{A logarithmic remainder bound lets \(σ\) times the weighted integral tend to its value at \(σ\)=1 from the right.}
\bpnodegroup{residues}{ana}{5}{Mellin transform and positive residues}{1}
\bpnoderow{pole}{checked}{Poles of the finite-family transform}{\href{https://github.com/Deicyde/riemann-hypothesis/blob/40e56bab026552e089d7afcfeff08d2defb9c005/src/Development/AutomorphicArithmetic/DedekindLFunctions/LandauTheorems/poleData.lean\#L248-L259}{\bpcode{Deduction.\allowbreak{}poleData}}}{poleData.lean, L248–259}{For each admissible finite family there exists a transform with \bpcode{PoleData} for its zero set.}
\bpnodegroup{criterion}{ana}{6}{The window criterion}{4}
\bpnoderow{form-b}{covered}{The analytic criterion for a level}{\href{https://github.com/Deicyde/riemann-hypothesis/blob/40e56bab026552e089d7afcfeff08d2defb9c005/src/Development/AutomorphicArithmetic/PrimitiveCosets/Imports.lean\#L8-L19}{\bpcode{UpperTheory.\allowbreak{}Comparison₂.\allowbreak{}formB}}}{Imports.lean, L8–19}{For every level, every \(λ\)>1 and 0<\(θ\)<1, arbitrarily late windows bounded by \(x^{θ}\) imply no zero right of \(θ\).}
\bpnoderow{landau}{checked}{The character-family power-window theorem}{\href{https://github.com/Deicyde/riemann-hypothesis/blob/40e56bab026552e089d7afcfeff08d2defb9c005/src/Development/AutomorphicArithmetic/DedekindLFunctions/LandauTheorems/ZeroFreeOfRpowWindow.lean\#L11-L14}{\bpcode{Deduction.\allowbreak{}zeroFree\_\allowbreak{}of\_\allowbreak{}rpow\_\allowbreak{}window}}}{ZeroFreeOfRpowWindow.lean, L11–14}{For a positive-weight family, a one-sided \(x^{θ}\) bound on arbitrarily late [Y,\(Y^{1+δ}\)] implies \bpcode{ZeroFree} F \(θ\).}
\bpnoderow{landau-lemma}{covered}{Apply the pole obstruction}{\href{https://github.com/Deicyde/riemann-hypothesis/blob/40e56bab026552e089d7afcfeff08d2defb9c005/src/Development/AutomorphicArithmetic/DedekindLFunctions/LandauTheorems/lemmaL.lean\#L13-L22}{\bpcode{Deduction.\allowbreak{}lemmaL}}}{lemmaL.lean, L13–22}{\bpcode{WindowHyp} on a finite character family implies \bpcode{ZeroFree} F \(θ\) in the given parameter range.}
\bpnoderow{oscillation}{covered}{Rule out a pole right of \(θ\)}{\href{https://github.com/Deicyde/riemann-hypothesis/blob/40e56bab026552e089d7afcfeff08d2defb9c005/src/Development/AutomorphicArithmetic/DedekindLFunctions/LandauTheorems/landauWindow.lean\#L83-L95}{\bpcode{Deduction.\allowbreak{}landau\_\allowbreak{}window}}}{landauWindow.lean, L83–95}{For the supplied \bpcode{PoleData} and \bpcode{WindowHypU}, any \(ρ\) in the specified pole set satisfies Re(\(ρ\))≤\(θ\).}
\bpnodegroup{choreography}{ana}{7}{The order of choices}{5}
\bpnoderow{zero-green}{read}{Select the window width 3/2}{\href{https://github.com/Deicyde/riemann-hypothesis/blob/40e56bab026552e089d7afcfeff08d2defb9c005/src/Development/AutomorphicArithmetic/ThetaCorrespondence/ThetaPairLifts/ThirdPlaceSeries/TopPlug.lean\#L12-L15}{\bpcode{UpperTheory.\allowbreak{}Top.\allowbreak{}zero\_\allowbreak{}free\_\allowbreak{}all\_\allowbreak{}of\_\allowbreak{}green\_\allowbreak{}v\_\allowbreak{}win}}}{TopPlug.lean, L12–15}{Green, \bpcode{VWin} and a nonempty definite-side package imply \bpcode{ZeroFreeAll}.}
\bpnoderow{zero-top}{checked}{Produce primitive zero-free levels}{\href{https://github.com/Deicyde/riemann-hypothesis/blob/40e56bab026552e089d7afcfeff08d2defb9c005/src/Development/AutomorphicArithmetic/CharacterBoxes/ZeroFreeAllOfTopG.lean\#L26-L34}{\bpcode{UpperTheory.\allowbreak{}Top.\allowbreak{}zeroFreeAll\_\allowbreak{}of\_\allowbreak{}topG}}}{ZeroFreeAllOfTopG.lean, L26–34}{For \(λ\)>1, \bpcode{TopG} and hosts for every primitive character imply \bpcode{ZeroFreeAll}.}
\bpnoderow{tied-cyc}{covered}{Assemble the cyclic levels}{\href{https://github.com/Deicyde/riemann-hypothesis/blob/40e56bab026552e089d7afcfeff08d2defb9c005/src/Formalism/SharpRestriction/RamificationFrobenius/TotientGaloisClasses/TopWin.lean\#L90-L93}{\bpcode{Formalism.\allowbreak{}HC.\allowbreak{}zeroFreeOn\_\allowbreak{}of\_\allowbreak{}tiedWin\_\allowbreak{}cyc}}}{TopWin.lean, L90–93}{Uniform window bounds, \bpcode{FormB}, and family existence for cyclic levels imply \bpcode{ZeroFreeOn} Q.}
\bpnoderow{tied}{covered}{Pass from levels to characters}{\href{https://github.com/Deicyde/riemann-hypothesis/blob/40e56bab026552e089d7afcfeff08d2defb9c005/src/Formalism/SharpRestriction/RamificationFrobenius/TotientGaloisClasses/TopWin.lean\#L83-L88}{\bpcode{Formalism.\allowbreak{}HC.\allowbreak{}zeroFreeOn\_\allowbreak{}of\_\allowbreak{}tiedWin}}}{TopWin.lean, L83–88}{Appropriate level data for each \(ψ\) satisfying Q imply \bpcode{ZeroFreeOn} Q.}
\bpnoderow{window-consumer}{covered}{Choose epsilon after a candidate zero}{\href{https://github.com/Deicyde/riemann-hypothesis/blob/40e56bab026552e089d7afcfeff08d2defb9c005/src/Formalism/SharpRestriction/RamificationFrobenius/TotientGaloisClasses/TopWin.lean\#L68-L81}{\bpcode{Formalism.\allowbreak{}HC.\allowbreak{}level\_\allowbreak{}of\_\allowbreak{}tiedWin}}}{TopWin.lean, L68–81}{At any level, the chosen uniform arithmetic window and the analytic criterion imply zero-freeness right of one half.}
\bpnodegroup{comparison}{est}{8}{The finite character comparison}{5}
\bpnoderow{prime-bound}{covered}{The prime-sum inequality on tied families}{\href{https://github.com/Deicyde/riemann-hypothesis/blob/40e56bab026552e089d7afcfeff08d2defb9c005/src/Development/AutomorphicArithmetic/PrimitiveCosets/LevelGmLeOnFamilyScales.lean\#L28-L30}{\bpcode{UpperTheory.\allowbreak{}Comparison₂.\allowbreak{}level\_\allowbreak{}gm\_\allowbreak{}le\_\allowbreak{}on\_\allowbreak{}family\_\allowbreak{}scales}}}{LevelGmLeOnFamilyScales.lean, L28–30}{For every level, tied compatible families yield the windowed upper bound on the weighted prime sum.}
\bpnoderow{artin-window}{checked}{Estimate the weighted level on a window}{\href{https://github.com/Deicyde/riemann-hypothesis/blob/40e56bab026552e089d7afcfeff08d2defb9c005/src/Development/AutomorphicArithmetic/PrimitiveCosets/WindowInequalityOfArtinInputs.lean\#L17-L25}{\bpcode{UpperTheory.\allowbreak{}Comparison₂.\allowbreak{}windowInequality\_\allowbreak{}of\_\allowbreak{}artinInputs}}}{WindowInequalityOfArtinInputs.lean, L17–25}{Identification of tied families with controlled host data gives \bpcode{ChebyshevInequalityOnWindows} at any level.}
\bpnoderow{core-gm}{checked}{From the height discrepancy to the prime sum}{\href{https://github.com/Deicyde/riemann-hypothesis/blob/40e56bab026552e089d7afcfeff08d2defb9c005/src/Development/AutomorphicArithmetic/PrimitiveCosets/CoreGm.lean\#L12-L74}{\bpcode{UpperTheory.\allowbreak{}Comparison₂.\allowbreak{}core\_\allowbreak{}gm}}}{CoreGm.lean, L12–74}{For large x, \bpcode{EAtH} and \bpcode{CoreAtH} imply \(G_{P}\)(x)≤\(x^{1/2+(}\)\bpcode{eD}+\bpcode{eL}+7)\(ε\)).}
\bpnoderow{ch-top}{checked}{Bound the discrepancy per index}{\href{https://github.com/Deicyde/riemann-hypothesis/blob/40e56bab026552e089d7afcfeff08d2defb9c005/src/Development/AutomorphicArithmetic/DedekindLFunctions/LandauTheorems/ChTop.lean\#L37-L72}{\bpcode{Deduction.\allowbreak{}ch\_\allowbreak{}top}}}{ChTop.lean, L37–72}{The named identifications and size, error, and norm windows give an upper bound for \bpcode{gapA}.}
\bpnoderow{ch-core}{checked}{Average over characters}{\href{https://github.com/Deicyde/riemann-hypothesis/blob/40e56bab026552e089d7afcfeff08d2defb9c005/src/Development/AutomorphicArithmetic/DedekindLFunctions/LandauTheorems/ChCore.lean\#L343-L351}{\bpcode{Deduction.\allowbreak{}ch\_\allowbreak{}core}}}{ChCore.lean, L343–351}{Under the size and character hypotheses, per-index discrepancy bounds imply a power bound for \bpcode{GmA} up to \bpcode{SigA}.}
\bpnodegroup{family}{est}{9}{The geometric family and the theta chain}{11}
\bpnoderow{goal-shape}{read}{Apply the shape to the same chosen data}{\href{https://github.com/Deicyde/riemann-hypothesis/blob/40e56bab026552e089d7afcfeff08d2defb9c005/src/Development/AutomorphicArithmetic/ThetaCorrespondence/ThetaPairLifts/ThirdPlaceSeries/ShapeOf.lean\#L31-L32}{\bpcode{GlobalTheory.\allowbreak{}Fc.\allowbreak{}forall\_\allowbreak{}goal\_\allowbreak{}of\_\allowbreak{}shape\_\allowbreak{}of\_\allowbreak{}hyps}}}{ShapeOf.lean, L31–32}{For propositions V, NR, G and a class predicate N, take \bpcode{TopShapeAlt} V NR N G, V, NR and \bpcode{HypsAlt} N \bpcode{DefiniteSideCondition} to G.}
\bpnoderow{shape-fields}{read}{Supply the actual fields of the shape}{\href{https://github.com/Deicyde/riemann-hypothesis/blob/40e56bab026552e089d7afcfeff08d2defb9c005/src/Development/AutomorphicArithmetic/ThetaCorrespondence/ThetaPairLifts/ThirdPlaceSeries/ShapeOf.lean\#L24-L30}{\bpcode{GlobalTheory.\allowbreak{}Fc.\allowbreak{}goal\_\allowbreak{}of\_\allowbreak{}shape\_\allowbreak{}of\_\allowbreak{}hyps}}}{ShapeOf.lean, L24–30}{\bpcode{TopShapeAlt} V NR N G with V, NR and \bpcode{HypsAlt} N \bpcode{DefiniteSideCondition} yields G.}
\bpnoderow{shape}{read}{Choose a common geometric family}{\href{https://github.com/Deicyde/riemann-hypothesis/blob/40e56bab026552e089d7afcfeff08d2defb9c005/src/Development/AutomorphicArithmetic/ThetaCorrespondence/ThetaPairLifts/ThirdPlaceSeries/TopShape.lean\#L33-L55}{\bpcode{UpperTheory.\allowbreak{}Top.\allowbreak{}shape\_\allowbreak{}closed\_\allowbreak{}link\_\allowbreak{}n\_\allowbreak{}of\_\allowbreak{}green}}}{TopShape.lean, L33–55}{Green implies the specified \bpcode{TopShapeAlt} ending in \bpcode{TopG}.}
\bpnoderow{family-design}{read}{Combine divisibility, slopes and windows}{\href{https://github.com/Deicyde/riemann-hypothesis/blob/40e56bab026552e089d7afcfeff08d2defb9c005/src/Development/AutomorphicArithmetic/CharacterBoxes/DesWin.lean\#L11-L31}{\bpcode{UpperTheory.\allowbreak{}Top.\allowbreak{}topDesWin}}}{DesWin.lean, L11–31}{The listed existence, local, rank, divisibility, global and slope controls produce \bpcode{FamilyClassConnected}.}
\bpnoderow{divin}{read}{Control the local divisibility loss}{\href{https://github.com/Deicyde/riemann-hypothesis/blob/40e56bab026552e089d7afcfeff08d2defb9c005/src/Development/AutomorphicArithmetic/CharacterBoxes/DivinNOfHypV.lean\#L117-L130}{\bpcode{UpperTheory.\allowbreak{}Top.\allowbreak{}divinN\_\allowbreak{}of\_\allowbreak{}hypV}}}{DivinNOfHypV.lean, L117–130}{For \(ε\)>0, \bpcode{HypV}, \bpcode{LinkN} and \bpcode{SmallLevelRankEq} imply \bpcode{DIVinN} on the family.}
\bpnoderow{vv-sum}{read}{Sum the local losses}{\href{https://github.com/Deicyde/riemann-hypothesis/blob/40e56bab026552e089d7afcfeff08d2defb9c005/src/Development/AutomorphicArithmetic/CharacterBoxes/DivinNOfHypV.lean\#L17-L57}{\bpcode{UpperTheory.\allowbreak{}Top.\allowbreak{}vvRecord}}}{DivinNOfHypV.lean, L17–57}{\bpcode{VVle} gives finite cohomology lengths, summability and the bound on the infinite sum of \bpcode{vvTerm}.}
\bpnoderow{rank}{read}{The geometric ranks}{\href{https://github.com/Deicyde/riemann-hypothesis/blob/40e56bab026552e089d7afcfeff08d2defb9c005/src/Development/AbelianVarietyTheory/Definitions/GlobalTwisting/PairRowIffQdRows.lean\#L28-L32}{\bpcode{GlobalTheory.\allowbreak{}EX.\allowbreak{}RankVal}}}{PairRowIffQdRows.lean, L28–32}{In particular, \(\sqrt{x δ₁}\)≤m≤2 \(\sqrt{x δ₁}\), with \(δ\)₁=\(x^{−ε/8}\).}
\bpnoderow{minor}{checked}{A local source of divisibility}{\href{https://github.com/Deicyde/riemann-hypothesis/blob/40e56bab026552e089d7afcfeff08d2defb9c005/src/Formalism/SharpRestriction/GrowthBounds/CoarseMultiplicity/KeyMinorDvd.lean\#L132-L155}{\bpcode{Formalism.\allowbreak{}MainNotions.\allowbreak{}ColumnSystems.\allowbreak{}Pf.\allowbreak{}key\_\allowbreak{}minor\_\allowbreak{}dvd}}}{KeyMinorDvd.lean, L132–155}{For w²=p and the displayed congruences, w divides a two-by-two minor.}
\bpnoderow{empty-data}{read}{An existence check on a period identity}{\href{https://github.com/Deicyde/riemann-hypothesis/blob/40e56bab026552e089d7afcfeff08d2defb9c005/src/Development/AutomorphicArithmetic/ThetaCorrespondence/ThetaPairLifts/ThirdPlaceSeries/VolLevMulWtMid.lean\#L24-L51}{\bpcode{GlobalTheory.\allowbreak{}Fam.\allowbreak{}volLev\_\allowbreak{}mul\_\allowbreak{}wt\_\allowbreak{}mid\_\allowbreak{}alt}}}{VolLevMulWtMid.lean, L24–51}{At a suitable intermediate stage there exists v≠0 such that level volume times character weight equals v at every admissible datum.}
\bpnoderow{slope}{read}{The lattice slope inequality}{\href{https://github.com/Deicyde/riemann-hypothesis/blob/40e56bab026552e089d7afcfeff08d2defb9c005/src/Formalism/LatticeSlopes/LatticeSlopesProofs/SlopeInequality.lean\#L111-L125}{\bpcode{Formalism.\allowbreak{}F.\allowbreak{}slopeInequality\_\allowbreak{}of}}}{SlopeInequality.lean, L111–125}{The sublattice bound \bpcode{SlopeLeMuMax} implies \bpcode{SlopeInequality} for injective norm-bounded lattice maps.}
\bpnoderow{paper-chain}{motivation}{The square-root multiplicity analogy}{\bpnodecl}{}{The author describes repeated theta lifts as a number-field analogue of multiplicity in the function-field argument.}
\bpnodegroup{hosts}{est}{10}{Hosts for every character}{1}
\bpnoderow{supply}{checked}{Produce cyclic character hosts}{\href{https://github.com/Deicyde/riemann-hypothesis/blob/40e56bab026552e089d7afcfeff08d2defb9c005/src/Development/AutomorphicArithmetic/CharacterBoxes/SupplyChi.lean\#L8-L16}{\bpcode{UpperTheory.\allowbreak{}Top.\allowbreak{}supplyChi}}}{SupplyChi.lean, L8–16}{For every primitive \(ψ\), \bpcode{VWinAt} \(ψ\) 401 implies host existence at width 3/2.}
\bpnodegroup{cohomology}{arith}{11}{The cohomological window}{25}
\bpnoderow{vv}{read}{Weighted lengths of local-global kernels}{\href{https://github.com/Deicyde/riemann-hypothesis/blob/40e56bab026552e089d7afcfeff08d2defb9c005/src/Development/AutomorphicArithmetic/CharacterBoxes/ExistsPageN.lean\#L12-L17}{\bpcode{UpperTheory.\allowbreak{}Top.\allowbreak{}VVle}}}{ExistsPageN.lean, L12–17}{For all admissible local data, lengths are finite for 1≤e<t and every finite prime sum of log(ℓ) times the sum of these lengths is ≤B.}
\bpnoderow{vwin-def}{covered}{The cohomological window}{\href{https://github.com/Deicyde/riemann-hypothesis/blob/40e56bab026552e089d7afcfeff08d2defb9c005/src/Development/AutomorphicArithmetic/CharacterBoxes/ExistsPageN.lean\#L44-L47}{\bpcode{UpperTheory.\allowbreak{}Top.\allowbreak{}VWinAt}}}{ExistsPageN.lean, L44–47}{For every \(ε\)∈(0,1/52) and admissible layer choice e(x), sufficiently large X and each prescribed pair q,d admit a window [Y,Y²] where \bpcode{VAt} holds.}
\bpnoderow{vwin-row}{read}{Cohomological window from the row}{\href{https://github.com/Deicyde/riemann-hypothesis/blob/40e56bab026552e089d7afcfeff08d2defb9c005/src/Development/AutomorphicArithmetic/AutomorphicVarieties/ModularityArgument/ThresholdArguments/DenseThresholdCases/VwinOfRow.lean\#L4-L5}{\bpcode{UpperTheory.\allowbreak{}One.\allowbreak{}vwin\_\allowbreak{}of\_\allowbreak{}row}}}{VwinOfRow.lean, L4–5}{\bpcode{Row0b} in the stated convention implies \bpcode{VWin}.}
\bpnoderow{window-inputs}{read}{Supply the automorphic window inputs}{\href{https://github.com/Deicyde/riemann-hypothesis/blob/40e56bab026552e089d7afcfeff08d2defb9c005/src/Development/AutomorphicArithmetic/AutomorphicVarieties/ModularityArgument/ThresholdArguments/DenseThresholdCases/TopMainZj.lean\#L71-L74}{\bpcode{UpperTheory.\allowbreak{}One.\allowbreak{}Join.\allowbreak{}automorphic\_\allowbreak{}window\_\allowbreak{}of\_\allowbreak{}inputs}}}{TopMainZj.lean, L71–74}{\bpcode{OpenImageCondition}, \bpcode{BaseChangeIsobaricN} and \bpcode{RowT} imply \bpcode{VWin}.}
\bpnoderow{depth-window}{read}{Build the rank-N package}{\href{https://github.com/Deicyde/riemann-hypothesis/blob/40e56bab026552e089d7afcfeff08d2defb9c005/src/Development/AutomorphicArithmetic/AutomorphicVarieties/ModularityArgument/ThresholdArguments/DenseThresholdCases/TopMainZj.lean\#L55-L58}{\bpcode{UpperTheory.\allowbreak{}One.\allowbreak{}Join.\allowbreak{}v\_\allowbreak{}window\_\allowbreak{}of\_\allowbreak{}depth\_\allowbreak{}zero\_\allowbreak{}residual\_\allowbreak{}bc\_\allowbreak{}cm}}}{TopMainZj.lean, L55–58}{The open-image, base-change, Galois and CM packages imply \bpcode{VWin}.}
\bpnoderow{lit-window}{read}{Window from rank-N automorphy}{\href{https://github.com/Deicyde/riemann-hypothesis/blob/40e56bab026552e089d7afcfeff08d2defb9c005/src/Development/AutomorphicArithmetic/AutomorphicVarieties/ModularityArgument/ThresholdArguments/DenseThresholdCases/WindowBoundOfLiteratureLevel.lean\#L33-L34}{\bpcode{UpperTheory.\allowbreak{}V.\allowbreak{}v\_\allowbreak{}window\_\allowbreak{}of\_\allowbreak{}literal\_\allowbreak{}automorphy\_\allowbreak{}n}}}{WindowBoundOfLiteratureLevel.lean, L33–34}{\bpcode{LiteratureRankNHypothesis} implies \bpcode{VWin}.}
\bpnoderow{patch-window}{read}{Pass through the descent hypotheses}{\href{https://github.com/Deicyde/riemann-hypothesis/blob/40e56bab026552e089d7afcfeff08d2defb9c005/src/Development/AutomorphicArithmetic/AutomorphicVarieties/ModularityArgument/ThresholdArguments/DenseThresholdCases/WindowBoundOfLiteratureLevel.lean\#L30-L31}{\bpcode{UpperTheory.\allowbreak{}V.\allowbreak{}HV.\allowbreak{}DescentVanishingReductions.\allowbreak{}v\_\allowbreak{}window\_\allowbreak{}of\_\allowbreak{}automorphy\_\allowbreak{}csd\_\allowbreak{}ss\_\allowbreak{}mono\_\allowbreak{}odd}}}{WindowBoundOfLiteratureLevel.lean, L30–31}{\bpcode{LiteratureRankNHypothesis} implies \bpcode{VWin}.}
\bpnoderow{four-produce}{read}{Produce the four counting bounds}{\href{https://github.com/Deicyde/riemann-hypothesis/blob/40e56bab026552e089d7afcfeff08d2defb9c005/src/Development/AutomorphicArithmetic/AutomorphicVarieties/ModularityArgument/ThresholdArguments/DenseThresholdCases/WindowBoundOfLiteratureLevel.lean\#L25-L29}{\bpcode{UpperTheory.\allowbreak{}V.\allowbreak{}HV.\allowbreak{}DescentVanishingReductions.\allowbreak{}vwin\_\allowbreak{}of\_\allowbreak{}hypotheses\_\allowbreak{}H\_\allowbreak{}E\_\allowbreak{}alt}}}{WindowBoundOfLiteratureLevel.lean, L25–29}{\bpcode{HypothesesH}′ imply \bpcode{VWin}.}
\bpnoderow{four-use}{read}{Combine the cohomology counts}{\href{https://github.com/Deicyde/riemann-hypothesis/blob/40e56bab026552e089d7afcfeff08d2defb9c005/src/Development/AutomorphicArithmetic/AutomorphicVarieties/ModularityArgument/ThresholdArguments/DenseThresholdCases/HVCore.lean\#L18-L22}{\bpcode{UpperTheory.\allowbreak{}V.\allowbreak{}HV.\allowbreak{}Core.\allowbreak{}of\_\allowbreak{}four\_\allowbreak{}alt}}}{HVCore.lean, L18–22}{Rank-N input, an existential K bound and \bpcode{CountI}, \bpcode{CountII}, \bpcode{CountIII}, \bpcode{CountV} imply \bpcode{VWin}.}
\bpnoderow{tower-window-from-lemma-k}{checked}{Assemble the arithmetic tower window}{\href{https://github.com/Deicyde/riemann-hypothesis/blob/40e56bab026552e089d7afcfeff08d2defb9c005/src/Development/AutomorphicArithmetic/AutomorphicVarieties/LayeredIsomorphisms/ThmBabOfLemmaK.lean\#L9-L12}{\bpcode{UpperTheory.\allowbreak{}V.\allowbreak{}Win.\allowbreak{}thmBab\_\allowbreak{}of\_\allowbreak{}lemmaK}}}{ThmBabOfLemmaK.lean, L9–12}{One package of constants works at every cutoff exponent Ct ≥ 1.}
\bpnoderow{tower-window-constants}{checked}{Choose one set of tower constants}{\href{https://github.com/Deicyde/riemann-hypothesis/blob/40e56bab026552e089d7afcfeff08d2defb9c005/src/Development/AutomorphicArithmetic/AutomorphicVarieties/LayeredIsomorphisms/ThmBabOf.lean\#L11-L103}{\bpcode{UpperTheory.\allowbreak{}V.\allowbreak{}Win.\allowbreak{}thmBab\_\allowbreak{}of}}}{ThmBabOf.lean, L11–103}{There exists a single \bpcode{BConst} K such that every Ct ≥ 1 satisfies \bpcode{ThmBab} K Ct.}
\bpnoderow{tower-bound-from-rank}{checked}{From a finitely generated tower to two length bounds}{\href{https://github.com/Deicyde/riemann-hypothesis/blob/40e56bab026552e089d7afcfeff08d2defb9c005/src/Development/AutomorphicArithmetic/AutomorphicVarieties/AdditiveValuationNorms/ABoundOfIsTower.lean\#L7-L11}{\bpcode{UpperTheory.\allowbreak{}V.\allowbreak{}PropA.\allowbreak{}aBound\_\allowbreak{}of\_\allowbreak{}isTower}}}{ABoundOfIsTower.lean, L7–11}{For each prime and eligible tower, \bpcode{ABound} consists of a bound away from finitely many centres and a bound in terms of a fixed layer.}
\bpnoderow{tower-operator-bound}{covered}{Control exceptional twists by the fixed layer}{\href{https://github.com/Deicyde/riemann-hypothesis/blob/40e56bab026552e089d7afcfeff08d2defb9c005/src/Development/AutomorphicArithmetic/AutomorphicVarieties/AdditiveValuationNorms/ABoundOfDescends.lean\#L36-L55}{\bpcode{UpperTheory.\allowbreak{}V.\allowbreak{}PropA.\allowbreak{}aBound\_\allowbreak{}of\_\allowbreak{}descends}}}{ABoundOfDescends.lean, L36–55}{The lengths are controlled by at most R ℓ-adic centres. For all twists they are also bounded by the fixed-layer H² length plus R times the excess layer index.}
\bpnoderow{tower-class-rank}{read}{Uniform ℓ-rank of class groups in the tower}{\href{https://github.com/Deicyde/riemann-hypothesis/blob/40e56bab026552e089d7afcfeff08d2defb9c005/src/Development/AutomorphicArithmetic/AutomorphicVarieties/GaloisLibrary/FerreroWashingtonIndex/FerreroWashington.lean\#L104-L119}{\bpcode{UpperTheory.\allowbreak{}V.\allowbreak{}FW.\allowbreak{}ferreroWashington}}}{FerreroWashington.lean, L104–119}{For each eligible abelian base K and prime ℓ there is an exponent r bounding the order of the quotient by ℓ-th powers by ℓʳ on every cyclotomic layer.}
\bpnoderow{tower-fixed-growth}{read}{Supply the absolute bound for the reference layer}{\href{https://github.com/Deicyde/riemann-hypothesis/blob/40e56bab026552e089d7afcfeff08d2defb9c005/src/Development/AutomorphicArithmetic/AutomorphicVarieties/ModularityArgument/ThresholdArguments/DenseThresholdCases/KFinal.lean\#L11-L11}{\bpcode{UpperTheory.\allowbreak{}V.\allowbreak{}ZK.\allowbreak{}lemmaK\_\allowbreak{}litN\_\allowbreak{}alt}}}{KFinal.lean, L11}{One real exponent bounds the logarithm of the order of $H^2(G_{L,S},\mathbb Z_\ell(e+1))$ over all eligible fields, coefficient primes and positive twists. Each such group is finite.}
\bpnoderow{counts-i}{read}{Case I gives one bound across coefficient primes}{\href{https://github.com/Deicyde/riemann-hypothesis/blob/40e56bab026552e089d7afcfeff08d2defb9c005/src/Development/AutomorphicArithmetic/AutomorphicVarieties/ModularityArgument/ThresholdArguments/DenseThresholdCases/LinkedUnlE.lean\#L49-L54}{\bpcode{UpperTheory.\allowbreak{}V.\allowbreak{}Z.\allowbreak{}J.\allowbreak{}count\_\allowbreak{}i\_\allowbreak{}of\_\allowbreak{}printed\_\allowbreak{}bound\_\allowbreak{}k\_\allowbreak{}kato\_\allowbreak{}automorphy}}}{LinkedUnlE.lean, L49–54}{Away from the exceptional set, when ℓ−1 does not divide twice the twist, the sum over every allowed finite set of coefficient primes obeys one polynomial bound.}
\bpnoderow{counts-ii-iii}{read}{Cases II and III from the rank-four symplectic argument}{\href{https://github.com/Deicyde/riemann-hypothesis/blob/40e56bab026552e089d7afcfeff08d2defb9c005/src/Development/AutomorphicArithmetic/AutomorphicVarieties/ModularityArgument/ThresholdArguments/DenseThresholdCases/AdjointLiftingOfCohomology.lean\#L33-L39}{\bpcode{UpperTheory.\allowbreak{}V.\allowbreak{}HV.\allowbreak{}DescentVanishingReductions.\allowbreak{}countII\_\allowbreak{}and\_\allowbreak{}countIII\_\allowbreak{}of\_\allowbreak{}printsH\_\allowbreak{}of\_\allowbreak{}sp4In}}}{AdjointLiftingOfCohomology.lean, L33–39}{Cases II and III have uniform polynomial bounds for each normalized character length, with finiteness. The hypotheses distinguish parity of character orders and divisibility of twice the twist.}
\bpnoderow{counts-v}{read}{Case V at twist one}{\href{https://github.com/Deicyde/riemann-hypothesis/blob/40e56bab026552e089d7afcfeff08d2defb9c005/src/Development/AutomorphicArithmetic/AutomorphicVarieties/ModularityArgument/ThresholdArguments/DenseThresholdCases/LinkedUnlE.lean\#L55-L61}{\bpcode{UpperTheory.\allowbreak{}V.\allowbreak{}Z.\allowbreak{}J.\allowbreak{}count\_\allowbreak{}v\_\allowbreak{}of\_\allowbreak{}printed\_\allowbreak{}bound\_\allowbreak{}k\_\allowbreak{}reduced\_\allowbreak{}unitary\_\allowbreak{}lift}}}{LinkedUnlE.lean, L55–61}{For exceptional odd primes at twist one, Case V bounds the normalized length when the two-adic conductor crosses the specified threshold, including equality.}
\bpnoderow{counts-common-integer}{read}{A nonzero integer controls one character orbit}{\href{https://github.com/Deicyde/riemann-hypothesis/blob/40e56bab026552e089d7afcfeff08d2defb9c005/src/Development/AutomorphicArithmetic/AutomorphicVarieties/ModularityArgument/ThresholdArguments/DenseThresholdCases/FPropGKU.lean\#L110-L172}{\bpcode{UpperTheory.\allowbreak{}V.\allowbreak{}Z.\allowbreak{}J.\allowbreak{}propGK\_\allowbreak{}of\_\allowbreak{}numberU3U\_\allowbreak{}of\_\allowbreak{}oneCharU}}}{FPropGKU.lean, L110–172}{For a fixed imaginary quadratic field, twist, character kernel and admissible auxiliary prime, one nonzero integer gives a logarithmic size bound and simultaneous valuation bounds at the other eligible primes.}
\bpnoderow{counts-window}{read}{Convert the joint window into \bpcode{VWin}}{\href{https://github.com/Deicyde/riemann-hypothesis/blob/40e56bab026552e089d7afcfeff08d2defb9c005/src/Development/AutomorphicArithmetic/AutomorphicVarieties/HeightAndVanishing/JointWinOfCounts.lean\#L30-L33}{\bpcode{UpperTheory.\allowbreak{}V.\allowbreak{}HV.\allowbreak{}vwin\_\allowbreak{}of\_\allowbreak{}counts}}}{JointWinOfCounts.lean, L30–33}{The tower window, sign bounds and five counting inputs yield \bpcode{VWin}, the cohomology estimate for each primitive character at cutoff exponent 401.}
\bpnoderow{counts-joint}{read}{Reuse the tower interval for the positive part}{\href{https://github.com/Deicyde/riemann-hypothesis/blob/40e56bab026552e089d7afcfeff08d2defb9c005/src/Development/AutomorphicArithmetic/AutomorphicVarieties/HeightAndVanishing/JointWinOfCounts.lean\#L25-L28}{\bpcode{UpperTheory.\allowbreak{}V.\allowbreak{}HV.\allowbreak{}jointWin\_\allowbreak{}of\_\allowbreak{}counts}}}{JointWinOfCounts.lean, L25–28}{For each eligible choice of character and parameters, one interval satisfies both the tower and positive-part bounds at every point.}
\bpnoderow{counts-height}{read}{Evaluate the pointwise estimate on admissible windows}{\href{https://github.com/Deicyde/riemann-hypothesis/blob/40e56bab026552e089d7afcfeff08d2defb9c005/src/Development/AutomorphicArithmetic/AutomorphicVarieties/HeightAndVanishing/JointWinOfCounts.lean\#L20-L23}{\bpcode{UpperTheory.\allowbreak{}V.\allowbreak{}HV.\allowbreak{}thmHV\_\allowbreak{}of\_\allowbreak{}counts}}}{JointWinOfCounts.lean, L20–23}{Every package K and positive cutoff exponent Ct satisfy the positive-part height theorem under the stated counting and cohomology inputs.}
\bpnoderow{counts-pointwise}{read}{Combine the two prime sums at a point}{\href{https://github.com/Deicyde/riemann-hypothesis/blob/40e56bab026552e089d7afcfeff08d2defb9c005/src/Development/AutomorphicArithmetic/AutomorphicVarieties/HeightAndVanishing/JointWinOfCounts.lean\#L16-L18}{\bpcode{UpperTheory.\allowbreak{}V.\allowbreak{}HV.\allowbreak{}pointHV\_\allowbreak{}of\_\allowbreak{}counts}}}{JointWinOfCounts.lean, L16–18}{After the character and layer indices are fixed, one polynomial bound controls the positive cohomological sum on eligible hosts, with an additive multiple of the layer budget.}
\bpnoderow{counts-inside}{read}{Add contributions inside the exceptional set}{\href{https://github.com/Deicyde/riemann-hypothesis/blob/40e56bab026552e089d7afcfeff08d2defb9c005/src/Development/AutomorphicArithmetic/AutomorphicVarieties/HeightAndVanishing/FinalLe.lean\#L86-L249}{\bpcode{UpperTheory.\allowbreak{}V.\allowbreak{}HV.\allowbreak{}sumIn}}}{FinalLe.lean, L86–249}{For fixed character, exceptional coefficient prime and layer index, the positive-part sum at that prime is at most (Nt)\(^{C}\) + C·B when the fixed layer satisfies its budget B.}
\bpnoderow{counts-outside}{read}{Bound the sum over arbitrary finite sets of other primes}{\href{https://github.com/Deicyde/riemann-hypothesis/blob/40e56bab026552e089d7afcfeff08d2defb9c005/src/Development/AutomorphicArithmetic/AutomorphicVarieties/HeightAndVanishing/SumOut.lean\#L50-L231}{\bpcode{UpperTheory.\allowbreak{}V.\allowbreak{}HV.\allowbreak{}sumOut}}}{SumOut.lean, L50–231}{Choose a common exponent and two auxiliary primes before the host. The positive-part sum over every allowed finite set of nonexceptional odd primes is at most (Nt)\(^{C}\).}
\bpnodegroup{primes}{arith}{12}{Auxiliary primes and Behrend towers}{2}
\bpnoderow{behrend}{checked}{Behrend’s collision-free indices}{\href{https://github.com/Deicyde/riemann-hypothesis/blob/40e56bab026552e089d7afcfeff08d2defb9c005/src/Formalism/FrobeniusOrbits/FrobeniusOrbitsProofs/ThickExists.lean\#L85-L95}{\bpcode{Formalism.\allowbreak{}FrobeniusOrbits.\allowbreak{}ThickExists.\allowbreak{}noRel\_\allowbreak{}image}}}{ThickExists.lean, L85–95}{For L≥1, encoding an integer sphere in base 4Ld gives a set with \bpcode{NoRel} for coefficients bounded by L.}
\bpnoderow{paper-behrend}{motivation}{The tower selection in the timeline}{\bpnodecl}{}{Entry 19 links the choice of a tower to Behrend-type sets.}
\bpnodegroup{galois}{arith}{13}{The Galois row}{6}
\bpnoderow{row-def}{read}{The local–global Galois row}{\href{https://github.com/Deicyde/riemann-hypothesis/blob/40e56bab026552e089d7afcfeff08d2defb9c005/src/Development/AutomorphicArithmetic/Foundations/GaloisRepresentations/TypeConditions/WeightsConjBis.lean\#L69-L73}{\bpcode{UpperTheory.\allowbreak{}GalStar.\allowbreak{}Row0b}}}{WeightsConjBis.lean, L69–73}{For a regular algebraic polarized cuspidal representation over a CM field and a prime ℓ, there exist a Galois representation, weight and Hodge–Tate multiset satisfying \bpcode{PropsB}.}
\bpnoderow{row-prod}{read}{Producer of the Galois row}{\href{https://github.com/Deicyde/riemann-hypothesis/blob/40e56bab026552e089d7afcfeff08d2defb9c005/src/Development/AutomorphicArithmetic/Foundations/GaloisRepresentations/IntegralStructures/RowSrc.lean\#L6-L6}{\bpcode{UpperTheory.\allowbreak{}Gal.\allowbreak{}Int.\allowbreak{}galoisLocalGlobalStrong\_\allowbreak{}recParameter}}}{RowSrc.lean, L6}{\bpcode{Row0b} \bpcode{recT} with no explicit global premise.}
\bpnoderow{row-geo}{read}{Galois representations from theta lifting}{\href{https://github.com/Deicyde/riemann-hypothesis/blob/40e56bab026552e089d7afcfeff08d2defb9c005/src/Development/AutomorphicArithmetic/Foundations/GaloisRepresentations/IntegralStructures/GaloisRepresentationsOfThetaLift.lean\#L6-L6}{\bpcode{UpperTheory.\allowbreak{}Gal.\allowbreak{}Int.\allowbreak{}galoisRepresentations\_\allowbreak{}of\_\allowbreak{}thetaLift}}}{GaloisRepresentationsOfThetaLift.lean, L6}{\bpcode{TLiftK} implies \bpcode{Row0b} \bpcode{recT}.}
\bpnoderow{tlift}{read}{Theta-lift existence at high rank}{\href{https://github.com/Deicyde/riemann-hypothesis/blob/40e56bab026552e089d7afcfeff08d2defb9c005/src/Development/AutomorphicArithmetic/Foundations/GaloisRepresentations/ThetaParseval/ThetaKernels/TLiftK.lean\#L4-L5}{\bpcode{UpperTheory.\allowbreak{}TP.\allowbreak{}K.\allowbreak{}tLiftK}}}{TLiftK.lean, L4–5}{For the admissible CM data of \bpcode{TLiftK}, for all sufficiently high dimensions and allowed weights a lift U′ exists.}
\bpnoderow{geo-one}{read}{Weight-one geometric realization}{\href{https://github.com/Deicyde/riemann-hypothesis/blob/40e56bab026552e089d7afcfeff08d2defb9c005/src/Development/AutomorphicArithmetic/Foundations/GaloisRepresentations/IntegralStructures/GeometricFlagBases/WeightOneGeometricRealization.lean\#L11-L11}{\bpcode{UpperTheory.\allowbreak{}Gal.\allowbreak{}Int.\allowbreak{}Geo.\allowbreak{}geoOne\_\allowbreak{}holds}}}{WeightOneGeometricRealization.lean, L11}{\bpcode{GeoOne} with no explicit premise.}
\bpnoderow{tlift-def}{read}{Quantifiers of the theta lift}{\href{https://github.com/Deicyde/riemann-hypothesis/blob/40e56bab026552e089d7afcfeff08d2defb9c005/src/Development/AutomorphicArithmetic/Foundations/GaloisRepresentations/ThetaParseval/HalfIndexLattices/TLiftK.lean\#L12-L23}{\bpcode{UpperTheory.\allowbreak{}TP.\allowbreak{}TLiftK}}}{TLiftK.lean, L12–23}{For each admissible CM setup, n≥1 and X, there exists d₀; for each d≥d₀, permissible W, Hecke characters, weights and U, an \bpcode{IsLift} witness U′ exists.}
\bpnodegroup{inputs}{arith}{14}{Seven inputs at the decisive call}{16}
\bpnoderow{every}{read}{Closed half-plane assembly}{\href{https://github.com/Deicyde/riemann-hypothesis/blob/40e56bab026552e089d7afcfeff08d2defb9c005/src/Development/AutomorphicArithmetic/CharacterBoxes/MainConclusion.lean\#L10-L11}{\bpcode{UpperTheory.\allowbreak{}FinalTheorems.\allowbreak{}everyCharacter}}}{MainConclusion.lean, L10–11}{\bpcode{EveryCharacter} with no explicit arithmetic premise.}
\bpnoderow{from-row}{read}{The half-plane theorem from one row}{\href{https://github.com/Deicyde/riemann-hypothesis/blob/40e56bab026552e089d7afcfeff08d2defb9c005/src/Development/AutomorphicArithmetic/CharacterBoxes/MainConclusion.lean\#L7-L8}{\bpcode{UpperTheory.\allowbreak{}FinalTheorems.\allowbreak{}everyCharacter\_\allowbreak{}of\_\allowbreak{}row}}}{MainConclusion.lean, L7–8}{\bpcode{Row0b} with the specified local parameter convention implies \bpcode{EveryCharacter}.}
\bpnoderow{rec}{read}{Existence of local Langlands parameters}{\href{https://github.com/Deicyde/riemann-hypothesis/blob/40e56bab026552e089d7afcfeff08d2defb9c005/src/Development/AutomorphicArithmetic/Foundations/LocalLanglands/LocalCentralUnits/RecTExists.lean\#L20-L31}{\bpcode{UpperTheory.\allowbreak{}Zero.\allowbreak{}LLC.\allowbreak{}recTExists}}}{RecTExists.lean, L20–31}{\bpcode{UpperTheory.\allowbreak{}Zero.\allowbreak{}LLC.\allowbreak{}RecTExists}.}
\bpnoderow{lit}{read}{Assemble the automorphic inputs}{\href{https://github.com/Deicyde/riemann-hypothesis/blob/40e56bab026552e089d7afcfeff08d2defb9c005/src/Development/AutomorphicArithmetic/ThetaCorrespondence/ThetaPairLifts/ThirdPlaceSeries/EveryCharacterOfLiteratureOfVWin.lean\#L810-L813}{\bpcode{UpperTheory.\allowbreak{}LiteratureTheorems.\allowbreak{}every\_\allowbreak{}character\_\allowbreak{}of\_\allowbreak{}literature\_\allowbreak{}of\_\allowbreak{}vWin}}}{EveryCharacterOfLiteratureOfVWin.lean, L810–813}{\bpcode{OpenLiteratureStatements} and \bpcode{VWin} imply \bpcode{EveryCharacter}.}
\bpnoderow{open-record}{read}{Three fields in the literature input}{\href{https://github.com/Deicyde/riemann-hypothesis/blob/40e56bab026552e089d7afcfeff08d2defb9c005/src/Development/AutomorphicArithmetic/ThetaCorrespondence/ThetaPairLifts/ThirdPlaceSeries/EveryCharacterOfLiteratureOfVWin.lean\#L215-L218}{\bpcode{UpperTheory.\allowbreak{}LiteratureTheorems.\allowbreak{}OpenLiteratureStatements}}}{EveryCharacterOfLiteratureOfVWin.lean, L215–218}{The record has fields Row, \bpcode{RecTExists} and Row, in that order.}
\bpnoderow{eigen-package}{read}{The unitary and eigenvalue package}{\href{https://github.com/Deicyde/riemann-hypothesis/blob/40e56bab026552e089d7afcfeff08d2defb9c005/src/Development/AutomorphicArithmetic/ThetaCorrespondence/ThetaPairLifts/ThirdPlaceSeries/EveryCharacterOfLiteratureOfVWin.lean\#L666-L672}{\bpcode{UpperTheory.\allowbreak{}LiteratureTheorems.\allowbreak{}published\_\allowbreak{}results\_\allowbreak{}eigen}}}{EveryCharacterOfLiteratureOfVWin.lean, L666–672}{\bpcode{OpenLiteratureStatements} imply \bpcode{UnitaryGroupThirdPlaceHypotheses}.}
\bpnoderow{theta-package}{read}{The local theta package}{\href{https://github.com/Deicyde/riemann-hypothesis/blob/40e56bab026552e089d7afcfeff08d2defb9c005/src/Development/AutomorphicArithmetic/ThetaCorrespondence/ThetaPairLifts/ThirdPlaceSeries/EveryCharacterOfLiteratureOfVWin.lean\#L727-L735}{\bpcode{UpperTheory.\allowbreak{}LiteratureTheorems.\allowbreak{}published\_\allowbreak{}results\_\allowbreak{}first}}}{EveryCharacterOfLiteratureOfVWin.lean, L727–735}{\bpcode{OpenLiteratureStatements} imply \bpcode{ThetaLocalInputs}.}
\bpnoderow{hodge-package}{read}{The Hodge-bundle package}{\href{https://github.com/Deicyde/riemann-hypothesis/blob/40e56bab026552e089d7afcfeff08d2defb9c005/src/Development/AutomorphicArithmetic/ThetaCorrespondence/ThetaPairLifts/ThirdPlaceSeries/EveryCharacterOfLiteratureOfVWin.lean\#L781-L786}{\bpcode{UpperTheory.\allowbreak{}LiteratureTheorems.\allowbreak{}published\_\allowbreak{}results\_\allowbreak{}arithmetic}}}{EveryCharacterOfLiteratureOfVWin.lean, L781–786}{\bpcode{OpenLiteratureStatements} imply \bpcode{PublishedHodgeBundleTheorems}.}
\bpnoderow{local-rec}{read}{The local parameter input at the call site}{\href{https://github.com/Deicyde/riemann-hypothesis/blob/40e56bab026552e089d7afcfeff08d2defb9c005/src/Development/AutomorphicArithmetic/ThetaCorrespondence/ThetaPairLifts/ThirdPlaceSeries/EveryCharacterOfLiteratureOfVWin.lean\#L313-L314}{\bpcode{UpperTheory.\allowbreak{}LiteratureTheorems.\allowbreak{}local\_\allowbreak{}langlands\_\allowbreak{}gl\_\allowbreak{}recT\_\allowbreak{}exists\_\allowbreak{}of\_\allowbreak{}literature}}}{EveryCharacterOfLiteratureOfVWin.lean, L313–314}{\bpcode{OpenLiteratureStatements} imply \bpcode{RecTExists} in the GL convention.}
\bpnoderow{local-rec-proj}{read}{Read the parameter field}{\href{https://github.com/Deicyde/riemann-hypothesis/blob/40e56bab026552e089d7afcfeff08d2defb9c005/src/Development/AutomorphicArithmetic/ThetaCorrespondence/ThetaPairLifts/ThirdPlaceSeries/EveryCharacterOfLiteratureOfVWin.lean\#L283-L284}{\bpcode{UpperTheory.\allowbreak{}LiteratureTheorems.\allowbreak{}local\_\allowbreak{}langlands\_\allowbreak{}recT\_\allowbreak{}exists\_\allowbreak{}of\_\allowbreak{}literature}}}{EveryCharacterOfLiteratureOfVWin.lean, L283–284}{\bpcode{OpenLiteratureStatements} imply \bpcode{RecTExists}.}
\bpnoderow{green}{read}{Regular characters yield the Green input}{\href{https://github.com/Deicyde/riemann-hypothesis/blob/40e56bab026552e089d7afcfeff08d2defb9c005/src/Development/AbelianVarietyTheory/Proofs/InertiaEigenvalues/SumRangeEqRoots.lean\#L225-L233}{\bpcode{GlobalTheory.\allowbreak{}GreenBridge.\allowbreak{}green\_\allowbreak{}of\_\allowbreak{}regular}}}{SumRangeEqRoots.lean, L225–233}{\bpcode{GreenRegularCharacters} imply \bpcode{GLq.\allowbreak{}Green}.}
\bpnoderow{linked-two-three}{read}{Bind the estimates at 2 and 3}{\href{https://github.com/Deicyde/riemann-hypothesis/blob/40e56bab026552e089d7afcfeff08d2defb9c005/src/Development/AutomorphicArithmetic/ThetaCorrespondence/ThetaPairLifts/ThirdPlaceSeries/At23OfAll.lean\#L24-L26}{\bpcode{UpperTheory.\allowbreak{}N.\allowbreak{}boundAtTwoThree\_\allowbreak{}linked\_\allowbreak{}of\_\allowbreak{}recT}}}{At23OfAll.lean, L24–26}{For local-global inputs, \bpcode{RecTExists} and a compatible dependence parameter, the bounds at two and three hold on linked data.}
\bpnoderow{theta-match}{read}{Match global integrals and local data}{\href{https://github.com/Deicyde/riemann-hypothesis/blob/40e56bab026552e089d7afcfeff08d2defb9c005/src/Development/AutomorphicArithmetic/ThetaCorrespondence/ThetaPairLifts/ThirdPlaceSeries/EveryCharacterOfLiteratureOfVWin.lean\#L565-L566}{\bpcode{UpperTheory.\allowbreak{}LiteratureTheorems.\allowbreak{}theta\_\allowbreak{}integrals\_\allowbreak{}agree\_\allowbreak{}and\_\allowbreak{}local\_\allowbreak{}matching\_\allowbreak{}of\_\allowbreak{}literature}}}{EveryCharacterOfLiteratureOfVWin.lean, L565–566}{\bpcode{OpenLiteratureStatements} imply \bpcode{ThetaIntegralsAgree} ∧ \bpcode{LocalMatching}.}
\bpnoderow{covol}{read}{Produce the covolume bound}{\href{https://github.com/Deicyde/riemann-hypothesis/blob/40e56bab026552e089d7afcfeff08d2defb9c005/src/Development/AutomorphicArithmetic/ThetaCorrespondence/ThetaPairLifts/ThirdPlaceSeries/TopCovolSlim2Use.lean\#L16-L17}{\bpcode{SelfIntersectionVolumeEstimates.\allowbreak{}topCovolSlim2\_\allowbreak{}use\_\allowbreak{}alt\_\allowbreak{}alt}}}{TopCovolSlim2Use.lean, L16–17}{The named geometric, arithmetic Hilbert–Samuel and Hodge-bundle hypotheses imply \bpcode{CovolumeBoundOfX}.}
\bpnoderow{seven}{read}{Seven inputs for every character}{\href{https://github.com/Deicyde/riemann-hypothesis/blob/40e56bab026552e089d7afcfeff08d2defb9c005/src/Development/AutomorphicArithmetic/ThetaCorrespondence/ThetaPairLifts/ThirdPlaceSeries/SiegelEisensteinSecondTermHypotheses.lean\#L69-L73}{\bpcode{UpperTheory.\allowbreak{}Lit.\allowbreak{}every\_\allowbreak{}character\_\allowbreak{}of\_\allowbreak{}theta\_\allowbreak{}integrals\_\allowbreak{}local\_\allowbreak{}matching}}}{SiegelEisensteinSecondTermHypotheses.lean, L69–73}{The seven groups of hypotheses imply \bpcode{EveryCharacter}.}
\bpnoderow{definite}{read}{Build a nonempty definite-side package}{\href{https://github.com/Deicyde/riemann-hypothesis/blob/40e56bab026552e089d7afcfeff08d2defb9c005/src/Development/AutomorphicArithmetic/ThetaCorrespondence/ThetaPairLifts/ThirdPlaceSeries/SiegelEisensteinSecondTermHypotheses.lean\#L44-L68}{\bpcode{UpperTheory.\allowbreak{}Lit.\allowbreak{}nonempty\_\allowbreak{}hyps\_\allowbreak{}definite\_\allowbreak{}side\_\allowbreak{}of\_\allowbreak{}closed\_\allowbreak{}bound}}}{SiegelEisensteinSecondTermHypotheses.lean, L44–68}{Unitary, theta, matching, covolume, local-parameter and linked-bound data yield Nonempty (\bpcode{HypsAlt} … \bpcode{DefiniteSideCondition}).}
\end{bpnodetable}

% Generated by scripts/build_blueprint_tex.py from data/model.json -- do not edit by hand.
\bpgraph{1772.25}{1772.25}{15.00}{25.50}{1.0000}{1.0000}{Blueprint graph: 109 nodes}{Layered dependency graph of the reading project's blueprint: suppliers on the left, the five endpoints on the right; bar colour = stratum, dot colour = verification tier.}


%% ------------------------------------------------------------------ 9. timeline
\clearpage
\section{The announcement, entry by entry}\label{sec:timeline}

The announcement lists \NumTimeline{} dated entries on its pages \TimelinePages. The dates are as written in
the announcement (a final ``+'' means \emph{from that date on}, drawn as a short arrow); the relation and
the assessment of each entry are the reading project's, and the component is the one of this map to
which the entry belongs. Dots are coloured by the stratum of that component.

\begin{center}
% Generated by scripts/build_blueprint_tex.py from data/model.json -- do not edit by hand.
\begin{tikzpicture}[x=1mm,y=1mm,bptimeline]
\draw[bptlaxis] (0.00,0) -- (94.80,0);
\node[bptlmonth] at (47.40,-6.2) {October 2026};
\draw[bptlaxis] (103.80,0) -- (142.80,0);
\node[bptlmonth] at (123.30,-6.2) {November 2026};
\node[bptlbreak] at (99.30,0) {};
\draw[bptltick] (4.00,0) -- ++(0,-1.2) node[below,bptlday] {6};
\draw[bptltick] (16.40,0) -- ++(0,-1.2) node[below,bptlday] {8};
\draw[bptltick] (22.60,0) -- ++(0,-1.2) node[below,bptlday] {9};
\draw[bptltick] (53.60,0) -- ++(0,-1.2) node[below,bptlday] {14};
\draw[bptltick] (84.60,0) -- ++(0,-1.2) node[below,bptlday] {19};
\draw[bptltick] (90.80,0) -- ++(0,-1.2) node[below,bptlday] {20};
\draw[bptltick] (107.80,0) -- ++(0,-1.2) node[below,bptlday] {15};
\draw[bptltick] (114.00,0) -- ++(0,-1.2) node[below,bptlday] {16};
\draw[bptlann] (138.80,-1.2) -- (138.80,44) node[above,bptlannlabel] {announcement};
\node[below,bptlday] at (138.80,-1.2) {20};
\node[bptldot=stratum-est] (tl1) at (4.00,3.60) {\hyperlink{tl:1}{1}};
\node[bptldot=stratum-est] (tl2) at (22.60,3.60) {\hyperlink{tl:2}{2}};
\node[bptldot=tier-motivation] (tl3) at (84.60,3.60) {\hyperlink{tl:3}{3}};
\draw[bptlfrom=tier-motivation] (tl3.east) -- ++(3.2,0);
\node[bptldot=tier-motivation] (tl4) at (4.00,8.20) {\hyperlink{tl:4}{4}};
\draw[bptlfrom=tier-motivation] (tl4.east) -- ++(3.2,0);
\node[bptldot=stratum-ana] (tl5) at (90.80,3.60) {\hyperlink{tl:5}{5}};
\node[bptldot=stratum-ana] (tl6) at (16.40,3.60) {\hyperlink{tl:6}{6}};
\node[bptldot=stratum-arith] (tl7) at (53.60,3.60) {\hyperlink{tl:7}{7}};
\node[bptldot=stratum-est] (tl8) at (107.80,3.60) {\hyperlink{tl:8}{8}};
\node[bptldot=stratum-est] (tl9) at (107.80,8.20) {\hyperlink{tl:9}{9}};
\node[bptldot=stratum-est] (tl10) at (107.80,12.80) {\hyperlink{tl:10}{10}};
\node[bptldot=stratum-arith] (tl11) at (4.00,12.80) {\hyperlink{tl:11}{11}};
\node[bptldot=stratum-arith] (tl12) at (107.80,17.40) {\hyperlink{tl:12}{12}};
\node[bptldot=stratum-est] (tl13) at (4.00,17.40) {\hyperlink{tl:13}{13}};
\node[bptldot=stratum-est] (tl14) at (107.80,22.00) {\hyperlink{tl:14}{14}};
\node[bptldot=stratum-est] (tl15) at (107.80,26.60) {\hyperlink{tl:15}{15}};
\node[bptldot=stratum-est] (tl16) at (53.60,8.20) {\hyperlink{tl:16}{16}};
\node[bptldot=stratum-arith] (tl17) at (107.80,31.20) {\hyperlink{tl:17}{17}};
\node[bptldot=stratum-arith] (tl18) at (107.80,35.80) {\hyperlink{tl:18}{18}};
\node[bptldot=stratum-arith] (tl19) at (84.60,8.20) {\hyperlink{tl:19}{19}};
\node[bptldot=stratum-dest] (tl20) at (114.00,3.60) {\hyperlink{tl:20}{20}};
\end{tikzpicture}

\end{center}
% Generated by scripts/build_blueprint_tex.py from data/model.json -- do not edit by hand.
\begin{bptimelinetable}
\bptlrow{1}{10/6/26}{6 October 2026}{Stepanov determinant}{3}{Historical motivation}{\bpcompref{family}}{The single-line entry records an origin of the argument. Its underlying function-field determinant is unspecified.}{stratum-est}
\bptlrow{2}{10/9/26}{9 October 2026}{Invariant differentials on a CM curve}{3}{Historical calculation}{\bpcompref{family}}{The entry describes inert-prime divisibility of iterated derivatives. The argument for this displayed exponent was not identified in the selected source path.}{stratum-est}
\bptlrow{3}{10/19/26+}{19 October 2026 onwards}{The finite-field sign argument}{3}{Historical analogy}{\bpnocomp}{The footnote compares a zero mode with an error bounded by iterated Parseval. Its objects and hypotheses differ from the Dirichlet L-function endpoint.}{tier-motivation}
\bptlrow{4}{10/6/26+}{6 October 2026 onwards}{Determinantal positivity with multiplicity one}{3}{Historical motivation}{\bpnocomp}{The short entry leaves W and S implicit. An exact theorem-to-type identification would require definitions of both.}{tier-motivation}
\bptlrow{5}{10/20/26}{20 October 2026}{The smoothed von Mangoldt sum}{3}{Formula matched}{\bpcompref{statistic}}{The displayed zeta error has the exact ramp and the correction (1 + Euler’s constant)x. Character twists enter the level average. The sentence continues onto page 4.}{stratum-ana}
\bptlrow{6}{10/8/26}{8 October 2026}{Two-support positivity test}{4}{Historical calculation}{\bpcompref{criterion}}{The two indicator functions are a proposed model for a zero-free exponent. The final power-window criterion has a different statement.}{stratum-ana}
\bptlrow{7}{10/14/26}{14 October 2026}{Interlacing weights and periods}{4}{Supporting comparison}{\bpcompref{inputs}}{A local Bessel model is constructed from two interlacing conditions. The announcement’s claim about determinant values 0 and 1 has further content.}{stratum-arith}
\bptlrow{8}{11/15/26}{15 November 2026}{A chain with square-root length}{4}{Quantitative comparison}{\bpcompref{family}}{The arithmetic rank specification puts m between one and two times x raised to 1/2 minus epsilon/16. The finite kernel recursion explains iteration.}{stratum-est}
\bptlrow{9}{11/15/26}{15 November 2026}{Norms, levels and arithmetic periods}{4}{Supporting comparison}{\bpcompref{family}}{One result relates a CM period to a logarithmic derivative. Another cancels a level volume against an index. The complete order-m cancellation requires more data.}{stratum-est}
\bptlrow{10}{11/15/26}{15 November 2026}{Divisibility of repeated lifts}{4}{Supporting comparison}{\bpcompref{family}}{A specific minor vanishes modulo a divisor and a lifted column has a higher-order remainder. Turning these local facts into the required global gain is an additional step.}{stratum-est}
\bptlrow{11}{10/6/26}{6 October 2026}{The tower of fields}{4}{Supporting comparison}{\bpcompref{cohomology}}{The formal window ranges over fields specified by a fixing subgroup in a cyclotomic extension and by layer parameter e(x).}{stratum-arith}
\bptlrow{12}{11/15/26}{15 November 2026}{Seesaw and Bessel periods}{4}{Supporting comparison}{\bpcompref{inputs}}{Conditional unfolding identifies a theta-lift period with an integral. A middle-step volume identity is independent of the admissible character. A full iterated comparison is a stronger assertion.}{stratum-arith}
\bptlrow{13}{10/6/26}{6 October 2026}{The slope inequality}{4}{Mathematical statement matched}{\bpcompref{family}}{The Euclidean lattice model bounds the slope of an injectively mapped source by the maximal target slope plus the logarithm of the map bound. Identification of the particular period lattice is separate.}{stratum-est}
\bptlrow{14}{11/15/26}{15 November 2026}{The provisional exponent two-thirds}{4}{Intermediate strategy}{\bpcompref{comparison}}{The author describes a barrier for an earlier estimate. The endpoint reads a one-half-plus-loss exponent. The earlier barrier is not derived in this comparison.}{stratum-est}
\bptlrow{15}{11/15/26}{15 November 2026}{A common integral structure across characters}{4}{Supporting comparison}{\bpcompref{family}}{The final deduction imposes local and global controls on the same family. A precise comparison with the gcd in the historical description remains open.}{stratum-est}
\bptlrow{16}{10/14/26}{14 October 2026}{Period pairings and approximate orthogonality}{4}{Supporting comparison}{\bpcompref{comparison}}{Unfolding relates a Poincare-series pairing to an integral of periods. The later host comparison turns suitable bounds into an estimate for the weighted prime sum.}{stratum-est}
\bptlrow{17}{11/15/26}{15 November 2026}{Cohomology and a weakening of Vandiver}{4}{Quantitative comparison}{\bpcompref{cohomology}}{The window requires finite lengths and an aggregate weighted bound. Further implications for individual K-groups are conjectural comments in the announcement.}{stratum-arith}
\bptlrow{18}{11/15/26}{15 November 2026}{Covolumes, lattice indices and congruences}{4}{Supporting comparison}{\bpcompref{inputs}}{A covolume bound is an explicit premise of the automorphic assembly. The local congruence lemma gives a representative higher-order equality.}{stratum-arith}
\bptlrow{19}{10/19/26}{19 October 2026}{A Behrend set in the tower construction}{4}{Supporting comparison}{\bpcompref{primes}}{Encoding a populated integer sphere in a large base excludes short affine relations and controls intersections with short orbits.}{stratum-arith}
\bptlrow{20}{11/16/26}{16 November 2026}{Dirichlet GRH}{4}{Endpoint statement matched}{\bpcompref{endpoints}}{The source declares that zeros of every Dirichlet L-function in the open critical strip are on the critical line. This entry records source-statement correspondence.}{stratum-dest}
\end{bptimelinetable}


%% ------------------------------------------------------------------ 10. sources, pins, reproduction
\clearpage
\section{Sources, pins and reproduction}\label{sec:pins}

\subsection{Pins}
% Generated by scripts/build_blueprint_tex.py from data/model.json -- do not edit by hand.
\begin{bppintable}
\bppinrow{Announcement PDF (SHA-256)}{\bphash{3fc97784b71ddd44ff8aa0a177e99622a1c61b13cfa843a0989ee4a7a0cfeb8d}}
\bppinrow{Lean archive (SHA-256 of the gzip stream, 14 parts)}{\bphash{9631b4f020199e580ba86178cb56c6e9b09f690c2369a80dc4334aa93f49c7fc}}
\bppinrow{Lean files in the archive}{175,187}
\bppinrow{Archive toolchain}{\bpcode{leanprover/\allowbreak{}lean4:v4.\allowbreak{}33.\allowbreak{}1}}
\bppinrow{Mathlib commit}{\href{https://github.com/leanprover-community/mathlib4/tree/0df444a360eaa60ab8c11dca51a86af692955474}{\bphash{0df444a360eaa60ab8c11dca51a86af692955474}}}
\bppinrow{Public mirror (all source links)}{\href{https://github.com/Deicyde/riemann-hypothesis/tree/40e56bab026552e089d7afcfeff08d2defb9c005}{\bpcode{Deicyde/\allowbreak{}riemann-hypothesis}} at \bphash{40e56bab026552e089d7afcfeff08d2defb9c005}}
\bppinrow{Reading project commit imported}{\bphash{348ec48748e0ea536aba7d64351b2ee2f7d116d7}}
\bppinrow{Import time}{2026-11-21 00:17:58 UTC}
\bppinrow{Status record as of}{2026-11-20 23:51:25 UTC, Helios job \bpcode{26647685}}
\bppinrow{Prime-window proof export}{61,697 declarations; 89,615,680 bytes compressed (SHA-256 \bphash{125bdd7041764c898f4d5b31653f3794abb61be09e484c1b791150f86cc084d0})}
\end{bppintable}


\subsection{The status record}\label{sec:status}
The reading project's status record as of \StatusAsOf, quoted verbatim (``we'' is the reading project).
It reports on the archive; this repository did not replay it.

\noindent\textbf{Headline.} ``\StatusHeadline''

\noindent\textbf{Checks made by the reading project.} ``\StatusNative''

\noindent\textbf{External report and authentication.} ``\StatusExternal''

\noindent\textbf{The implication at the centre.} ``\StatusDetail''

\noindent\textbf{Build job} (Helios \StatusJob). ``\StatusFull.

\noindent External sources, as recorded by the reading project:
% Generated by scripts/build_blueprint_tex.py from data/model.json -- do not edit by hand.
\begin{bpsourcelist}
\item \href{https://math.mit.edu/~drew/grh/con-leche-compare.log}{Sutherland’s five named con-leche results}
\item \href{https://math.mit.edu/~drew/grh/lean4lean-raised-fuel.log}{Sutherland’s five raised-fuel lean4lean reruns}
\item \href{https://math.mit.edu/~drew/grh/}{Full solution and dated operator report}
\item \href{https://photos.proofsandreasons.io/wright/alp-replay.json}{Second complete original-toolchain replay}
\item \href{https://math.mit.edu/~drew/grh/audit.log}{Five individual axiom lists}
\item \href{https://math.mit.edu/~drew/grh/nanoda.log}{Rust process log}
\item \href{https://math.mit.edu/~drew/grh/SHA256SUMS}{Download checksums}
\item \href{https://leanprover.zulipchat.com/\#narrow/channel/113488-general/topic/Discussion.3A.20Alp.C3.B6ge.20Proves.20the.20Riemann.20Hypothesis/near/642313194}{Earlier challenge comparison}
\item \href{https://leanprover.zulipchat.com/\#narrow/channel/113488-general/topic/Discussion.3A.20Alp.C3.B6ge.20Proves.20the.20Riemann.20Hypothesis/near/642338051}{Independent five-target comparison}
\item \href{https://github.com/Infatoshi/rh-verification-receipts/blob/ac99f18cef21952de094234fab95db342ed16638/VERIFICATION_RESULT.txt}{Independent rebuild with a separately generated proof export}
\end{bpsourcelist}


\subsection{Provenance of the imported data}
The data were imported from the reading project at commit \bphash{\SourceCommit} on \ImportedAt;
the files and their SHA-256 prefixes are:
% Generated by scripts/build_blueprint_tex.py from data/model.json -- do not edit by hand.
\begin{bpprovenance}
\bpprovrow{\bpcode{data/\allowbreak{}source/\allowbreak{}map.\allowbreak{}json}}{\bpcode{blueprint/\allowbreak{}map.\allowbreak{}json}}{26ad5dd085d155ba}
\bpprovrow{\bpcode{data/\allowbreak{}source/\allowbreak{}status.\allowbreak{}json}}{\bpcode{evidence/\allowbreak{}status.\allowbreak{}json}}{d4afbf5694ca44ca}
\bpprovrow{\bpcode{data/\allowbreak{}source/\allowbreak{}sources.\allowbreak{}json}}{\bpcode{evidence/\allowbreak{}sources.\allowbreak{}json}}{19f64f8f967813cd}
\bpprovrow{\bpcode{data/\allowbreak{}source/\allowbreak{}timeline.\allowbreak{}json}}{\bpcode{research/\allowbreak{}agents/\allowbreak{}arithmetic/\allowbreak{}timeline.\allowbreak{}json}}{8cc2d4c4dedc8f7c}
\bpprovrow{\bpcode{data/\allowbreak{}source/\allowbreak{}native-scopes.\allowbreak{}json}}{\bpcode{dashboard/\allowbreak{}data/\allowbreak{}native-scopes.\allowbreak{}json}}{13b55c110571a997}
\bpprovrow{\bpcode{data/\allowbreak{}source/\allowbreak{}helios-progress.\allowbreak{}json}}{\bpcode{evidence/\allowbreak{}helios/\allowbreak{}progress.\allowbreak{}json}}{44d9fd51367c85aa}
\bpprovrow{\bpcode{lean/\allowbreak{}Reading/\allowbreak{}WindowLogic.\allowbreak{}lean}}{\bpcode{formal/\allowbreak{}standalone/\allowbreak{}WindowLogic.\allowbreak{}lean}}{cd648265b6aa58d8}
\bpprovrow{\bpcode{lean/\allowbreak{}Reading/\allowbreak{}BehrendSphere.\allowbreak{}lean}}{\bpcode{formal/\allowbreak{}behrend/\allowbreak{}BehrendSphere.\allowbreak{}lean}}{a026ef4bc8c5eee7}
\bpprovrow{\bpcode{references/\allowbreak{}archive-bridges/\allowbreak{}WindowBridge.\allowbreak{}lean}}{\bpcode{formal/\allowbreak{}window-bridge/\allowbreak{}WindowBridge.\allowbreak{}lean}}{e77eb05549e7622b}
\bpprovrow{\bpcode{references/\allowbreak{}archive-bridges/\allowbreak{}GeometricTransport.\allowbreak{}lean}}{\bpcode{formal/\allowbreak{}window-bridge/\allowbreak{}GeometricTransport.\allowbreak{}lean}}{f6213d4c45e70a1a}
\bpprovrow{\bpcode{references/\allowbreak{}archive-bridges/\allowbreak{}DirichletBridge.\allowbreak{}lean}}{\bpcode{formal/\allowbreak{}dirichlet-bridge/\allowbreak{}DirichletBridge.\allowbreak{}lean}}{59d543a5550f953b}
\bpprovrow{\bpcode{references/\allowbreak{}archive-bridges/\allowbreak{}EndpointAudit.\allowbreak{}lean}}{\bpcode{formal/\allowbreak{}upstream/\allowbreak{}EndpointAudit.\allowbreak{}lean}}{7946558ed7e5f2da}
\end{bpprovenance}


\subsection{Reproduction}
Everything in this PDF is regenerated from the repository by the commands below (run from its root).
The generator is deterministic; the PDF date is fixed to the status time through
\bpcode{SOURCE\_DATE\_EPOCH}.
\begin{bpshell}
python3 -B scripts/model.py                 # data/model.json from the model and data/source
python3 -B scripts/build_figures.py         # assets/figures/*.svg and facts.json
python3 -B scripts/build_blueprint_tex.py   # blueprint/generated/*.tex
make -C blueprint                           # SVG -> PDF (rsvg-convert), LuaLaTeX via latexmk
(cd lean && lake build)                     # the skeleton and its axiom audit (lean/lean-toolchain)
\end{bpshell}
The LaTeX build needs TeX~Live (LuaLaTeX with the fonts Libertinus, Inter and JetBrains Mono, all part of
TeX~Live) and \bpcode{rsvg-convert}; intermediates go to \bpcode{blueprint/build/}.

\end{document}
